\documentclass[11pt,a4paper]{article}
\usepackage[a4paper,margin=20mm]{geometry}
\usepackage{fontspec}
\setmainfont{Noto Serif}
\setsansfont{Noto Sans}
\usepackage{amsmath,amssymb}
\usepackage{booktabs,tabularx,array}
\usepackage[table]{xcolor}
\usepackage{hyperref}
\definecolor{primary}{RGB}{27,79,114}
\definecolor{accent}{RGB}{40,116,166}
\definecolor{rowtint}{RGB}{235,245,251}
\hypersetup{colorlinks=true,linkcolor=primary,urlcolor=primary}
\setlength{\parskip}{5pt plus 2pt minus 1pt}
\setlength{\parindent}{0pt}
\begin{document}
\begin{center}
  {\LARGE\bfseries\color{primary} PHIẾU HỌC TẬP TRẮC NGHIỆM TOÁN 8}\\[4pt]
  {\normalsize\itshape\color{accent} Bộ sách: Kết nối tri thức với cuộc sống (Tập 1 \& Tập 2 – 10 Chương)}\\[4pt]
  {\small Cấu trúc mỗi chương: 20 câu Dễ (Câu 1--20) $\bullet$ 20 câu Trung bình (Câu 21--40) $\bullet$ 10 câu Khó (Câu 41--50) -- Tổng 50 câu/chương}
\end{center}
\vspace{6pt}\hrule\vspace{10pt}
\section*{\color{primary} CHƯƠNG I: ĐA THỨC}
\textit{Trọng tâm kiến thức SGK: Đơn thức và đa thức nhiều biến, đơn thức đồng dạng, các phép tính cộng, trừ, nhân đa thức và chia đa thức cho đơn thức.}\par\vspace{4pt}
\subsection*{\color{accent} Phần I. Mức độ Dễ (Nhận biết -- Câu 1 đến Câu 20)}
\textbf{\color{primary} Câu 1.} Bậc của đơn thức thu gọn A = -2x$^{2y}$$^3$z là:\par
\hspace*{1.5em}\textbf{A.} Bậc 2\quad\quad \textbf{B.} Bậc 6\quad\quad \textbf{C.} Bậc 6 (1)\quad\quad \textbf{D.} Bậc 5\par\vspace{2pt}
\textbf{\color{primary} Câu 2.} Thực hiện yêu cầu nhận biết sau: Thu gọn tổng hai đơn thức đồng dạng 5x$^2$y$^3$ + 3x$^2$y$^3$ ta được.\par
\hspace*{1.5em}\textbf{A.} 15x$^2$y$^3$\quad\quad \textbf{B.} 8x$^2$y$^3$\quad\quad \textbf{C.} 8x$^4$y$^6$\quad\quad \textbf{D.} -2x$^2$y$^3$\par\vspace{2pt}
\textbf{\color{primary} Câu 3.} Trong bài học thuộc Chương I, xét nội dung: Thương của phép chia đơn thức 18x$^4$y$^3$ cho đơn thức 2x$^2$y là. Đáp án đúng là:\par
\hspace*{1.5em}\textbf{A.} 9x$^2$y$^2$\quad\quad \textbf{B.} 12x$^2$y$^2$\quad\quad \textbf{C.} 9x$^6$y$^4$\quad\quad \textbf{D.} 9xy$^2$\par\vspace{2pt}
\textbf{\color{primary} Câu 4.} Dựa vào kiến thức chuẩn SGK của Chương I (Đa thức), bậc của đơn thức thu gọn A = -5x$^{5y}$$^3$z là:\par
\hspace*{1.5em}\textbf{A.} Bậc 15\quad\quad \textbf{B.} Bậc 5\quad\quad \textbf{C.} Bậc 8\quad\quad \textbf{D.} Bậc 9\par\vspace{2pt}
\textbf{\color{primary} Câu 5.} Chọn kết quả đúng nhất khi xét: Thu gọn tổng hai đơn thức đồng dạng 8x$^2$y$^3$ + 9x$^2$y$^3$ ta được:\par
\hspace*{1.5em}\textbf{A.} 72x$^2$y$^3$\quad\quad \textbf{B.} 17x$^2$y$^3$\quad\quad \textbf{C.} 1x$^2$y$^3$\quad\quad \textbf{D.} 17x$^4$y$^6$\par\vspace{2pt}
\textbf{\color{primary} Câu 6.} Khi làm bài tập cơ bản về đa thức, một học sinh gặp câu hỏi: Thương của phép chia đơn thức 36x$^4$y$^3$ cho đơn thức 2x$^2$y là. Kết quả nào dưới đây đúng?\par
\hspace*{1.5em}\textbf{A.} 18xy$^2$\quad\quad \textbf{B.} 18x$^2$y$^2$\quad\quad \textbf{C.} 24x$^2$y$^2$\quad\quad \textbf{D.} 18x$^6$y$^4$\par\vspace{2pt}
\textbf{\color{primary} Câu 7.} Xác định giá trị hoặc khẳng định đúng cho bài toán sau: Bậc của đơn thức thu gọn A = -8x$^{8y}$$^3$z là:\par
\hspace*{1.5em}\textbf{A.} Bậc 12\quad\quad \textbf{B.} Bậc 11\quad\quad \textbf{C.} Bậc 8\quad\quad \textbf{D.} Bậc 24\par\vspace{2pt}
\textbf{\color{primary} Câu 8.} Áp dụng trực tiếp định nghĩa và quy tắc trong Chương I: Thu gọn tổng hai đơn thức đồng dạng 11x$^2$y$^3$ + 15x$^2$y$^3$ ta được:\par
\hspace*{1.5em}\textbf{A.} 4x$^2$y$^3$\quad\quad \textbf{B.} 26x$^2$y$^3$\quad\quad \textbf{C.} 26x$^4$y$^6$\quad\quad \textbf{D.} 165x$^2$y$^3$\par\vspace{2pt}
\textbf{\color{primary} Câu 9.} Kiểm tra nhanh kiến thức nền tảng: Thương của phép chia đơn thức 54x$^4$y$^3$ cho đơn thức 2x$^2$y là:\par
\hspace*{1.5em}\textbf{A.} 27xy$^2$\quad\quad \textbf{B.} 27x$^2$y$^2$\quad\quad \textbf{C.} 36x$^2$y$^2$\quad\quad \textbf{D.} 27x$^6$y$^4$\par\vspace{2pt}
\textbf{\color{primary} Câu 10.} Cho dữ kiện nhận biết sau đây: Bậc của đơn thức thu gọn A = -11x$^{11y}$$^3$z là. Phương án chính xác là:\par
\hspace*{1.5em}\textbf{A.} Bậc 11\quad\quad \textbf{B.} Bậc 15\quad\quad \textbf{C.} Bậc 14\quad\quad \textbf{D.} Bậc 33\par\vspace{2pt}
\textbf{\color{primary} Câu 11.} Hãy chỉ ra kết quả đúng trong bốn lựa chọn A, B, C, D đối với yêu cầu: Thu gọn tổng hai đơn thức đồng dạng 14x$^2$y$^3$ + 21x$^2$y$^3$ ta được:\par
\hspace*{1.5em}\textbf{A.} 294x$^2$y$^3$\quad\quad \textbf{B.} 35x$^4$y$^6$\quad\quad \textbf{C.} 7x$^2$y$^3$\quad\quad \textbf{D.} 35x$^2$y$^3$\par\vspace{2pt}
\textbf{\color{primary} Câu 12.} Câu hỏi ôn tập lý thuyết và kỹ năng cơ bản (Chương I): Thương của phép chia đơn thức 72x$^4$y$^3$ cho đơn thức 2x$^2$y là:\par
\hspace*{1.5em}\textbf{A.} 36x$^2$y$^2$\quad\quad \textbf{B.} 48x$^2$y$^2$\quad\quad \textbf{C.} 36x$^6$y$^4$\quad\quad \textbf{D.} 36xy$^2$\par\vspace{2pt}
\textbf{\color{primary} Câu 13.} Theo nội dung sách giáo khoa Kết nối tri thức với cuộc sống (Chương I), bậc của đơn thức thu gọn A = -14x$^{14y}$$^3$z là:\par
\hspace*{1.5em}\textbf{A.} Bậc 14\quad\quad \textbf{B.} Bậc 17\quad\quad \textbf{C.} Bậc 42\quad\quad \textbf{D.} Bậc 18\par\vspace{2pt}
\textbf{\color{primary} Câu 14.} Tìm đáp số đúng của bài toán cơ bản sau: Thu gọn tổng hai đơn thức đồng dạng 17x$^2$y$^3$ + 27x$^2$y$^3$ ta được:\par
\hspace*{1.5em}\textbf{A.} 44x$^2$y$^3$\quad\quad \textbf{B.} 10x$^2$y$^3$\quad\quad \textbf{C.} 459x$^2$y$^3$\quad\quad \textbf{D.} 44x$^4$y$^6$\par\vspace{2pt}
\textbf{\color{primary} Câu 15.} Xét đối tượng toán học trong Chương I: Thương của phép chia đơn thức 90x$^4$y$^3$ cho đơn thức 2x$^2$y là. Giá trị hoặc tính chất tương ứng là:\par
\hspace*{1.5em}\textbf{A.} 45x$^2$y$^2$\quad\quad \textbf{B.} 45xy$^2$\quad\quad \textbf{C.} 45x$^6$y$^4$\quad\quad \textbf{D.} 60x$^2$y$^2$\par\vspace{2pt}
\textbf{\color{primary} Câu 16.} Hoàn thành phép xác định cơ bản sau đây: Bậc của đơn thức thu gọn A = -17x$^{17y}$$^3$z là:\par
\hspace*{1.5em}\textbf{A.} Bậc 17\quad\quad \textbf{B.} Bậc 51\quad\quad \textbf{C.} Bậc 20\quad\quad \textbf{D.} Bậc 21\par\vspace{2pt}
\textbf{\color{primary} Câu 17.} Trong các khẳng định hoặc giá trị dưới đây, đâu là đáp án đúng cho câu hỏi: Thu gọn tổng hai đơn thức đồng dạng 20x$^2$y$^3$ + 33x$^2$y$^3$ ta được?\par
\hspace*{1.5em}\textbf{A.} 53x$^2$y$^3$\quad\quad \textbf{B.} 13x$^2$y$^3$\quad\quad \textbf{C.} 53x$^4$y$^6$\quad\quad \textbf{D.} 660x$^2$y$^3$\par\vspace{2pt}
\textbf{\color{primary} Câu 18.} Bài tập khởi động Chương I: Thương của phép chia đơn thức 108x$^4$y$^3$ cho đơn thức 2x$^2$y là:\par
\hspace*{1.5em}\textbf{A.} 54xy$^2$\quad\quad \textbf{B.} 72x$^2$y$^2$\quad\quad \textbf{C.} 54x$^2$y$^2$\quad\quad \textbf{D.} 54x$^6$y$^4$\par\vspace{2pt}
\textbf{\color{primary} Câu 19.} Quan sát dữ kiện đề bài: Bậc của đơn thức thu gọn A = -20x$^{20y}$$^3$z là. Hãy chọn phương án đúng:\par
\hspace*{1.5em}\textbf{A.} Bậc 20\quad\quad \textbf{B.} Bậc 24\quad\quad \textbf{C.} Bậc 23\quad\quad \textbf{D.} Bậc 60\par\vspace{2pt}
\textbf{\color{primary} Câu 20.} Sử dụng công thức hoặc tính chất đầu bài của Chương I để giải quyết yêu cầu: Thu gọn tổng hai đơn thức đồng dạng 23x$^2$y$^3$ + 39x$^2$y$^3$ ta được:\par
\hspace*{1.5em}\textbf{A.} 62x$^4$y$^6$\quad\quad \textbf{B.} 16x$^2$y$^3$\quad\quad \textbf{C.} 62x$^2$y$^3$\quad\quad \textbf{D.} 897x$^2$y$^3$\par\vspace{2pt}
\subsection*{\color{accent} Phần II. Mức độ Trung bình (Thông hiểu -- Câu 21 đến Câu 40)}
\textbf{\color{primary} Câu 21.} Tính giá trị của đa thức P = x$^2$y - 2xy + 5 tại x = 2 và y = -1:\par
\hspace*{1.5em}\textbf{A.} 9\quad\quad \textbf{B.} 5\quad\quad \textbf{C.} -3\quad\quad \textbf{D.} 1\par\vspace{2pt}
\textbf{\color{primary} Câu 22.} Vận dụng quy tắc trong Chương I (Đa thức) để giải bài toán: Rút gọn biểu thức M = (x + 3)(x - 3) - x$^2$ ta được:\par
\hspace*{1.5em}\textbf{A.} -9\quad\quad \textbf{B.} 2x$^2$ + 0x - 9\quad\quad \textbf{C.} 6x - 9\quad\quad \textbf{D.} 0x - 9\par\vspace{2pt}
\textbf{\color{primary} Câu 23.} Biến đổi và tính toán theo kiến thức Chương I: Bậc của đa thức Q = 3x$^4$y - 4x$^2$y$^2$ + 5xy - 3x$^4$y + 7 sau khi đã thu gọn là. Kết quả thu được là:\par
\hspace*{1.5em}\textbf{A.} Bậc 4\quad\quad \textbf{B.} Bậc 0\quad\quad \textbf{C.} Bậc 2\quad\quad \textbf{D.} Bậc 5\par\vspace{2pt}
\textbf{\color{primary} Câu 24.} Xét bài toán mức độ thông hiểu sau: Tính giá trị của đa thức P = x$^2$y - 2xy + 8 tại x = 2 và y = -1. Đáp án đúng là:\par
\hspace*{1.5em}\textbf{A.} 12\quad\quad \textbf{B.} 8\quad\quad \textbf{C.} 4\quad\quad \textbf{D.} 0\par\vspace{2pt}
\textbf{\color{primary} Câu 25.} Để hoàn thành phiếu học tập Chương I, hãy xác định: Rút gọn biểu thức M = (x + 6)(x - 3) - x$^2$ ta được:\par
\hspace*{1.5em}\textbf{A.} 2x$^2$ + 3x - 18\quad\quad \textbf{B.} -18\quad\quad \textbf{C.} 3x - 18\quad\quad \textbf{D.} 9x - 18\par\vspace{2pt}
\textbf{\color{primary} Câu 26.} Phân tích dữ kiện bài toán: Bậc của đa thức Q = 3x$^4$y - 7x$^2$y$^2$ + 5xy - 3x$^4$y + 7 sau khi đã thu gọn là. Giá trị cần tìm bằng:\par
\hspace*{1.5em}\textbf{A.} Bậc 5\quad\quad \textbf{B.} Bậc 0\quad\quad \textbf{C.} Bậc 4\quad\quad \textbf{D.} Bậc 2\par\vspace{2pt}
\textbf{\color{primary} Câu 27.} Áp dụng tính chất đã học ở Chương I vào bài tập sau: Tính giá trị của đa thức P = x$^2$y - 2xy + 11 tại x = 2 và y = -1:\par
\hspace*{1.5em}\textbf{A.} 15\quad\quad \textbf{B.} 7\quad\quad \textbf{C.} 11\quad\quad \textbf{D.} 3\par\vspace{2pt}
\textbf{\color{primary} Câu 28.} Thực hiện các bước suy luận hợp lý để giải quyết yêu cầu: Rút gọn biểu thức M = (x + 9)(x - 3) - x$^2$ ta được:\par
\hspace*{1.5em}\textbf{A.} 12x - 27\quad\quad \textbf{B.} 6x - 27\quad\quad \textbf{C.} -27\quad\quad \textbf{D.} 2x$^2$ + 6x - 27\par\vspace{2pt}
\textbf{\color{primary} Câu 29.} Cho bài toán tự luyện (Đa thức): Bậc của đa thức Q = 3x$^4$y - 10x$^2$y$^2$ + 5xy - 3x$^4$y + 7 sau khi đã thu gọn là. Phương án nào sau đây chính xác?\par
\hspace*{1.5em}\textbf{A.} Bậc 0\quad\quad \textbf{B.} Bậc 5\quad\quad \textbf{C.} Bậc 2\quad\quad \textbf{D.} Bậc 4\par\vspace{2pt}
\textbf{\color{primary} Câu 30.} Kết hợp các phép biến đổi trong Chương I, hãy tìm kết quả của bài toán: Tính giá trị của đa thức P = x$^2$y - 2xy + 14 tại x = 2 và y = -1:\par
\hspace*{1.5em}\textbf{A.} 6\quad\quad \textbf{B.} 10\quad\quad \textbf{C.} 18\quad\quad \textbf{D.} 14\par\vspace{2pt}
\textbf{\color{primary} Câu 31.} Một bài kiểm tra định kỳ nội dung Chương I nêu câu hỏi: Rút gọn biểu thức M = (x + 12)(x - 3) - x$^2$ ta được. Hãy chọn đáp số đúng:\par
\hspace*{1.5em}\textbf{A.} -36\quad\quad \textbf{B.} 15x - 36\quad\quad \textbf{C.} 2x$^2$ + 9x - 36\quad\quad \textbf{D.} 9x - 36\par\vspace{2pt}
\textbf{\color{primary} Câu 32.} Tìm lời giải đúng cho bài toán thông hiểu: Bậc của đa thức Q = 3x$^4$y - 13x$^2$y$^2$ + 5xy - 3x$^4$y + 7 sau khi đã thu gọn là:\par
\hspace*{1.5em}\textbf{A.} Bậc 0\quad\quad \textbf{B.} Bậc 2\quad\quad \textbf{C.} Bậc 4\quad\quad \textbf{D.} Bậc 5\par\vspace{2pt}
\textbf{\color{primary} Câu 33.} Sử dụng mối liên hệ giữa các đại lượng trong Chương I: Tính giá trị của đa thức P = x$^2$y - 2xy + 17 tại x = 2 và y = -1:\par
\hspace*{1.5em}\textbf{A.} 9\quad\quad \textbf{B.} 17\quad\quad \textbf{C.} 21\quad\quad \textbf{D.} 13\par\vspace{2pt}
\textbf{\color{primary} Câu 34.} Hãy tính toán và chọn phương án đúng đối với bài tập: Rút gọn biểu thức M = (x + 15)(x - 3) - x$^2$ ta được:\par
\hspace*{1.5em}\textbf{A.} 2x$^2$ + 12x - 45\quad\quad \textbf{B.} 18x - 45\quad\quad \textbf{C.} -45\quad\quad \textbf{D.} 12x - 45\par\vspace{2pt}
\textbf{\color{primary} Câu 35.} Dựa trên các tính chất trọng tâm của Chương I, xét bài toán: Bậc của đa thức Q = 3x$^4$y - 16x$^2$y$^2$ + 5xy - 3x$^4$y + 7 sau khi đã thu gọn là:\par
\hspace*{1.5em}\textbf{A.} Bậc 2\quad\quad \textbf{B.} Bậc 4\quad\quad \textbf{C.} Bậc 0\quad\quad \textbf{D.} Bậc 5\par\vspace{2pt}
\textbf{\color{primary} Câu 36.} Thực hành giải toán (Đa thức): Tính giá trị của đa thức P = x$^2$y - 2xy + 20 tại x = 2 và y = -1. Đáp án nào dưới đây thỏa mãn?\par
\hspace*{1.5em}\textbf{A.} 20\quad\quad \textbf{B.} 12\quad\quad \textbf{C.} 16\quad\quad \textbf{D.} 24\par\vspace{2pt}
\textbf{\color{primary} Câu 37.} Kiểm tra kỹ năng biến đổi và suy luận qua bài toán: Rút gọn biểu thức M = (x + 18)(x - 3) - x$^2$ ta được:\par
\hspace*{1.5em}\textbf{A.} -54\quad\quad \textbf{B.} 2x$^2$ + 15x - 54\quad\quad \textbf{C.} 21x - 54\quad\quad \textbf{D.} 15x - 54\par\vspace{2pt}
\textbf{\color{primary} Câu 38.} Cho giả thiết thuộc phạm vi Chương I: Bậc của đa thức Q = 3x$^4$y - 19x$^2$y$^2$ + 5xy - 3x$^4$y + 7 sau khi đã thu gọn là. Kết luận đúng là:\par
\hspace*{1.5em}\textbf{A.} Bậc 5\quad\quad \textbf{B.} Bậc 4\quad\quad \textbf{C.} Bậc 2\quad\quad \textbf{D.} Bậc 0\par\vspace{2pt}
\textbf{\color{primary} Câu 39.} Rút gọn hoặc tính giá trị theo yêu cầu sau của Chương I: Tính giá trị của đa thức P = x$^2$y - 2xy + 23 tại x = 2 và y = -1:\par
\hspace*{1.5em}\textbf{A.} 27\quad\quad \textbf{B.} 23\quad\quad \textbf{C.} 19\quad\quad \textbf{D.} 15\par\vspace{2pt}
\textbf{\color{primary} Câu 40.} Xác định phương án đúng sau khi thực hiện phép tính / suy luận: Rút gọn biểu thức M = (x + 21)(x - 3) - x$^2$ ta được:\par
\hspace*{1.5em}\textbf{A.} 18x - 63\quad\quad \textbf{B.} 2x$^2$ + 18x - 63\quad\quad \textbf{C.} -63\quad\quad \textbf{D.} 24x - 63\par\vspace{2pt}
\subsection*{\color{accent} Phần III. Mức độ Khó (Vận dụng \& Vận dụng cao -- Câu 41 đến Câu 50)}
\textbf{\color{primary} Câu 41.} Tìm số tự nhiên n để đơn thức A = 5x$^{n+1}$y$^4$ chia hết cho đơn thức B = 2x$^3$y$^{n-1}$:\par
\hspace*{1.5em}\textbf{A.} n $\in$ \{2; 3; 4; 5\}\quad\quad \textbf{B.} n $\ge$ 2\quad\quad \textbf{C.} n $\in$ \{3; 4; 5\}\quad\quad \textbf{D.} n = 3\par\vspace{2pt}
\textbf{\color{primary} Câu 42.} Bài toán vận dụng (Chương I – Đa thức): Chứng minh giá trị của biểu thức T = x(2x - 3) - 2x(x - 1) + x + 13 không phụ thuộc vào biến x. Giá trị hằng số của T bằng. Kết quả đúng là:\par
\hspace*{1.5em}\textbf{A.} -13\quad\quad \textbf{B.} 15\quad\quad \textbf{C.} 13\quad\quad \textbf{D.} 0\par\vspace{2pt}
\textbf{\color{primary} Câu 43.} Vận dụng kiến thức nâng cao trong phạm vi Chương I để giải bài toán: Tìm số tự nhiên n để đơn thức A = 5x$^{n+1}$y$^4$ chia hết cho đơn thức B = 2x$^3$y$^{n-1}$:\par
\hspace*{1.5em}\textbf{A.} n $\ge$ 2\quad\quad \textbf{B.} n $\in$ \{3; 4; 5\}\quad\quad \textbf{C.} n $\in$ \{2; 3; 4; 5\}\quad\quad \textbf{D.} n = 3\par\vspace{2pt}
\textbf{\color{primary} Câu 44.} Phân tích và giải bài toán tổng hợp sau của Chương I: Chứng minh giá trị của biểu thức T = x(2x - 3) - 2x(x - 1) + x + 19 không phụ thuộc vào biến x. Giá trị hằng số của T bằng:\par
\hspace*{1.5em}\textbf{A.} 21\quad\quad \textbf{B.} 19\quad\quad \textbf{C.} -19\quad\quad \textbf{D.} 0\par\vspace{2pt}
\textbf{\color{primary} Câu 45.} Thử thách tư duy trong phạm vi kiến thức Chương I: Tìm số tự nhiên n để đơn thức A = 5x$^{n+1}$y$^4$ chia hết cho đơn thức B = 2x$^3$y$^{n-1}$. Đáp số chính xác là:\par
\hspace*{1.5em}\textbf{A.} n $\in$ \{3; 4; 5\}\quad\quad \textbf{B.} n = 3\quad\quad \textbf{C.} n $\in$ \{2; 3; 4; 5\}\quad\quad \textbf{D.} n $\ge$ 2\par\vspace{2pt}
\textbf{\color{primary} Câu 46.} Sử dụng phương pháp giải toán của Chương I để tìm đáp án cho bài toán: Chứng minh giá trị của biểu thức T = x(2x - 3) - 2x(x - 1) + x + 25 không phụ thuộc vào biến x. Giá trị hằng số của T bằng:\par
\hspace*{1.5em}\textbf{A.} -25\quad\quad \textbf{B.} 0\quad\quad \textbf{C.} 25\quad\quad \textbf{D.} 27\par\vspace{2pt}
\textbf{\color{primary} Câu 47.} Xét bài toán vận dụng thực tiễn / nâng cao thuộc Chương I: Tìm số tự nhiên n để đơn thức A = 5x$^{n+1}$y$^4$ chia hết cho đơn thức B = 2x$^3$y$^{n-1}$:\par
\hspace*{1.5em}\textbf{A.} n $\ge$ 2\quad\quad \textbf{B.} n $\in$ \{3; 4; 5\}\quad\quad \textbf{C.} n $\in$ \{2; 3; 4; 5\}\quad\quad \textbf{D.} n = 3\par\vspace{2pt}
\textbf{\color{primary} Câu 48.} Tìm kết quả cuối cùng của bài toán phân loại học sinh khá giỏi (Chương I): Chứng minh giá trị của biểu thức T = x(2x - 3) - 2x(x - 1) + x + 31 không phụ thuộc vào biến x. Giá trị hằng số của T bằng:\par
\hspace*{1.5em}\textbf{A.} 31\quad\quad \textbf{B.} 0\quad\quad \textbf{C.} 33\quad\quad \textbf{D.} -31\par\vspace{2pt}
\textbf{\color{primary} Câu 49.} Kết hợp các tính chất đã học trong Chương I, hãy giải quyết bài toán: Tìm số tự nhiên n để đơn thức A = 5x$^{n+1}$y$^4$ chia hết cho đơn thức B = 2x$^3$y$^{n-1}$:\par
\hspace*{1.5em}\textbf{A.} n $\ge$ 2\quad\quad \textbf{B.} n = 3\quad\quad \textbf{C.} n $\in$ \{2; 3; 4; 5\}\quad\quad \textbf{D.} n $\in$ \{3; 4; 5\}\par\vspace{2pt}
\textbf{\color{primary} Câu 50.} Câu hỏi vận dụng cao cuối Chương I: Chứng minh giá trị của biểu thức T = x(2x - 3) - 2x(x - 1) + x + 37 không phụ thuộc vào biến x. Giá trị hằng số của T bằng. Phương án đúng là:\par
\hspace*{1.5em}\textbf{A.} -37\quad\quad \textbf{B.} 37\quad\quad \textbf{C.} 39\quad\quad \textbf{D.} 0\par\vspace{2pt}
\subsection*{\color{primary} Bảng đáp án Chương I: Đa thức}
\begin{center}
\renewcommand{\arraystretch}{1.25}
\begin{tabularx}{\textwidth}{*{10}{>{\centering\arraybackslash}X}}
\toprule
\textbf{\color{primary} Cột 1} & \textbf{\color{primary} Cột 2} & \textbf{\color{primary} Cột 3} & \textbf{\color{primary} Cột 4} & \textbf{\color{primary} Cột 5} & \textbf{\color{primary} Cột 6} & \textbf{\color{primary} Cột 7} & \textbf{\color{primary} Cột 8} & \textbf{\color{primary} Cột 9} & \textbf{\color{primary} Cột 10} \\
\midrule
\textbf{1.B} & \textbf{2.B} & \textbf{3.A} & \textbf{4.D} & \textbf{5.B} & \textbf{6.B} & \textbf{7.A} & \textbf{8.B} & \textbf{9.B} & \textbf{10.B} \\
\rowcolor{rowtint} \textbf{11.D} & \textbf{12.A} & \textbf{13.D} & \textbf{14.A} & \textbf{15.A} & \textbf{16.D} & \textbf{17.A} & \textbf{18.C} & \textbf{19.B} & \textbf{20.C} \\
\textbf{21.B} & \textbf{22.D} & \textbf{23.A} & \textbf{24.B} & \textbf{25.C} & \textbf{26.C} & \textbf{27.C} & \textbf{28.B} & \textbf{29.D} & \textbf{30.D} \\
\rowcolor{rowtint} \textbf{31.D} & \textbf{32.C} & \textbf{33.B} & \textbf{34.D} & \textbf{35.B} & \textbf{36.A} & \textbf{37.D} & \textbf{38.B} & \textbf{39.B} & \textbf{40.A} \\
\textbf{41.A} & \textbf{42.C} & \textbf{43.C} & \textbf{44.B} & \textbf{45.C} & \textbf{46.C} & \textbf{47.C} & \textbf{48.A} & \textbf{49.C} & \textbf{50.B} \\
\bottomrule
\end{tabularx}
\end{center}\vspace{10pt}
\section*{\color{primary} CHƯƠNG II: HẰNG ĐẲNG THỨC ĐÁNG NHỚ VÀ ỨNG DỤNG}
\textit{Trọng tâm kiến thức SGK: Bảy hằng đẳng thức đáng nhớ (bình phương của một tổng/hiệu, hiệu hai bình phương, lập phương của một tổng/hiệu, tổng/hiệu hai lập phương) và phân tích đa thức thành nhân tử.}\par\vspace{4pt}
\subsection*{\color{accent} Phần I. Mức độ Dễ (Nhận biết -- Câu 1 đến Câu 20)}
\textbf{\color{primary} Câu 1.} Khai triển hằng đẳng thức (x + 2)$^2$ ta được:\par
\hspace*{1.5em}\textbf{A.} x$^2$ + 2x + 4\quad\quad \textbf{B.} x$^2$ + 4\quad\quad \textbf{C.} x$^2$ - 4x + 4\quad\quad \textbf{D.} x$^2$ + 4x + 4\par\vspace{2pt}
\textbf{\color{primary} Câu 2.} Thực hiện yêu cầu nhận biết sau: Viết biểu thức x$^2$ - 16 dưới dạng tích ta được.\par
\hspace*{1.5em}\textbf{A.} (x + 4)$^2$\quad\quad \textbf{B.} (x - 4)$^2$\quad\quad \textbf{C.} (x - 4)(x + 4)\quad\quad \textbf{D.} (x - 16)(x + 1)\par\vspace{2pt}
\textbf{\color{primary} Câu 3.} Trong bài học thuộc Chương II, xét nội dung: Phân tích đa thức 4x$^2$ - 8x thành nhân tử ta được. Đáp án đúng là:\par
\hspace*{1.5em}\textbf{A.} 4x(x - 2)\quad\quad \textbf{B.} 4x(x + 2)\quad\quad \textbf{C.} x(4x - 2)\quad\quad \textbf{D.} 4(x$^2$ - 2)\par\vspace{2pt}
\textbf{\color{primary} Câu 4.} Dựa vào kiến thức chuẩn SGK của Chương II (Hằng đẳng thức đáng nhớ và ứng dụng), khai triển hằng đẳng thức (x + 5)$^2$ ta được:\par
\hspace*{1.5em}\textbf{A.} x$^2$ + 25\quad\quad \textbf{B.} x$^2$ + 10x + 25\quad\quad \textbf{C.} x$^2$ + 5x + 25\quad\quad \textbf{D.} x$^2$ - 10x + 25\par\vspace{2pt}
\textbf{\color{primary} Câu 5.} Chọn kết quả đúng nhất khi xét: Viết biểu thức x$^2$ - 49 dưới dạng tích ta được:\par
\hspace*{1.5em}\textbf{A.} (x - 7)(x + 7)\quad\quad \textbf{B.} (x - 49)(x + 1)\quad\quad \textbf{C.} (x + 7)$^2$\quad\quad \textbf{D.} (x - 7)$^2$\par\vspace{2pt}
\textbf{\color{primary} Câu 6.} Khi làm bài tập cơ bản về hằng đẳng thức đáng nhớ và ứng dụng, một học sinh gặp câu hỏi: Phân tích đa thức 7x$^2$ - 14x thành nhân tử ta được. Kết quả nào dưới đây đúng?\par
\hspace*{1.5em}\textbf{A.} x(7x - 2)\quad\quad \textbf{B.} 7(x$^2$ - 2)\quad\quad \textbf{C.} 7x(x + 2)\quad\quad \textbf{D.} 7x(x - 2)\par\vspace{2pt}
\textbf{\color{primary} Câu 7.} Xác định giá trị hoặc khẳng định đúng cho bài toán sau: Khai triển hằng đẳng thức (x + 8)$^2$ ta được:\par
\hspace*{1.5em}\textbf{A.} x$^2$ + 8x + 64\quad\quad \textbf{B.} x$^2$ + 16x + 64\quad\quad \textbf{C.} x$^2$ - 16x + 64\quad\quad \textbf{D.} x$^2$ + 64\par\vspace{2pt}
\textbf{\color{primary} Câu 8.} Áp dụng trực tiếp định nghĩa và quy tắc trong Chương II: Viết biểu thức x$^2$ - 100 dưới dạng tích ta được:\par
\hspace*{1.5em}\textbf{A.} (x - 10)$^2$\quad\quad \textbf{B.} (x - 10)(x + 10)\quad\quad \textbf{C.} (x - 100)(x + 1)\quad\quad \textbf{D.} (x + 10)$^2$\par\vspace{2pt}
\textbf{\color{primary} Câu 9.} Kiểm tra nhanh kiến thức nền tảng: Phân tích đa thức 10x$^2$ - 20x thành nhân tử ta được:\par
\hspace*{1.5em}\textbf{A.} 10x(x + 2)\quad\quad \textbf{B.} 10(x$^2$ - 2)\quad\quad \textbf{C.} 10x(x - 2)\quad\quad \textbf{D.} x(10x - 2)\par\vspace{2pt}
\textbf{\color{primary} Câu 10.} Cho dữ kiện nhận biết sau đây: Khai triển hằng đẳng thức (x + 11)$^2$ ta được. Phương án chính xác là:\par
\hspace*{1.5em}\textbf{A.} x$^2$ - 22x + 121\quad\quad \textbf{B.} x$^2$ + 22x + 121\quad\quad \textbf{C.} x$^2$ + 11x + 121\quad\quad \textbf{D.} x$^2$ + 121\par\vspace{2pt}
\textbf{\color{primary} Câu 11.} Hãy chỉ ra kết quả đúng trong bốn lựa chọn A, B, C, D đối với yêu cầu: Viết biểu thức x$^2$ - 169 dưới dạng tích ta được:\par
\hspace*{1.5em}\textbf{A.} (x - 13)$^2$\quad\quad \textbf{B.} (x - 169)(x + 1)\quad\quad \textbf{C.} (x + 13)$^2$\quad\quad \textbf{D.} (x - 13)(x + 13)\par\vspace{2pt}
\textbf{\color{primary} Câu 12.} Câu hỏi ôn tập lý thuyết và kỹ năng cơ bản (Chương II): Phân tích đa thức 13x$^2$ - 26x thành nhân tử ta được:\par
\hspace*{1.5em}\textbf{A.} 13x(x + 2)\quad\quad \textbf{B.} 13(x$^2$ - 2)\quad\quad \textbf{C.} x(13x - 2)\quad\quad \textbf{D.} 13x(x - 2)\par\vspace{2pt}
\textbf{\color{primary} Câu 13.} Theo nội dung sách giáo khoa Kết nối tri thức với cuộc sống (Chương II), khai triển hằng đẳng thức (x + 14)$^2$ ta được:\par
\hspace*{1.5em}\textbf{A.} x$^2$ + 14x + 196\quad\quad \textbf{B.} x$^2$ + 28x + 196\quad\quad \textbf{C.} x$^2$ - 28x + 196\quad\quad \textbf{D.} x$^2$ + 196\par\vspace{2pt}
\textbf{\color{primary} Câu 14.} Tìm đáp số đúng của bài toán cơ bản sau: Viết biểu thức x$^2$ - 256 dưới dạng tích ta được:\par
\hspace*{1.5em}\textbf{A.} (x - 16)$^2$\quad\quad \textbf{B.} (x - 16)(x + 16)\quad\quad \textbf{C.} (x - 256)(x + 1)\quad\quad \textbf{D.} (x + 16)$^2$\par\vspace{2pt}
\textbf{\color{primary} Câu 15.} Xét đối tượng toán học trong Chương II: Phân tích đa thức 16x$^2$ - 32x thành nhân tử ta được. Giá trị hoặc tính chất tương ứng là:\par
\hspace*{1.5em}\textbf{A.} 16x(x + 2)\quad\quad \textbf{B.} 16x(x - 2)\quad\quad \textbf{C.} x(16x - 2)\quad\quad \textbf{D.} 16(x$^2$ - 2)\par\vspace{2pt}
\textbf{\color{primary} Câu 16.} Hoàn thành phép xác định cơ bản sau đây: Khai triển hằng đẳng thức (x + 17)$^2$ ta được:\par
\hspace*{1.5em}\textbf{A.} x$^2$ + 17x + 289\quad\quad \textbf{B.} x$^2$ - 34x + 289\quad\quad \textbf{C.} x$^2$ + 289\quad\quad \textbf{D.} x$^2$ + 34x + 289\par\vspace{2pt}
\textbf{\color{primary} Câu 17.} Trong các khẳng định hoặc giá trị dưới đây, đâu là đáp án đúng cho câu hỏi: Viết biểu thức x$^2$ - 361 dưới dạng tích ta được?\par
\hspace*{1.5em}\textbf{A.} (x - 19)(x + 19)\quad\quad \textbf{B.} (x + 19)$^2$\quad\quad \textbf{C.} (x - 361)(x + 1)\quad\quad \textbf{D.} (x - 19)$^2$\par\vspace{2pt}
\textbf{\color{primary} Câu 18.} Bài tập khởi động Chương II: Phân tích đa thức 19x$^2$ - 38x thành nhân tử ta được:\par
\hspace*{1.5em}\textbf{A.} 19x(x - 2)\quad\quad \textbf{B.} 19(x$^2$ - 2)\quad\quad \textbf{C.} x(19x - 2)\quad\quad \textbf{D.} 19x(x + 2)\par\vspace{2pt}
\textbf{\color{primary} Câu 19.} Quan sát dữ kiện đề bài: Khai triển hằng đẳng thức (x + 20)$^2$ ta được. Hãy chọn phương án đúng:\par
\hspace*{1.5em}\textbf{A.} x$^2$ + 40x + 400\quad\quad \textbf{B.} x$^2$ - 40x + 400\quad\quad \textbf{C.} x$^2$ + 400\quad\quad \textbf{D.} x$^2$ + 20x + 400\par\vspace{2pt}
\textbf{\color{primary} Câu 20.} Sử dụng công thức hoặc tính chất đầu bài của Chương II để giải quyết yêu cầu: Viết biểu thức x$^2$ - 484 dưới dạng tích ta được:\par
\hspace*{1.5em}\textbf{A.} (x - 484)(x + 1)\quad\quad \textbf{B.} (x - 22)$^2$\quad\quad \textbf{C.} (x - 22)(x + 22)\quad\quad \textbf{D.} (x + 22)$^2$\par\vspace{2pt}
\subsection*{\color{accent} Phần II. Mức độ Trung bình (Thông hiểu -- Câu 21 đến Câu 40)}
\textbf{\color{primary} Câu 21.} Tính nhanh giá trị của biểu thức P = x$^2$ - 6x + 9 tại x = 103:\par
\hspace*{1.5em}\textbf{A.} 1000\quad\quad \textbf{B.} 10200\quad\quad \textbf{C.} 100\quad\quad \textbf{D.} 10000\par\vspace{2pt}
\textbf{\color{primary} Câu 22.} Vận dụng quy tắc trong Chương II (Hằng đẳng thức đáng nhớ và ứng dụng) để giải bài toán: Tính nhanh hiệu hai bình phương: 56$^2$ - 46$^2$ bằng:\par
\hspace*{1.5em}\textbf{A.} 1020\quad\quad \textbf{B.} 102\quad\quad \textbf{C.} 2040\quad\quad \textbf{D.} 100\par\vspace{2pt}
\textbf{\color{primary} Câu 23.} Biến đổi và tính toán theo kiến thức Chương II: Phân tích đa thức x$^2$ - 12x + 36 - y$^2$ thành nhân tử ta được. Kết quả thu được là:\par
\hspace*{1.5em}\textbf{A.} (x - y)$^2$ - 36\quad\quad \textbf{B.} (x - 6 - y)(x - 6 + y)\quad\quad \textbf{C.} (x + 6 - y)(x + 6 + y)\quad\quad \textbf{D.} (x - 6)(x + y)\par\vspace{2pt}
\textbf{\color{primary} Câu 24.} Xét bài toán mức độ thông hiểu sau: Tính nhanh giá trị của biểu thức P = x$^2$ - 12x + 36 tại x = 106. Đáp án đúng là:\par
\hspace*{1.5em}\textbf{A.} 1000\quad\quad \textbf{B.} 100\quad\quad \textbf{C.} 10000\quad\quad \textbf{D.} 10200\par\vspace{2pt}
\textbf{\color{primary} Câu 25.} Để hoàn thành phiếu học tập Chương II, hãy xác định: Tính nhanh hiệu hai bình phương: 59$^2$ - 49$^2$ bằng:\par
\hspace*{1.5em}\textbf{A.} 2160\quad\quad \textbf{B.} 100\quad\quad \textbf{C.} 1080\quad\quad \textbf{D.} 108\par\vspace{2pt}
\textbf{\color{primary} Câu 26.} Phân tích dữ kiện bài toán: Phân tích đa thức x$^2$ - 18x + 81 - y$^2$ thành nhân tử ta được. Giá trị cần tìm bằng:\par
\hspace*{1.5em}\textbf{A.} (x - y)$^2$ - 81\quad\quad \textbf{B.} (x + 9 - y)(x + 9 + y)\quad\quad \textbf{C.} (x - 9)(x + y)\quad\quad \textbf{D.} (x - 9 - y)(x - 9 + y)\par\vspace{2pt}
\textbf{\color{primary} Câu 27.} Áp dụng tính chất đã học ở Chương II vào bài tập sau: Tính nhanh giá trị của biểu thức P = x$^2$ - 18x + 81 tại x = 109:\par
\hspace*{1.5em}\textbf{A.} 1000\quad\quad \textbf{B.} 100\quad\quad \textbf{C.} 10000\quad\quad \textbf{D.} 10200\par\vspace{2pt}
\textbf{\color{primary} Câu 28.} Thực hiện các bước suy luận hợp lý để giải quyết yêu cầu: Tính nhanh hiệu hai bình phương: 62$^2$ - 52$^2$ bằng:\par
\hspace*{1.5em}\textbf{A.} 100\quad\quad \textbf{B.} 1140\quad\quad \textbf{C.} 2280\quad\quad \textbf{D.} 114\par\vspace{2pt}
\textbf{\color{primary} Câu 29.} Cho bài toán tự luyện (Hằng đẳng thức đáng nhớ và ứng dụng): Phân tích đa thức x$^2$ - 24x + 144 - y$^2$ thành nhân tử ta được. Phương án nào sau đây chính xác?\par
\hspace*{1.5em}\textbf{A.} (x - y)$^2$ - 144\quad\quad \textbf{B.} (x + 12 - y)(x + 12 + y)\quad\quad \textbf{C.} (x - 12 - y)(x - 12 + y)\quad\quad \textbf{D.} (x - 12)(x + y)\par\vspace{2pt}
\textbf{\color{primary} Câu 30.} Kết hợp các phép biến đổi trong Chương II, hãy tìm kết quả của bài toán: Tính nhanh giá trị của biểu thức P = x$^2$ - 24x + 144 tại x = 112:\par
\hspace*{1.5em}\textbf{A.} 10000\quad\quad \textbf{B.} 1000\quad\quad \textbf{C.} 100\quad\quad \textbf{D.} 10200\par\vspace{2pt}
\textbf{\color{primary} Câu 31.} Một bài kiểm tra định kỳ nội dung Chương II nêu câu hỏi: Tính nhanh hiệu hai bình phương: 65$^2$ - 55$^2$ bằng. Hãy chọn đáp số đúng:\par
\hspace*{1.5em}\textbf{A.} 120\quad\quad \textbf{B.} 2400\quad\quad \textbf{C.} 1200\quad\quad \textbf{D.} 100\par\vspace{2pt}
\textbf{\color{primary} Câu 32.} Tìm lời giải đúng cho bài toán thông hiểu: Phân tích đa thức x$^2$ - 30x + 225 - y$^2$ thành nhân tử ta được:\par
\hspace*{1.5em}\textbf{A.} (x + 15 - y)(x + 15 + y)\quad\quad \textbf{B.} (x - 15)(x + y)\quad\quad \textbf{C.} (x - y)$^2$ - 225\quad\quad \textbf{D.} (x - 15 - y)(x - 15 + y)\par\vspace{2pt}
\textbf{\color{primary} Câu 33.} Sử dụng mối liên hệ giữa các đại lượng trong Chương II: Tính nhanh giá trị của biểu thức P = x$^2$ - 30x + 225 tại x = 115:\par
\hspace*{1.5em}\textbf{A.} 1000\quad\quad \textbf{B.} 100\quad\quad \textbf{C.} 10200\quad\quad \textbf{D.} 10000\par\vspace{2pt}
\textbf{\color{primary} Câu 34.} Hãy tính toán và chọn phương án đúng đối với bài tập: Tính nhanh hiệu hai bình phương: 68$^2$ - 58$^2$ bằng:\par
\hspace*{1.5em}\textbf{A.} 100\quad\quad \textbf{B.} 2520\quad\quad \textbf{C.} 1260\quad\quad \textbf{D.} 126\par\vspace{2pt}
\textbf{\color{primary} Câu 35.} Dựa trên các tính chất trọng tâm của Chương II, xét bài toán: Phân tích đa thức x$^2$ - 36x + 324 - y$^2$ thành nhân tử ta được:\par
\hspace*{1.5em}\textbf{A.} (x - 18 - y)(x - 18 + y)\quad\quad \textbf{B.} (x + 18 - y)(x + 18 + y)\quad\quad \textbf{C.} (x - 18)(x + y)\quad\quad \textbf{D.} (x - y)$^2$ - 324\par\vspace{2pt}
\textbf{\color{primary} Câu 36.} Thực hành giải toán (Hằng đẳng thức đáng nhớ và ứng dụng): Tính nhanh giá trị của biểu thức P = x$^2$ - 36x + 324 tại x = 118. Đáp án nào dưới đây thỏa mãn?\par
\hspace*{1.5em}\textbf{A.} 1000\quad\quad \textbf{B.} 100\quad\quad \textbf{C.} 10200\quad\quad \textbf{D.} 10000\par\vspace{2pt}
\textbf{\color{primary} Câu 37.} Kiểm tra kỹ năng biến đổi và suy luận qua bài toán: Tính nhanh hiệu hai bình phương: 71$^2$ - 61$^2$ bằng:\par
\hspace*{1.5em}\textbf{A.} 1320\quad\quad \textbf{B.} 2640\quad\quad \textbf{C.} 132\quad\quad \textbf{D.} 100\par\vspace{2pt}
\textbf{\color{primary} Câu 38.} Cho giả thiết thuộc phạm vi Chương II: Phân tích đa thức x$^2$ - 42x + 441 - y$^2$ thành nhân tử ta được. Kết luận đúng là:\par
\hspace*{1.5em}\textbf{A.} (x + 21 - y)(x + 21 + y)\quad\quad \textbf{B.} (x - 21)(x + y)\quad\quad \textbf{C.} (x - y)$^2$ - 441\quad\quad \textbf{D.} (x - 21 - y)(x - 21 + y)\par\vspace{2pt}
\textbf{\color{primary} Câu 39.} Rút gọn hoặc tính giá trị theo yêu cầu sau của Chương II: Tính nhanh giá trị của biểu thức P = x$^2$ - 42x + 441 tại x = 121:\par
\hspace*{1.5em}\textbf{A.} 10200\quad\quad \textbf{B.} 10000\quad\quad \textbf{C.} 1000\quad\quad \textbf{D.} 100\par\vspace{2pt}
\textbf{\color{primary} Câu 40.} Xác định phương án đúng sau khi thực hiện phép tính / suy luận: Tính nhanh hiệu hai bình phương: 74$^2$ - 64$^2$ bằng:\par
\hspace*{1.5em}\textbf{A.} 1380\quad\quad \textbf{B.} 138\quad\quad \textbf{C.} 2760\quad\quad \textbf{D.} 100\par\vspace{2pt}
\subsection*{\color{accent} Phần III. Mức độ Khó (Vận dụng \& Vận dụng cao -- Câu 41 đến Câu 50)}
\textbf{\color{primary} Câu 41.} Tìm giá trị nhỏ nhất của đa thức A = x$^2$ - 6x + 14.:\par
\hspace*{1.5em}\textbf{A.} 5\quad\quad \textbf{B.} -5\quad\quad \textbf{C.} 3\quad\quad \textbf{D.} 14\par\vspace{2pt}
\textbf{\color{primary} Câu 42.} Bài toán vận dụng (Chương II – Hằng đẳng thức đáng nhớ và ứng dụng): Cho x + y = 5 và x$\cdot$y = 2. Giá trị của biểu thức S = x$^2$ + y$^2$ bằng. Kết quả đúng là:\par
\hspace*{1.5em}\textbf{A.} 21\quad\quad \textbf{B.} 29\quad\quad \textbf{C.} 25\quad\quad \textbf{D.} 23\par\vspace{2pt}
\textbf{\color{primary} Câu 43.} Vận dụng kiến thức nâng cao trong phạm vi Chương II để giải bài toán: Tìm giá trị nhỏ nhất của đa thức A = x$^2$ - 10x + 32.:\par
\hspace*{1.5em}\textbf{A.} 7\quad\quad \textbf{B.} 5\quad\quad \textbf{C.} -7\quad\quad \textbf{D.} 32\par\vspace{2pt}
\textbf{\color{primary} Câu 44.} Phân tích và giải bài toán tổng hợp sau của Chương II: Cho x + y = 7 và x$\cdot$y = 2. Giá trị của biểu thức S = x$^2$ + y$^2$ bằng:\par
\hspace*{1.5em}\textbf{A.} 47\quad\quad \textbf{B.} 49\quad\quad \textbf{C.} 45\quad\quad \textbf{D.} 53\par\vspace{2pt}
\textbf{\color{primary} Câu 45.} Thử thách tư duy trong phạm vi kiến thức Chương II: Tìm giá trị nhỏ nhất của đa thức A = x$^2$ - 14x + 58.. Đáp số chính xác là:\par
\hspace*{1.5em}\textbf{A.} 9\quad\quad \textbf{B.} -9\quad\quad \textbf{C.} 7\quad\quad \textbf{D.} 58\par\vspace{2pt}
\textbf{\color{primary} Câu 46.} Sử dụng phương pháp giải toán của Chương II để tìm đáp án cho bài toán: Cho x + y = 9 và x$\cdot$y = 2. Giá trị của biểu thức S = x$^2$ + y$^2$ bằng:\par
\hspace*{1.5em}\textbf{A.} 77\quad\quad \textbf{B.} 85\quad\quad \textbf{C.} 81\quad\quad \textbf{D.} 79\par\vspace{2pt}
\textbf{\color{primary} Câu 47.} Xét bài toán vận dụng thực tiễn / nâng cao thuộc Chương II: Tìm giá trị nhỏ nhất của đa thức A = x$^2$ - 18x + 92.:\par
\hspace*{1.5em}\textbf{A.} 9\quad\quad \textbf{B.} -11\quad\quad \textbf{C.} 92\quad\quad \textbf{D.} 11\par\vspace{2pt}
\textbf{\color{primary} Câu 48.} Tìm kết quả cuối cùng của bài toán phân loại học sinh khá giỏi (Chương II): Cho x + y = 11 và x$\cdot$y = 2. Giá trị của biểu thức S = x$^2$ + y$^2$ bằng:\par
\hspace*{1.5em}\textbf{A.} 125\quad\quad \textbf{B.} 119\quad\quad \textbf{C.} 117\quad\quad \textbf{D.} 121\par\vspace{2pt}
\textbf{\color{primary} Câu 49.} Kết hợp các tính chất đã học trong Chương II, hãy giải quyết bài toán: Tìm giá trị nhỏ nhất của đa thức A = x$^2$ - 22x + 134.:\par
\hspace*{1.5em}\textbf{A.} -13\quad\quad \textbf{B.} 134\quad\quad \textbf{C.} 13\quad\quad \textbf{D.} 11\par\vspace{2pt}
\textbf{\color{primary} Câu 50.} Câu hỏi vận dụng cao cuối Chương II: Cho x + y = 13 và x$\cdot$y = 2. Giá trị của biểu thức S = x$^2$ + y$^2$ bằng. Phương án đúng là:\par
\hspace*{1.5em}\textbf{A.} 167\quad\quad \textbf{B.} 173\quad\quad \textbf{C.} 169\quad\quad \textbf{D.} 165\par\vspace{2pt}
\subsection*{\color{primary} Bảng đáp án Chương II: Hằng đẳng thức đáng nhớ và ứng dụng}
\begin{center}
\renewcommand{\arraystretch}{1.25}
\begin{tabularx}{\textwidth}{*{10}{>{\centering\arraybackslash}X}}
\toprule
\textbf{\color{primary} Cột 1} & \textbf{\color{primary} Cột 2} & \textbf{\color{primary} Cột 3} & \textbf{\color{primary} Cột 4} & \textbf{\color{primary} Cột 5} & \textbf{\color{primary} Cột 6} & \textbf{\color{primary} Cột 7} & \textbf{\color{primary} Cột 8} & \textbf{\color{primary} Cột 9} & \textbf{\color{primary} Cột 10} \\
\midrule
\textbf{1.D} & \textbf{2.C} & \textbf{3.A} & \textbf{4.B} & \textbf{5.A} & \textbf{6.D} & \textbf{7.B} & \textbf{8.B} & \textbf{9.C} & \textbf{10.B} \\
\rowcolor{rowtint} \textbf{11.D} & \textbf{12.D} & \textbf{13.B} & \textbf{14.B} & \textbf{15.B} & \textbf{16.D} & \textbf{17.A} & \textbf{18.A} & \textbf{19.A} & \textbf{20.C} \\
\textbf{21.D} & \textbf{22.A} & \textbf{23.B} & \textbf{24.C} & \textbf{25.C} & \textbf{26.D} & \textbf{27.C} & \textbf{28.B} & \textbf{29.C} & \textbf{30.A} \\
\rowcolor{rowtint} \textbf{31.C} & \textbf{32.D} & \textbf{33.D} & \textbf{34.C} & \textbf{35.A} & \textbf{36.D} & \textbf{37.A} & \textbf{38.D} & \textbf{39.B} & \textbf{40.A} \\
\textbf{41.A} & \textbf{42.A} & \textbf{43.A} & \textbf{44.C} & \textbf{45.A} & \textbf{46.A} & \textbf{47.D} & \textbf{48.C} & \textbf{49.C} & \textbf{50.D} \\
\bottomrule
\end{tabularx}
\end{center}\vspace{10pt}
\section*{\color{primary} CHƯƠNG III: TỨ GIÁC}
\textit{Trọng tâm kiến thức SGK: Tổng các góc của một tứ giác, hình thang và hình thang cân, hình bình hành, hình chữ nhật, hình thoi, hình vuông (tính chất và dấu hiệu nhận biết).}\par\vspace{4pt}
\subsection*{\color{accent} Phần I. Mức độ Dễ (Nhận biết -- Câu 1 đến Câu 20)}
\textbf{\color{primary} Câu 1.} Cho tứ giác ABCD có góc A = 70°, góc B = 95°, góc C = 105°. Số đo góc D bằng:\par
\hspace*{1.5em}\textbf{A.} 90°\quad\quad \textbf{B.} 80°\quad\quad \textbf{C.} 100°\quad\quad \textbf{D.} 110°\par\vspace{2pt}
\textbf{\color{primary} Câu 2.} Hãy trả lời câu hỏi cơ bản sau: Hình bình hành có hai đường chéo bằng nhau là hình gì?\par
\hspace*{1.5em}\textbf{A.} Hình chữ nhật\quad\quad \textbf{B.} Hình tứ giác thường\quad\quad \textbf{C.} Hình thoi\quad\quad \textbf{D.} Hình thang vuông\par\vspace{2pt}
\textbf{\color{primary} Câu 3.} Trong bài học thuộc Chương III, xét nội dung: Hình bình hành có hai đường chéo vuông góc với nhau là hình gì. Đáp án đúng là:\par
\hspace*{1.5em}\textbf{A.} Hình thang cân\quad\quad \textbf{B.} Hình thoi\quad\quad \textbf{C.} Hình chữ nhật\quad\quad \textbf{D.} Hình thang vuông\par\vspace{2pt}
\textbf{\color{primary} Câu 4.} Dựa vào kiến thức chuẩn SGK của Chương III (Tứ giác), cho tứ giác ABCD có góc A = 73°, góc B = 95°, góc C = 102°. Số đo góc D bằng:\par
\hspace*{1.5em}\textbf{A.} 100°\quad\quad \textbf{B.} 110°\quad\quad \textbf{C.} 80°\quad\quad \textbf{D.} 90°\par\vspace{2pt}
\textbf{\color{primary} Câu 5.} Chọn phương án trả lời đúng nhất: Hình bình hành có hai đường chéo bằng nhau là hình gì?\par
\hspace*{1.5em}\textbf{A.} Hình chữ nhật\quad\quad \textbf{B.} Hình thoi\quad\quad \textbf{C.} Hình thang vuông\quad\quad \textbf{D.} Hình tứ giác thường\par\vspace{2pt}
\textbf{\color{primary} Câu 6.} Khi làm bài tập cơ bản về tứ giác, một học sinh gặp câu hỏi: Hình bình hành có hai đường chéo vuông góc với nhau là hình gì. Kết quả nào dưới đây đúng?\par
\hspace*{1.5em}\textbf{A.} Hình thoi\quad\quad \textbf{B.} Hình thang cân\quad\quad \textbf{C.} Hình chữ nhật\quad\quad \textbf{D.} Hình thang vuông\par\vspace{2pt}
\textbf{\color{primary} Câu 7.} Xác định giá trị hoặc khẳng định đúng cho bài toán sau: Cho tứ giác ABCD có góc A = 76°, góc B = 95°, góc C = 99°. Số đo góc D bằng:\par
\hspace*{1.5em}\textbf{A.} 110°\quad\quad \textbf{B.} 80°\quad\quad \textbf{C.} 100°\quad\quad \textbf{D.} 90°\par\vspace{2pt}
\textbf{\color{primary} Câu 8.} Áp dụng trực tiếp định nghĩa và quy tắc trong Chương III: Hình bình hành có hai đường chéo bằng nhau là hình gì:\par
\hspace*{1.5em}\textbf{A.} Hình chữ nhật\quad\quad \textbf{B.} Hình thoi\quad\quad \textbf{C.} Hình thang vuông\quad\quad \textbf{D.} Hình tứ giác thường\par\vspace{2pt}
\textbf{\color{primary} Câu 9.} Kiểm tra nhanh kiến thức nền tảng: Hình bình hành có hai đường chéo vuông góc với nhau là hình gì?\par
\hspace*{1.5em}\textbf{A.} Hình chữ nhật\quad\quad \textbf{B.} Hình thang cân\quad\quad \textbf{C.} Hình thoi\quad\quad \textbf{D.} Hình thang vuông\par\vspace{2pt}
\textbf{\color{primary} Câu 10.} Cho dữ kiện nhận biết sau đây: Cho tứ giác ABCD có góc A = 79°, góc B = 95°, góc C = 96°. Số đo góc D bằng. Phương án chính xác là:\par
\hspace*{1.5em}\textbf{A.} 100°\quad\quad \textbf{B.} 80°\quad\quad \textbf{C.} 110°\quad\quad \textbf{D.} 90°\par\vspace{2pt}
\textbf{\color{primary} Câu 11.} Hãy chỉ ra kết quả đúng trong bốn lựa chọn A, B, C, D đối với yêu cầu: Hình bình hành có hai đường chéo bằng nhau là hình gì:\par
\hspace*{1.5em}\textbf{A.} Hình chữ nhật\quad\quad \textbf{B.} Hình thoi\quad\quad \textbf{C.} Hình thang vuông\quad\quad \textbf{D.} Hình tứ giác thường\par\vspace{2pt}
\textbf{\color{primary} Câu 12.} Câu hỏi ôn tập lý thuyết và kỹ năng cơ bản (Chương III): Hình bình hành có hai đường chéo vuông góc với nhau là hình gì:\par
\hspace*{1.5em}\textbf{A.} Hình thang vuông\quad\quad \textbf{B.} Hình chữ nhật\quad\quad \textbf{C.} Hình thoi\quad\quad \textbf{D.} Hình thang cân\par\vspace{2pt}
\textbf{\color{primary} Câu 13.} Theo nội dung sách giáo khoa Kết nối tri thức với cuộc sống (Chương III), cho tứ giác ABCD có góc A = 82°, góc B = 95°, góc C = 93°. Số đo góc D bằng:\par
\hspace*{1.5em}\textbf{A.} 90°\quad\quad \textbf{B.} 80°\quad\quad \textbf{C.} 110°\quad\quad \textbf{D.} 100°\par\vspace{2pt}
\textbf{\color{primary} Câu 14.} Tìm đáp số đúng của bài toán cơ bản sau: Hình bình hành có hai đường chéo bằng nhau là hình gì:\par
\hspace*{1.5em}\textbf{A.} Hình chữ nhật\quad\quad \textbf{B.} Hình thang vuông\quad\quad \textbf{C.} Hình thoi\quad\quad \textbf{D.} Hình tứ giác thường\par\vspace{2pt}
\textbf{\color{primary} Câu 15.} Xét đối tượng toán học trong Chương III: Hình bình hành có hai đường chéo vuông góc với nhau là hình gì. Giá trị hoặc tính chất tương ứng là:\par
\hspace*{1.5em}\textbf{A.} Hình thang cân\quad\quad \textbf{B.} Hình thoi\quad\quad \textbf{C.} Hình thang vuông\quad\quad \textbf{D.} Hình chữ nhật\par\vspace{2pt}
\textbf{\color{primary} Câu 16.} Hoàn thành phép xác định cơ bản sau đây: Cho tứ giác ABCD có góc A = 85°, góc B = 95°, góc C = 90°. Số đo góc D bằng:\par
\hspace*{1.5em}\textbf{A.} 90°\quad\quad \textbf{B.} 80°\quad\quad \textbf{C.} 110°\quad\quad \textbf{D.} 100°\par\vspace{2pt}
\textbf{\color{primary} Câu 17.} Trong các khẳng định hoặc giá trị dưới đây, đâu là đáp án đúng cho câu hỏi: Hình bình hành có hai đường chéo bằng nhau là hình gì?\par
\hspace*{1.5em}\textbf{A.} Hình chữ nhật\quad\quad \textbf{B.} Hình thang vuông\quad\quad \textbf{C.} Hình thoi\quad\quad \textbf{D.} Hình tứ giác thường\par\vspace{2pt}
\textbf{\color{primary} Câu 18.} Bài tập khởi động Chương III: Hình bình hành có hai đường chéo vuông góc với nhau là hình gì:\par
\hspace*{1.5em}\textbf{A.} Hình thoi\quad\quad \textbf{B.} Hình thang cân\quad\quad \textbf{C.} Hình chữ nhật\quad\quad \textbf{D.} Hình thang vuông\par\vspace{2pt}
\textbf{\color{primary} Câu 19.} Quan sát dữ kiện đề bài: Cho tứ giác ABCD có góc A = 88°, góc B = 95°, góc C = 87°. Số đo góc D bằng. Hãy chọn phương án đúng:\par
\hspace*{1.5em}\textbf{A.} 80°\quad\quad \textbf{B.} 90°\quad\quad \textbf{C.} 110°\quad\quad \textbf{D.} 100°\par\vspace{2pt}
\textbf{\color{primary} Câu 20.} Sử dụng công thức hoặc tính chất đầu bài của Chương III để giải quyết yêu cầu: Hình bình hành có hai đường chéo bằng nhau là hình gì:\par
\hspace*{1.5em}\textbf{A.} Hình tứ giác thường\quad\quad \textbf{B.} Hình thoi\quad\quad \textbf{C.} Hình chữ nhật\quad\quad \textbf{D.} Hình thang vuông\par\vspace{2pt}
\subsection*{\color{accent} Phần II. Mức độ Trung bình (Thông hiểu -- Câu 21 đến Câu 40)}
\textbf{\color{primary} Câu 21.} Cho tam giác ABC vuông tại A có cạnh huyền BC = 10 cm. Độ dài đường trung tuyến AM ứng với cạnh huyền BC bằng:\par
\hspace*{1.5em}\textbf{A.} 20 cm\quad\quad \textbf{B.} 10 cm\quad\quad \textbf{C.} 5 cm\quad\quad \textbf{D.} 3 cm\par\vspace{2pt}
\textbf{\color{primary} Câu 22.} Vận dụng quy tắc trong Chương III (Tứ giác) để giải bài toán: Cho hình bình hành ABCD có góc A - góc B = 24°. Số đo góc A của hình bình hành đó là:\par
\hspace*{1.5em}\textbf{A.} 78°\quad\quad \textbf{B.} 102°\quad\quad \textbf{C.} 112°\quad\quad \textbf{D.} 92°\par\vspace{2pt}
\textbf{\color{primary} Câu 23.} Biến đổi và tính toán theo kiến thức Chương III: Một hình thoi có độ dài hai đường chéo là 18 cm và 24 cm. Độ dài cạnh của hình thoi đó (tính theo định lí Pythagore) bằng. Kết quả thu được là:\par
\hspace*{1.5em}\textbf{A.} 42 cm\quad\quad \textbf{B.} 15 cm\quad\quad \textbf{C.} 30 cm\quad\quad \textbf{D.} 21 cm\par\vspace{2pt}
\textbf{\color{primary} Câu 24.} Xét bài toán mức độ thông hiểu sau: Cho tam giác ABC vuông tại A có cạnh huyền BC = 16 cm. Độ dài đường trung tuyến AM ứng với cạnh huyền BC bằng. Đáp án đúng là:\par
\hspace*{1.5em}\textbf{A.} 16 cm\quad\quad \textbf{B.} 6 cm\quad\quad \textbf{C.} 32 cm\quad\quad \textbf{D.} 8 cm\par\vspace{2pt}
\textbf{\color{primary} Câu 25.} Để hoàn thành phiếu học tập Chương III, hãy xác định: Cho hình bình hành ABCD có góc A - góc B = 36°. Số đo góc A của hình bình hành đó là:\par
\hspace*{1.5em}\textbf{A.} 98°\quad\quad \textbf{B.} 118°\quad\quad \textbf{C.} 72°\quad\quad \textbf{D.} 108°\par\vspace{2pt}
\textbf{\color{primary} Câu 26.} Phân tích dữ kiện bài toán: Một hình thoi có độ dài hai đường chéo là 18 cm và 24 cm. Độ dài cạnh của hình thoi đó (tính theo định lí Pythagore) bằng. Giá trị cần tìm bằng:\par
\hspace*{1.5em}\textbf{A.} 15 cm\quad\quad \textbf{B.} 30 cm\quad\quad \textbf{C.} 42 cm\quad\quad \textbf{D.} 21 cm\par\vspace{2pt}
\textbf{\color{primary} Câu 27.} Áp dụng tính chất đã học ở Chương III vào bài tập sau: Cho tam giác ABC vuông tại A có cạnh huyền BC = 22 cm. Độ dài đường trung tuyến AM ứng với cạnh huyền BC bằng:\par
\hspace*{1.5em}\textbf{A.} 22 cm\quad\quad \textbf{B.} 9 cm\quad\quad \textbf{C.} 11 cm\quad\quad \textbf{D.} 44 cm\par\vspace{2pt}
\textbf{\color{primary} Câu 28.} Thực hiện các bước suy luận hợp lý để giải quyết yêu cầu: Cho hình bình hành ABCD có góc A - góc B = 48°. Số đo góc A của hình bình hành đó là:\par
\hspace*{1.5em}\textbf{A.} 104°\quad\quad \textbf{B.} 124°\quad\quad \textbf{C.} 66°\quad\quad \textbf{D.} 114°\par\vspace{2pt}
\textbf{\color{primary} Câu 29.} Cho bài toán tự luyện (Tứ giác): Một hình thoi có độ dài hai đường chéo là 18 cm và 24 cm. Độ dài cạnh của hình thoi đó (tính theo định lí Pythagore) bằng. Phương án nào sau đây chính xác?\par
\hspace*{1.5em}\textbf{A.} 30 cm\quad\quad \textbf{B.} 21 cm\quad\quad \textbf{C.} 42 cm\quad\quad \textbf{D.} 15 cm\par\vspace{2pt}
\textbf{\color{primary} Câu 30.} Kết hợp các phép biến đổi trong Chương III, hãy tìm kết quả của bài toán: Cho tam giác ABC vuông tại A có cạnh huyền BC = 28 cm. Độ dài đường trung tuyến AM ứng với cạnh huyền BC bằng:\par
\hspace*{1.5em}\textbf{A.} 56 cm\quad\quad \textbf{B.} 28 cm\quad\quad \textbf{C.} 14 cm\quad\quad \textbf{D.} 12 cm\par\vspace{2pt}
\textbf{\color{primary} Câu 31.} Một bài kiểm tra định kỳ nội dung Chương III nêu câu hỏi: Cho hình bình hành ABCD có góc A - góc B = 60°. Số đo góc A của hình bình hành đó là. Hãy chọn đáp số đúng:\par
\hspace*{1.5em}\textbf{A.} 120°\quad\quad \textbf{B.} 130°\quad\quad \textbf{C.} 110°\quad\quad \textbf{D.} 60°\par\vspace{2pt}
\textbf{\color{primary} Câu 32.} Tìm lời giải đúng cho bài toán thông hiểu: Một hình thoi có độ dài hai đường chéo là 18 cm và 24 cm. Độ dài cạnh của hình thoi đó (tính theo định lí Pythagore) bằng:\par
\hspace*{1.5em}\textbf{A.} 21 cm\quad\quad \textbf{B.} 15 cm\quad\quad \textbf{C.} 30 cm\quad\quad \textbf{D.} 42 cm\par\vspace{2pt}
\textbf{\color{primary} Câu 33.} Sử dụng mối liên hệ giữa các đại lượng trong Chương III: Cho tam giác ABC vuông tại A có cạnh huyền BC = 34 cm. Độ dài đường trung tuyến AM ứng với cạnh huyền BC bằng:\par
\hspace*{1.5em}\textbf{A.} 34 cm\quad\quad \textbf{B.} 15 cm\quad\quad \textbf{C.} 68 cm\quad\quad \textbf{D.} 17 cm\par\vspace{2pt}
\textbf{\color{primary} Câu 34.} Hãy tính toán và chọn phương án đúng đối với bài tập: Cho hình bình hành ABCD có góc A - góc B = 72°. Số đo góc A của hình bình hành đó là:\par
\hspace*{1.5em}\textbf{A.} 54°\quad\quad \textbf{B.} 116°\quad\quad \textbf{C.} 136°\quad\quad \textbf{D.} 126°\par\vspace{2pt}
\textbf{\color{primary} Câu 35.} Dựa trên các tính chất trọng tâm của Chương III, xét bài toán: Một hình thoi có độ dài hai đường chéo là 18 cm và 24 cm. Độ dài cạnh của hình thoi đó (tính theo định lí Pythagore) bằng:\par
\hspace*{1.5em}\textbf{A.} 21 cm\quad\quad \textbf{B.} 15 cm\quad\quad \textbf{C.} 42 cm\quad\quad \textbf{D.} 30 cm\par\vspace{2pt}
\textbf{\color{primary} Câu 36.} Thực hành giải toán (Tứ giác): Cho tam giác ABC vuông tại A có cạnh huyền BC = 40 cm. Độ dài đường trung tuyến AM ứng với cạnh huyền BC bằng. Đáp án nào dưới đây thỏa mãn?\par
\hspace*{1.5em}\textbf{A.} 18 cm\quad\quad \textbf{B.} 80 cm\quad\quad \textbf{C.} 20 cm\quad\quad \textbf{D.} 40 cm\par\vspace{2pt}
\textbf{\color{primary} Câu 37.} Kiểm tra kỹ năng biến đổi và suy luận qua bài toán: Cho hình bình hành ABCD có góc A - góc B = 84°. Số đo góc A của hình bình hành đó là:\par
\hspace*{1.5em}\textbf{A.} 122°\quad\quad \textbf{B.} 142°\quad\quad \textbf{C.} 132°\quad\quad \textbf{D.} 48°\par\vspace{2pt}
\textbf{\color{primary} Câu 38.} Cho giả thiết thuộc phạm vi Chương III: Một hình thoi có độ dài hai đường chéo là 18 cm và 24 cm. Độ dài cạnh của hình thoi đó (tính theo định lí Pythagore) bằng. Kết luận đúng là:\par
\hspace*{1.5em}\textbf{A.} 15 cm\quad\quad \textbf{B.} 42 cm\quad\quad \textbf{C.} 30 cm\quad\quad \textbf{D.} 21 cm\par\vspace{2pt}
\textbf{\color{primary} Câu 39.} Rút gọn hoặc tính giá trị theo yêu cầu sau của Chương III: Cho tam giác ABC vuông tại A có cạnh huyền BC = 46 cm. Độ dài đường trung tuyến AM ứng với cạnh huyền BC bằng:\par
\hspace*{1.5em}\textbf{A.} 23 cm\quad\quad \textbf{B.} 21 cm\quad\quad \textbf{C.} 46 cm\quad\quad \textbf{D.} 92 cm\par\vspace{2pt}
\textbf{\color{primary} Câu 40.} Xác định phương án đúng sau khi thực hiện phép tính / suy luận: Cho hình bình hành ABCD có góc A - góc B = 96°. Số đo góc A của hình bình hành đó là:\par
\hspace*{1.5em}\textbf{A.} 148°\quad\quad \textbf{B.} 128°\quad\quad \textbf{C.} 138°\quad\quad \textbf{D.} 42°\par\vspace{2pt}
\subsection*{\color{accent} Phần III. Mức độ Khó (Vận dụng \& Vận dụng cao -- Câu 41 đến Câu 50)}
\textbf{\color{primary} Câu 41.} Cho tứ giác ABCD có hai đường chéo AC và BD vuông góc với nhau, biết AC = 6 cm, BD = 8 cm. Gọi M, N, P, Q lần lượt là trung điểm của AB, BC, CD, DA. Tứ giác MNPQ là hình chữ nhật có diện tích bằng:\par
\hspace*{1.5em}\textbf{A.} 48 cm$^2$\quad\quad \textbf{B.} 12 cm$^2$\quad\quad \textbf{C.} 14 cm$^2$\quad\quad \textbf{D.} 24 cm$^2$\par\vspace{2pt}
\textbf{\color{primary} Câu 42.} Bài toán vận dụng (Chương III – Tứ giác): Cho hình vuông ABCD có độ dài đường chéo bằng 10 cm. Diện tích của hình vuông ABCD bằng. Kết quả đúng là:\par
\hspace*{1.5em}\textbf{A.} 40 cm$^2$\quad\quad \textbf{B.} 50 cm$^2$\quad\quad \textbf{C.} 25 cm$^2$\quad\quad \textbf{D.} 100 cm$^2$\par\vspace{2pt}
\textbf{\color{primary} Câu 43.} Vận dụng kiến thức nâng cao trong phạm vi Chương III để giải bài toán: Cho tứ giác ABCD có hai đường chéo AC và BD vuông góc với nhau, biết AC = 10 cm, BD = 12 cm. Gọi M, N, P, Q lần lượt là trung điểm của AB, BC, CD, DA. Tứ giác MNPQ là hình chữ nhật có diện tích bằng:\par
\hspace*{1.5em}\textbf{A.} 22 cm$^2$\quad\quad \textbf{B.} 30 cm$^2$\quad\quad \textbf{C.} 120 cm$^2$\quad\quad \textbf{D.} 60 cm$^2$\par\vspace{2pt}
\textbf{\color{primary} Câu 44.} Phân tích và giải bài toán tổng hợp sau của Chương III: Cho hình vuông ABCD có độ dài đường chéo bằng 14 cm. Diện tích của hình vuông ABCD bằng:\par
\hspace*{1.5em}\textbf{A.} 56 cm$^2$\quad\quad \textbf{B.} 49 cm$^2$\quad\quad \textbf{C.} 196 cm$^2$\quad\quad \textbf{D.} 98 cm$^2$\par\vspace{2pt}
\textbf{\color{primary} Câu 45.} Thử thách tư duy trong phạm vi kiến thức Chương III: Cho tứ giác ABCD có hai đường chéo AC và BD vuông góc với nhau, biết AC = 14 cm, BD = 16 cm. Gọi M, N, P, Q lần lượt là trung điểm của AB, BC, CD, DA. Tứ giác MNPQ là hình chữ nhật có diện tích bằng. Đáp số chính xác là:\par
\hspace*{1.5em}\textbf{A.} 112 cm$^2$\quad\quad \textbf{B.} 224 cm$^2$\quad\quad \textbf{C.} 30 cm$^2$\quad\quad \textbf{D.} 56 cm$^2$\par\vspace{2pt}
\textbf{\color{primary} Câu 46.} Sử dụng phương pháp giải toán của Chương III để tìm đáp án cho bài toán: Cho hình vuông ABCD có độ dài đường chéo bằng 18 cm. Diện tích của hình vuông ABCD bằng:\par
\hspace*{1.5em}\textbf{A.} 72 cm$^2$\quad\quad \textbf{B.} 162 cm$^2$\quad\quad \textbf{C.} 81 cm$^2$\quad\quad \textbf{D.} 324 cm$^2$\par\vspace{2pt}
\textbf{\color{primary} Câu 47.} Xét bài toán vận dụng thực tiễn / nâng cao thuộc Chương III: Cho tứ giác ABCD có hai đường chéo AC và BD vuông góc với nhau, biết AC = 18 cm, BD = 20 cm. Gọi M, N, P, Q lần lượt là trung điểm của AB, BC, CD, DA. Tứ giác MNPQ là hình chữ nhật có diện tích bằng:\par
\hspace*{1.5em}\textbf{A.} 360 cm$^2$\quad\quad \textbf{B.} 38 cm$^2$\quad\quad \textbf{C.} 90 cm$^2$\quad\quad \textbf{D.} 180 cm$^2$\par\vspace{2pt}
\textbf{\color{primary} Câu 48.} Tìm kết quả cuối cùng của bài toán phân loại học sinh khá giỏi (Chương III): Cho hình vuông ABCD có độ dài đường chéo bằng 22 cm. Diện tích của hình vuông ABCD bằng:\par
\hspace*{1.5em}\textbf{A.} 484 cm$^2$\quad\quad \textbf{B.} 242 cm$^2$\quad\quad \textbf{C.} 121 cm$^2$\quad\quad \textbf{D.} 88 cm$^2$\par\vspace{2pt}
\textbf{\color{primary} Câu 49.} Kết hợp các tính chất đã học trong Chương III, hãy giải quyết bài toán: Cho tứ giác ABCD có hai đường chéo AC và BD vuông góc với nhau, biết AC = 22 cm, BD = 24 cm. Gọi M, N, P, Q lần lượt là trung điểm của AB, BC, CD, DA. Tứ giác MNPQ là hình chữ nhật có diện tích bằng:\par
\hspace*{1.5em}\textbf{A.} 528 cm$^2$\quad\quad \textbf{B.} 264 cm$^2$\quad\quad \textbf{C.} 46 cm$^2$\quad\quad \textbf{D.} 132 cm$^2$\par\vspace{2pt}
\textbf{\color{primary} Câu 50.} Câu hỏi vận dụng cao cuối Chương III: Cho hình vuông ABCD có độ dài đường chéo bằng 26 cm. Diện tích của hình vuông ABCD bằng. Phương án đúng là:\par
\hspace*{1.5em}\textbf{A.} 169 cm$^2$\quad\quad \textbf{B.} 104 cm$^2$\quad\quad \textbf{C.} 338 cm$^2$\quad\quad \textbf{D.} 676 cm$^2$\par\vspace{2pt}
\subsection*{\color{primary} Bảng đáp án Chương III: Tứ giác}
\begin{center}
\renewcommand{\arraystretch}{1.25}
\begin{tabularx}{\textwidth}{*{10}{>{\centering\arraybackslash}X}}
\toprule
\textbf{\color{primary} Cột 1} & \textbf{\color{primary} Cột 2} & \textbf{\color{primary} Cột 3} & \textbf{\color{primary} Cột 4} & \textbf{\color{primary} Cột 5} & \textbf{\color{primary} Cột 6} & \textbf{\color{primary} Cột 7} & \textbf{\color{primary} Cột 8} & \textbf{\color{primary} Cột 9} & \textbf{\color{primary} Cột 10} \\
\midrule
\textbf{1.A} & \textbf{2.A} & \textbf{3.B} & \textbf{4.D} & \textbf{5.A} & \textbf{6.A} & \textbf{7.D} & \textbf{8.A} & \textbf{9.C} & \textbf{10.D} \\
\rowcolor{rowtint} \textbf{11.A} & \textbf{12.C} & \textbf{13.A} & \textbf{14.A} & \textbf{15.B} & \textbf{16.A} & \textbf{17.A} & \textbf{18.A} & \textbf{19.B} & \textbf{20.C} \\
\textbf{21.C} & \textbf{22.B} & \textbf{23.B} & \textbf{24.D} & \textbf{25.D} & \textbf{26.A} & \textbf{27.C} & \textbf{28.D} & \textbf{29.D} & \textbf{30.C} \\
\rowcolor{rowtint} \textbf{31.A} & \textbf{32.B} & \textbf{33.D} & \textbf{34.D} & \textbf{35.B} & \textbf{36.C} & \textbf{37.C} & \textbf{38.A} & \textbf{39.A} & \textbf{40.C} \\
\textbf{41.B} & \textbf{42.B} & \textbf{43.B} & \textbf{44.D} & \textbf{45.D} & \textbf{46.B} & \textbf{47.C} & \textbf{48.B} & \textbf{49.D} & \textbf{50.C} \\
\bottomrule
\end{tabularx}
\end{center}\vspace{10pt}
\section*{\color{primary} CHƯƠNG IV: ĐỊNH LÍ THALÈS}
\textit{Trọng tâm kiến thức SGK: Định lí Thalès thuận và đảo trong tam giác, đường trung bình của tam giác, tính chất đường phân giác của tam giác.}\par\vspace{4pt}
\subsection*{\color{accent} Phần I. Mức độ Dễ (Nhận biết -- Câu 1 đến Câu 20)}
\textbf{\color{primary} Câu 1.} Cho tam giác ABC có cạnh BC = 8 cm. Gọi M, N lần lượt là trung điểm của AB và AC. Độ dài đường trung bình MN bằng:\par
\hspace*{1.5em}\textbf{A.} 16 cm\quad\quad \textbf{B.} 4 cm\quad\quad \textbf{C.} 8 cm\quad\quad \textbf{D.} 2 cm\par\vspace{2pt}
\textbf{\color{primary} Câu 2.} Thực hiện yêu cầu nhận biết sau: Cho đoạn thẳng AB = 6 cm và CD = 10 cm. Tỉ số của hai đoạn thẳng AB và CD là.\par
\hspace*{1.5em}\textbf{A.} 3/5\quad\quad \textbf{B.} 3/8\quad\quad \textbf{C.} 2/5\quad\quad \textbf{D.} 5/3\par\vspace{2pt}
\textbf{\color{primary} Câu 3.} Trong bài học thuộc Chương IV, xét nội dung: Đường trung bình của một tam giác thì có tính chất nào sau đây. Đáp án đúng là:\par
\begin{itemize}\setlength{\itemsep}{1pt}\setlength{\parskip}{0pt}
  \item[] \textbf{A.} Song song và bằng cạnh thứ ba
  \item[] \textbf{B.} Song song với cạnh thứ ba và bằng nửa cạnh ấy
  \item[] \textbf{C.} Đi qua trọng tâm của tam giác
  \item[] \textbf{D.} Vuông góc với cạnh thứ ba và bằng nửa cạnh ấy
\end{itemize}\vspace{2pt}
\textbf{\color{primary} Câu 4.} Dựa vào kiến thức chuẩn SGK của Chương IV (Định lí Thalès), cho tam giác ABC có cạnh BC = 14 cm. Gọi M, N lần lượt là trung điểm của AB và AC. Độ dài đường trung bình MN bằng:\par
\hspace*{1.5em}\textbf{A.} 28 cm\quad\quad \textbf{B.} 5 cm\quad\quad \textbf{C.} 14 cm\quad\quad \textbf{D.} 7 cm\par\vspace{2pt}
\textbf{\color{primary} Câu 5.} Chọn kết quả đúng nhất khi xét: Cho đoạn thẳng AB = 15 cm và CD = 25 cm. Tỉ số của hai đoạn thẳng AB và CD là:\par
\hspace*{1.5em}\textbf{A.} 3/8\quad\quad \textbf{B.} 2/5\quad\quad \textbf{C.} 3/5\quad\quad \textbf{D.} 5/3\par\vspace{2pt}
\textbf{\color{primary} Câu 6.} Khi làm bài tập cơ bản về định lí thalès, một học sinh gặp câu hỏi: Đường trung bình của một tam giác thì có tính chất nào sau đây. Kết quả nào dưới đây đúng?\par
\begin{itemize}\setlength{\itemsep}{1pt}\setlength{\parskip}{0pt}
  \item[] \textbf{A.} Song song và bằng cạnh thứ ba
  \item[] \textbf{B.} Vuông góc với cạnh thứ ba và bằng nửa cạnh ấy
  \item[] \textbf{C.} Đi qua trọng tâm của tam giác
  \item[] \textbf{D.} Song song với cạnh thứ ba và bằng nửa cạnh ấy
\end{itemize}\vspace{2pt}
\textbf{\color{primary} Câu 7.} Xác định giá trị hoặc khẳng định đúng cho bài toán sau: Cho tam giác ABC có cạnh BC = 20 cm. Gọi M, N lần lượt là trung điểm của AB và AC. Độ dài đường trung bình MN bằng:\par
\hspace*{1.5em}\textbf{A.} 20 cm\quad\quad \textbf{B.} 8 cm\quad\quad \textbf{C.} 10 cm\quad\quad \textbf{D.} 40 cm\par\vspace{2pt}
\textbf{\color{primary} Câu 8.} Áp dụng trực tiếp định nghĩa và quy tắc trong Chương IV: Cho đoạn thẳng AB = 24 cm và CD = 40 cm. Tỉ số của hai đoạn thẳng AB và CD là:\par
\hspace*{1.5em}\textbf{A.} 3/8\quad\quad \textbf{B.} 3/5\quad\quad \textbf{C.} 5/3\quad\quad \textbf{D.} 2/5\par\vspace{2pt}
\textbf{\color{primary} Câu 9.} Kiểm tra nhanh kiến thức nền tảng: Đường trung bình của một tam giác thì có tính chất nào sau đây?\par
\begin{itemize}\setlength{\itemsep}{1pt}\setlength{\parskip}{0pt}
  \item[] \textbf{A.} Vuông góc với cạnh thứ ba và bằng nửa cạnh ấy
  \item[] \textbf{B.} Đi qua trọng tâm của tam giác
  \item[] \textbf{C.} Song song với cạnh thứ ba và bằng nửa cạnh ấy
  \item[] \textbf{D.} Song song và bằng cạnh thứ ba
\end{itemize}\vspace{2pt}
\textbf{\color{primary} Câu 10.} Cho dữ kiện nhận biết sau đây: Cho tam giác ABC có cạnh BC = 26 cm. Gọi M, N lần lượt là trung điểm của AB và AC. Độ dài đường trung bình MN bằng. Phương án chính xác là:\par
\hspace*{1.5em}\textbf{A.} 13 cm\quad\quad \textbf{B.} 26 cm\quad\quad \textbf{C.} 52 cm\quad\quad \textbf{D.} 11 cm\par\vspace{2pt}
\textbf{\color{primary} Câu 11.} Hãy chỉ ra kết quả đúng trong bốn lựa chọn A, B, C, D đối với yêu cầu: Cho đoạn thẳng AB = 33 cm và CD = 55 cm. Tỉ số của hai đoạn thẳng AB và CD là:\par
\hspace*{1.5em}\textbf{A.} 5/3\quad\quad \textbf{B.} 3/5\quad\quad \textbf{C.} 2/5\quad\quad \textbf{D.} 3/8\par\vspace{2pt}
\textbf{\color{primary} Câu 12.} Câu hỏi ôn tập lý thuyết và kỹ năng cơ bản (Chương IV): Đường trung bình của một tam giác thì có tính chất nào sau đây:\par
\begin{itemize}\setlength{\itemsep}{1pt}\setlength{\parskip}{0pt}
  \item[] \textbf{A.} Vuông góc với cạnh thứ ba và bằng nửa cạnh ấy
  \item[] \textbf{B.} Đi qua trọng tâm của tam giác
  \item[] \textbf{C.} Song song và bằng cạnh thứ ba
  \item[] \textbf{D.} Song song với cạnh thứ ba và bằng nửa cạnh ấy
\end{itemize}\vspace{2pt}
\textbf{\color{primary} Câu 13.} Theo nội dung sách giáo khoa Kết nối tri thức với cuộc sống (Chương IV), cho tam giác ABC có cạnh BC = 32 cm. Gọi M, N lần lượt là trung điểm của AB và AC. Độ dài đường trung bình MN bằng:\par
\hspace*{1.5em}\textbf{A.} 32 cm\quad\quad \textbf{B.} 16 cm\quad\quad \textbf{C.} 64 cm\quad\quad \textbf{D.} 14 cm\par\vspace{2pt}
\textbf{\color{primary} Câu 14.} Tìm đáp số đúng của bài toán cơ bản sau: Cho đoạn thẳng AB = 42 cm và CD = 70 cm. Tỉ số của hai đoạn thẳng AB và CD là:\par
\hspace*{1.5em}\textbf{A.} 3/8\quad\quad \textbf{B.} 3/5\quad\quad \textbf{C.} 2/5\quad\quad \textbf{D.} 5/3\par\vspace{2pt}
\textbf{\color{primary} Câu 15.} Xét đối tượng toán học trong Chương IV: Đường trung bình của một tam giác thì có tính chất nào sau đây. Giá trị hoặc tính chất tương ứng là:\par
\begin{itemize}\setlength{\itemsep}{1pt}\setlength{\parskip}{0pt}
  \item[] \textbf{A.} Đi qua trọng tâm của tam giác
  \item[] \textbf{B.} Vuông góc với cạnh thứ ba và bằng nửa cạnh ấy
  \item[] \textbf{C.} Song song với cạnh thứ ba và bằng nửa cạnh ấy
  \item[] \textbf{D.} Song song và bằng cạnh thứ ba
\end{itemize}\vspace{2pt}
\textbf{\color{primary} Câu 16.} Hoàn thành phép xác định cơ bản sau đây: Cho tam giác ABC có cạnh BC = 38 cm. Gọi M, N lần lượt là trung điểm của AB và AC. Độ dài đường trung bình MN bằng:\par
\hspace*{1.5em}\textbf{A.} 38 cm\quad\quad \textbf{B.} 19 cm\quad\quad \textbf{C.} 76 cm\quad\quad \textbf{D.} 17 cm\par\vspace{2pt}
\textbf{\color{primary} Câu 17.} Trong các khẳng định hoặc giá trị dưới đây, đâu là đáp án đúng cho câu hỏi: Cho đoạn thẳng AB = 51 cm và CD = 85 cm. Tỉ số của hai đoạn thẳng AB và CD là?\par
\hspace*{1.5em}\textbf{A.} 2/5\quad\quad \textbf{B.} 3/5\quad\quad \textbf{C.} 5/3\quad\quad \textbf{D.} 3/8\par\vspace{2pt}
\textbf{\color{primary} Câu 18.} Bài tập khởi động Chương IV: Đường trung bình của một tam giác thì có tính chất nào sau đây:\par
\begin{itemize}\setlength{\itemsep}{1pt}\setlength{\parskip}{0pt}
  \item[] \textbf{A.} Đi qua trọng tâm của tam giác
  \item[] \textbf{B.} Song song với cạnh thứ ba và bằng nửa cạnh ấy
  \item[] \textbf{C.} Vuông góc với cạnh thứ ba và bằng nửa cạnh ấy
  \item[] \textbf{D.} Song song và bằng cạnh thứ ba
\end{itemize}\vspace{2pt}
\textbf{\color{primary} Câu 19.} Quan sát dữ kiện đề bài: Cho tam giác ABC có cạnh BC = 44 cm. Gọi M, N lần lượt là trung điểm của AB và AC. Độ dài đường trung bình MN bằng. Hãy chọn phương án đúng:\par
\hspace*{1.5em}\textbf{A.} 44 cm\quad\quad \textbf{B.} 22 cm\quad\quad \textbf{C.} 88 cm\quad\quad \textbf{D.} 20 cm\par\vspace{2pt}
\textbf{\color{primary} Câu 20.} Sử dụng công thức hoặc tính chất đầu bài của Chương IV để giải quyết yêu cầu: Cho đoạn thẳng AB = 60 cm và CD = 100 cm. Tỉ số của hai đoạn thẳng AB và CD là:\par
\hspace*{1.5em}\textbf{A.} 3/8\quad\quad \textbf{B.} 2/5\quad\quad \textbf{C.} 3/5\quad\quad \textbf{D.} 5/3\par\vspace{2pt}
\subsection*{\color{accent} Phần II. Mức độ Trung bình (Thông hiểu -- Câu 21 đến Câu 40)}
\textbf{\color{primary} Câu 21.} Cho tam giác ABC, đường thẳng d song song với BC cắt AB, AC lần lượt tại D và E. Biết AD = 2 cm, DB = 4 cm và AE = 3 cm. Độ dài đoạn EC bằng:\par
\hspace*{1.5em}\textbf{A.} 3 cm\quad\quad \textbf{B.} 6 cm\quad\quad \textbf{C.} 4 cm\quad\quad \textbf{D.} 9 cm\par\vspace{2pt}
\textbf{\color{primary} Câu 22.} Vận dụng quy tắc trong Chương IV (Định lí Thalès) để giải bài toán: Cho tam giác ABC có tia phân giác trong AD (D $\in$ BC). Biết AB = 6 cm, AC = 9 cm và BD = 6 cm. Độ dài đoạn DC bằng:\par
\hspace*{1.5em}\textbf{A.} 7 cm\quad\quad \textbf{B.} 12 cm\quad\quad \textbf{C.} 9 cm\quad\quad \textbf{D.} 6 cm\par\vspace{2pt}
\textbf{\color{primary} Câu 23.} Biến đổi và tính toán theo kiến thức Chương IV: Cho tam giác ABC có chu vi bằng 32 cm. Gọi D, E, F lần lượt là trung điểm của các cạnh AB, BC, CA. Chu vi của tam giác DEF bằng. Kết quả thu được là:\par
\hspace*{1.5em}\textbf{A.} 64 cm\quad\quad \textbf{B.} 8 cm\quad\quad \textbf{C.} 32 cm\quad\quad \textbf{D.} 16 cm\par\vspace{2pt}
\textbf{\color{primary} Câu 24.} Xét bài toán mức độ thông hiểu sau: Cho tam giác ABC, đường thẳng d song song với BC cắt AB, AC lần lượt tại D và E. Biết AD = 5 cm, DB = 10 cm và AE = 6 cm. Độ dài đoạn EC bằng. Đáp án đúng là:\par
\hspace*{1.5em}\textbf{A.} 12 cm\quad\quad \textbf{B.} 10 cm\quad\quad \textbf{C.} 6 cm\quad\quad \textbf{D.} 18 cm\par\vspace{2pt}
\textbf{\color{primary} Câu 25.} Để hoàn thành phiếu học tập Chương IV, hãy xác định: Cho tam giác ABC có tia phân giác trong AD (D $\in$ BC). Biết AB = 12 cm, AC = 18 cm và BD = 6 cm. Độ dài đoạn DC bằng:\par
\hspace*{1.5em}\textbf{A.} 9 cm\quad\quad \textbf{B.} 7 cm\quad\quad \textbf{C.} 6 cm\quad\quad \textbf{D.} 12 cm\par\vspace{2pt}
\textbf{\color{primary} Câu 26.} Phân tích dữ kiện bài toán: Cho tam giác ABC có chu vi bằng 44 cm. Gọi D, E, F lần lượt là trung điểm của các cạnh AB, BC, CA. Chu vi của tam giác DEF bằng. Giá trị cần tìm bằng:\par
\hspace*{1.5em}\textbf{A.} 44 cm\quad\quad \textbf{B.} 11 cm\quad\quad \textbf{C.} 88 cm\quad\quad \textbf{D.} 22 cm\par\vspace{2pt}
\textbf{\color{primary} Câu 27.} Áp dụng tính chất đã học ở Chương IV vào bài tập sau: Cho tam giác ABC, đường thẳng d song song với BC cắt AB, AC lần lượt tại D và E. Biết AD = 8 cm, DB = 16 cm và AE = 9 cm. Độ dài đoạn EC bằng:\par
\hspace*{1.5em}\textbf{A.} 9 cm\quad\quad \textbf{B.} 27 cm\quad\quad \textbf{C.} 18 cm\quad\quad \textbf{D.} 16 cm\par\vspace{2pt}
\textbf{\color{primary} Câu 28.} Thực hiện các bước suy luận hợp lý để giải quyết yêu cầu: Cho tam giác ABC có tia phân giác trong AD (D $\in$ BC). Biết AB = 18 cm, AC = 27 cm và BD = 6 cm. Độ dài đoạn DC bằng:\par
\hspace*{1.5em}\textbf{A.} 12 cm\quad\quad \textbf{B.} 6 cm\quad\quad \textbf{C.} 9 cm\quad\quad \textbf{D.} 7 cm\par\vspace{2pt}
\textbf{\color{primary} Câu 29.} Cho bài toán tự luyện (Định lí Thalès): Cho tam giác ABC có chu vi bằng 56 cm. Gọi D, E, F lần lượt là trung điểm của các cạnh AB, BC, CA. Chu vi của tam giác DEF bằng. Phương án nào sau đây chính xác?\par
\hspace*{1.5em}\textbf{A.} 28 cm\quad\quad \textbf{B.} 56 cm\quad\quad \textbf{C.} 112 cm\quad\quad \textbf{D.} 14 cm\par\vspace{2pt}
\textbf{\color{primary} Câu 30.} Kết hợp các phép biến đổi trong Chương IV, hãy tìm kết quả của bài toán: Cho tam giác ABC, đường thẳng d song song với BC cắt AB, AC lần lượt tại D và E. Biết AD = 11 cm, DB = 22 cm và AE = 12 cm. Độ dài đoạn EC bằng:\par
\hspace*{1.5em}\textbf{A.} 36 cm\quad\quad \textbf{B.} 24 cm\quad\quad \textbf{C.} 12 cm\quad\quad \textbf{D.} 22 cm\par\vspace{2pt}
\textbf{\color{primary} Câu 31.} Một bài kiểm tra định kỳ nội dung Chương IV nêu câu hỏi: Cho tam giác ABC có tia phân giác trong AD (D $\in$ BC). Biết AB = 24 cm, AC = 36 cm và BD = 6 cm. Độ dài đoạn DC bằng. Hãy chọn đáp số đúng:\par
\hspace*{1.5em}\textbf{A.} 7 cm\quad\quad \textbf{B.} 12 cm\quad\quad \textbf{C.} 9 cm\quad\quad \textbf{D.} 6 cm\par\vspace{2pt}
\textbf{\color{primary} Câu 32.} Tìm lời giải đúng cho bài toán thông hiểu: Cho tam giác ABC có chu vi bằng 68 cm. Gọi D, E, F lần lượt là trung điểm của các cạnh AB, BC, CA. Chu vi của tam giác DEF bằng:\par
\hspace*{1.5em}\textbf{A.} 68 cm\quad\quad \textbf{B.} 34 cm\quad\quad \textbf{C.} 136 cm\quad\quad \textbf{D.} 17 cm\par\vspace{2pt}
\textbf{\color{primary} Câu 33.} Sử dụng mối liên hệ giữa các đại lượng trong Chương IV: Cho tam giác ABC, đường thẳng d song song với BC cắt AB, AC lần lượt tại D và E. Biết AD = 14 cm, DB = 28 cm và AE = 15 cm. Độ dài đoạn EC bằng:\par
\hspace*{1.5em}\textbf{A.} 45 cm\quad\quad \textbf{B.} 28 cm\quad\quad \textbf{C.} 15 cm\quad\quad \textbf{D.} 30 cm\par\vspace{2pt}
\textbf{\color{primary} Câu 34.} Hãy tính toán và chọn phương án đúng đối với bài tập: Cho tam giác ABC có tia phân giác trong AD (D $\in$ BC). Biết AB = 30 cm, AC = 45 cm và BD = 6 cm. Độ dài đoạn DC bằng:\par
\hspace*{1.5em}\textbf{A.} 9 cm\quad\quad \textbf{B.} 7 cm\quad\quad \textbf{C.} 6 cm\quad\quad \textbf{D.} 12 cm\par\vspace{2pt}
\textbf{\color{primary} Câu 35.} Dựa trên các tính chất trọng tâm của Chương IV, xét bài toán: Cho tam giác ABC có chu vi bằng 80 cm. Gọi D, E, F lần lượt là trung điểm của các cạnh AB, BC, CA. Chu vi của tam giác DEF bằng:\par
\hspace*{1.5em}\textbf{A.} 80 cm\quad\quad \textbf{B.} 160 cm\quad\quad \textbf{C.} 20 cm\quad\quad \textbf{D.} 40 cm\par\vspace{2pt}
\textbf{\color{primary} Câu 36.} Thực hành giải toán (Định lí Thalès): Cho tam giác ABC, đường thẳng d song song với BC cắt AB, AC lần lượt tại D và E. Biết AD = 17 cm, DB = 34 cm và AE = 18 cm. Độ dài đoạn EC bằng. Đáp án nào dưới đây thỏa mãn?\par
\hspace*{1.5em}\textbf{A.} 34 cm\quad\quad \textbf{B.} 36 cm\quad\quad \textbf{C.} 18 cm\quad\quad \textbf{D.} 54 cm\par\vspace{2pt}
\textbf{\color{primary} Câu 37.} Kiểm tra kỹ năng biến đổi và suy luận qua bài toán: Cho tam giác ABC có tia phân giác trong AD (D $\in$ BC). Biết AB = 36 cm, AC = 54 cm và BD = 6 cm. Độ dài đoạn DC bằng:\par
\hspace*{1.5em}\textbf{A.} 9 cm\quad\quad \textbf{B.} 7 cm\quad\quad \textbf{C.} 12 cm\quad\quad \textbf{D.} 6 cm\par\vspace{2pt}
\textbf{\color{primary} Câu 38.} Cho giả thiết thuộc phạm vi Chương IV: Cho tam giác ABC có chu vi bằng 92 cm. Gọi D, E, F lần lượt là trung điểm của các cạnh AB, BC, CA. Chu vi của tam giác DEF bằng. Kết luận đúng là:\par
\hspace*{1.5em}\textbf{A.} 23 cm\quad\quad \textbf{B.} 46 cm\quad\quad \textbf{C.} 184 cm\quad\quad \textbf{D.} 92 cm\par\vspace{2pt}
\textbf{\color{primary} Câu 39.} Rút gọn hoặc tính giá trị theo yêu cầu sau của Chương IV: Cho tam giác ABC, đường thẳng d song song với BC cắt AB, AC lần lượt tại D và E. Biết AD = 20 cm, DB = 40 cm và AE = 21 cm. Độ dài đoạn EC bằng:\par
\hspace*{1.5em}\textbf{A.} 42 cm\quad\quad \textbf{B.} 21 cm\quad\quad \textbf{C.} 40 cm\quad\quad \textbf{D.} 63 cm\par\vspace{2pt}
\textbf{\color{primary} Câu 40.} Xác định phương án đúng sau khi thực hiện phép tính / suy luận: Cho tam giác ABC có tia phân giác trong AD (D $\in$ BC). Biết AB = 42 cm, AC = 63 cm và BD = 6 cm. Độ dài đoạn DC bằng:\par
\hspace*{1.5em}\textbf{A.} 6 cm\quad\quad \textbf{B.} 12 cm\quad\quad \textbf{C.} 7 cm\quad\quad \textbf{D.} 9 cm\par\vspace{2pt}
\subsection*{\color{accent} Phần III. Mức độ Khó (Vận dụng \& Vận dụng cao -- Câu 41 đến Câu 50)}
\textbf{\color{primary} Câu 41.} Cho tam giác ABC có AB = 6 cm, AC = 10 cm, BC = 16 cm. Đường phân giác AD chia cạnh BC thành hai đoạn BD và CD. Độ dài đoạn BD bằng:\par
\hspace*{1.5em}\textbf{A.} 6 cm\quad\quad \textbf{B.} 4 cm\quad\quad \textbf{C.} 8 cm\quad\quad \textbf{D.} 10 cm\par\vspace{2pt}
\textbf{\color{primary} Câu 42.} Bài toán vận dụng (Chương IV – Định lí Thalès): Cho tam giác ABC, điểm M thuộc cạnh AB sao cho AM = 2$\cdot$MB. Qua M kẻ đường thẳng song song với BC cắt AC tại N. Biết AC = 12 cm, độ dài đoạn thẳng AN bằng. Kết quả đúng là:\par
\hspace*{1.5em}\textbf{A.} 8 cm\quad\quad \textbf{B.} 4 cm\quad\quad \textbf{C.} 6.0 cm\quad\quad \textbf{D.} 12 cm\par\vspace{2pt}
\textbf{\color{primary} Câu 43.} Vận dụng kiến thức nâng cao trong phạm vi Chương IV để giải bài toán: Cho tam giác ABC có AB = 12 cm, AC = 20 cm, BC = 16 cm. Đường phân giác AD chia cạnh BC thành hai đoạn BD và CD. Độ dài đoạn BD bằng:\par
\hspace*{1.5em}\textbf{A.} 4 cm\quad\quad \textbf{B.} 8 cm\quad\quad \textbf{C.} 6 cm\quad\quad \textbf{D.} 10 cm\par\vspace{2pt}
\textbf{\color{primary} Câu 44.} Phân tích và giải bài toán tổng hợp sau của Chương IV: Cho tam giác ABC, điểm M thuộc cạnh AB sao cho AM = 2$\cdot$MB. Qua M kẻ đường thẳng song song với BC cắt AC tại N. Biết AC = 18 cm, độ dài đoạn thẳng AN bằng:\par
\hspace*{1.5em}\textbf{A.} 18 cm\quad\quad \textbf{B.} 12 cm\quad\quad \textbf{C.} 9.0 cm\quad\quad \textbf{D.} 6 cm\par\vspace{2pt}
\textbf{\color{primary} Câu 45.} Thử thách tư duy trong phạm vi kiến thức Chương IV: Cho tam giác ABC có AB = 9 cm, AC = 15 cm, BC = 16 cm. Đường phân giác AD chia cạnh BC thành hai đoạn BD và CD. Độ dài đoạn BD bằng. Đáp số chính xác là:\par
\hspace*{1.5em}\textbf{A.} 6 cm\quad\quad \textbf{B.} 4 cm\quad\quad \textbf{C.} 10 cm\quad\quad \textbf{D.} 8 cm\par\vspace{2pt}
\textbf{\color{primary} Câu 46.} Sử dụng phương pháp giải toán của Chương IV để tìm đáp án cho bài toán: Cho tam giác ABC, điểm M thuộc cạnh AB sao cho AM = 2$\cdot$MB. Qua M kẻ đường thẳng song song với BC cắt AC tại N. Biết AC = 24 cm, độ dài đoạn thẳng AN bằng:\par
\hspace*{1.5em}\textbf{A.} 24 cm\quad\quad \textbf{B.} 12.0 cm\quad\quad \textbf{C.} 16 cm\quad\quad \textbf{D.} 8 cm\par\vspace{2pt}
\textbf{\color{primary} Câu 47.} Xét bài toán vận dụng thực tiễn / nâng cao thuộc Chương IV: Cho tam giác ABC có AB = 6 cm, AC = 10 cm, BC = 16 cm. Đường phân giác AD chia cạnh BC thành hai đoạn BD và CD. Độ dài đoạn BD bằng:\par
\hspace*{1.5em}\textbf{A.} 10 cm\quad\quad \textbf{B.} 8 cm\quad\quad \textbf{C.} 4 cm\quad\quad \textbf{D.} 6 cm\par\vspace{2pt}
\textbf{\color{primary} Câu 48.} Tìm kết quả cuối cùng của bài toán phân loại học sinh khá giỏi (Chương IV): Cho tam giác ABC, điểm M thuộc cạnh AB sao cho AM = 2$\cdot$MB. Qua M kẻ đường thẳng song song với BC cắt AC tại N. Biết AC = 30 cm, độ dài đoạn thẳng AN bằng:\par
\hspace*{1.5em}\textbf{A.} 10 cm\quad\quad \textbf{B.} 30 cm\quad\quad \textbf{C.} 15.0 cm\quad\quad \textbf{D.} 20 cm\par\vspace{2pt}
\textbf{\color{primary} Câu 49.} Kết hợp các tính chất đã học trong Chương IV, hãy giải quyết bài toán: Cho tam giác ABC có AB = 12 cm, AC = 20 cm, BC = 16 cm. Đường phân giác AD chia cạnh BC thành hai đoạn BD và CD. Độ dài đoạn BD bằng:\par
\hspace*{1.5em}\textbf{A.} 10 cm\quad\quad \textbf{B.} 4 cm\quad\quad \textbf{C.} 8 cm\quad\quad \textbf{D.} 6 cm\par\vspace{2pt}
\textbf{\color{primary} Câu 50.} Câu hỏi vận dụng cao cuối Chương IV: Cho tam giác ABC, điểm M thuộc cạnh AB sao cho AM = 2$\cdot$MB. Qua M kẻ đường thẳng song song với BC cắt AC tại N. Biết AC = 36 cm, độ dài đoạn thẳng AN bằng. Phương án đúng là:\par
\hspace*{1.5em}\textbf{A.} 36 cm\quad\quad \textbf{B.} 12 cm\quad\quad \textbf{C.} 24 cm\quad\quad \textbf{D.} 18.0 cm\par\vspace{2pt}
\subsection*{\color{primary} Bảng đáp án Chương IV: Định lí Thalès}
\begin{center}
\renewcommand{\arraystretch}{1.25}
\begin{tabularx}{\textwidth}{*{10}{>{\centering\arraybackslash}X}}
\toprule
\textbf{\color{primary} Cột 1} & \textbf{\color{primary} Cột 2} & \textbf{\color{primary} Cột 3} & \textbf{\color{primary} Cột 4} & \textbf{\color{primary} Cột 5} & \textbf{\color{primary} Cột 6} & \textbf{\color{primary} Cột 7} & \textbf{\color{primary} Cột 8} & \textbf{\color{primary} Cột 9} & \textbf{\color{primary} Cột 10} \\
\midrule
\textbf{1.B} & \textbf{2.A} & \textbf{3.B} & \textbf{4.D} & \textbf{5.C} & \textbf{6.D} & \textbf{7.C} & \textbf{8.B} & \textbf{9.C} & \textbf{10.A} \\
\rowcolor{rowtint} \textbf{11.B} & \textbf{12.D} & \textbf{13.B} & \textbf{14.B} & \textbf{15.C} & \textbf{16.B} & \textbf{17.B} & \textbf{18.B} & \textbf{19.B} & \textbf{20.C} \\
\textbf{21.B} & \textbf{22.C} & \textbf{23.D} & \textbf{24.A} & \textbf{25.A} & \textbf{26.D} & \textbf{27.C} & \textbf{28.C} & \textbf{29.A} & \textbf{30.B} \\
\rowcolor{rowtint} \textbf{31.C} & \textbf{32.B} & \textbf{33.D} & \textbf{34.A} & \textbf{35.D} & \textbf{36.B} & \textbf{37.A} & \textbf{38.B} & \textbf{39.A} & \textbf{40.D} \\
\textbf{41.A} & \textbf{42.A} & \textbf{43.C} & \textbf{44.B} & \textbf{45.A} & \textbf{46.C} & \textbf{47.D} & \textbf{48.D} & \textbf{49.D} & \textbf{50.C} \\
\bottomrule
\end{tabularx}
\end{center}\vspace{10pt}
\section*{\color{primary} CHƯƠNG V: DỮ LIỆU VÀ BIỂU ĐỒ}
\textit{Trọng tâm kiến thức SGK: Thu thập và phân loại dữ liệu (dữ liệu liên tục, rời rạc), lựa chọn dạng biểu đồ để biểu diễn dữ liệu và phân tích dữ liệu thống kê.}\par\vspace{4pt}
\subsection*{\color{accent} Phần I. Mức độ Dễ (Nhận biết -- Câu 1 đến Câu 20)}
\textbf{\color{primary} Câu 1.} Trong các dữ liệu định lượng sau, dữ liệu nào là dữ liệu liên tục?\par
\begin{itemize}\setlength{\itemsep}{1pt}\setlength{\parskip}{0pt}
  \item[] \textbf{A.} Số quyển vở trong cặp của mỗi bạn
  \item[] \textbf{B.} Số học sinh giỏi của mỗi lớp khối 8
  \item[] \textbf{C.} Số bàn thắng ghi được trong mỗi trận đấu
  \item[] \textbf{D.} Chiều cao (tính bằng cm) của cây trồng trong vườn thí nghiệm
\end{itemize}\vspace{2pt}
\textbf{\color{primary} Câu 2.} Hãy trả lời câu hỏi cơ bản sau: Để so sánh trực quan hai tập dữ liệu cùng loại (ví dụ số học sinh nam và nữ của 4 lớp 8), ta nên chọn loại biểu đồ nào?\par
\hspace*{1.5em}\textbf{A.} Biểu đồ hình quạt tròn\quad\quad \textbf{B.} Biểu đồ đoạn thẳng đơn\quad\quad \textbf{C.} Biểu đồ cột kép\quad\quad \textbf{D.} Danh sách liệt kê\par\vspace{2pt}
\textbf{\color{primary} Câu 3.} Trong bài học thuộc Chương V, xét nội dung: Trong các dữ liệu định lượng sau, dữ liệu nào là dữ liệu rời rạc. Đáp án đúng là:\par
\begin{itemize}\setlength{\itemsep}{1pt}\setlength{\parskip}{0pt}
  \item[] \textbf{A.} Nhiệt độ trung bình mỗi ngày
  \item[] \textbf{B.} Số thành viên trong mỗi hộ gia đình ở tổ dân phố
  \item[] \textbf{C.} Khối lượng của các quả dưa hấu
  \item[] \textbf{D.} Thời gian chạy 100m của các vận động viên
\end{itemize}\vspace{2pt}
\textbf{\color{primary} Câu 4.} Dựa vào kiến thức chuẩn SGK của Chương V (Dữ liệu và biểu đồ), bảng thống kê số cây trồng của 3 tổ trong lớp 8A lần lượt là 15 cây, 18 cây và 21 cây. Tổng số cây cả 3 tổ trồng được là:\par
\hspace*{1.5em}\textbf{A.} 54 cây\quad\quad \textbf{B.} 18 cây\quad\quad \textbf{C.} 59 cây\quad\quad \textbf{D.} 49 cây\par\vspace{2pt}
\textbf{\color{primary} Câu 5.} Chọn kết quả đúng nhất khi xét: Biểu đồ hình quạt tròn biểu diễn cơ cấu chi tiêu của một gia đình, biết mục Giáo dục chiếm 22\%, mục Ăn uống chiếm 45\%. Tổng tỉ lệ phần trăm của các mục còn lại là:\par
\hspace*{1.5em}\textbf{A.} 43\%\quad\quad \textbf{B.} 51\%\quad\quad \textbf{C.} 23\%\quad\quad \textbf{D.} 33\%\par\vspace{2pt}
\textbf{\color{primary} Câu 6.} Khi làm bài tập cơ bản về dữ liệu và biểu đồ, một học sinh gặp câu hỏi: Trong các dữ liệu định lượng sau, dữ liệu nào là dữ liệu liên tục. Kết quả nào dưới đây đúng?\par
\begin{itemize}\setlength{\itemsep}{1pt}\setlength{\parskip}{0pt}
  \item[] \textbf{A.} Số bàn thắng ghi được trong mỗi trận đấu
  \item[] \textbf{B.} Số quyển vở trong cặp của mỗi bạn
  \item[] \textbf{C.} Chiều cao (tính bằng cm) của cây trồng trong vườn thí nghiệm
  \item[] \textbf{D.} Số học sinh giỏi của mỗi lớp khối 8
\end{itemize}\vspace{2pt}
\textbf{\color{primary} Câu 7.} Xác định giá trị hoặc khẳng định đúng cho bài toán sau: Để so sánh trực quan hai tập dữ liệu cùng loại (ví dụ số học sinh nam và nữ của 4 lớp 8), ta nên chọn loại biểu đồ nào:\par
\hspace*{1.5em}\textbf{A.} Biểu đồ đoạn thẳng đơn\quad\quad \textbf{B.} Biểu đồ cột kép\quad\quad \textbf{C.} Danh sách liệt kê\quad\quad \textbf{D.} Biểu đồ hình quạt tròn\par\vspace{2pt}
\textbf{\color{primary} Câu 8.} Áp dụng trực tiếp định nghĩa và quy tắc trong Chương V: Trong các dữ liệu định lượng sau, dữ liệu nào là dữ liệu rời rạc:\par
\begin{itemize}\setlength{\itemsep}{1pt}\setlength{\parskip}{0pt}
  \item[] \textbf{A.} Nhiệt độ trung bình mỗi ngày
  \item[] \textbf{B.} Thời gian chạy 100m của các vận động viên
  \item[] \textbf{C.} Khối lượng của các quả dưa hấu
  \item[] \textbf{D.} Số thành viên trong mỗi hộ gia đình ở tổ dân phố
\end{itemize}\vspace{2pt}
\textbf{\color{primary} Câu 9.} Kiểm tra nhanh kiến thức nền tảng: Bảng thống kê số cây trồng của 3 tổ trong lớp 8A lần lượt là 20 cây, 23 cây và 26 cây. Tổng số cây cả 3 tổ trồng được là:\par
\hspace*{1.5em}\textbf{A.} 69 cây\quad\quad \textbf{B.} 64 cây\quad\quad \textbf{C.} 74 cây\quad\quad \textbf{D.} 23 cây\par\vspace{2pt}
\textbf{\color{primary} Câu 10.} Cho dữ kiện nhận biết sau đây: Biểu đồ hình quạt tròn biểu diễn cơ cấu chi tiêu của một gia đình, biết mục Giáo dục chiếm 27\%, mục Ăn uống chiếm 45\%. Tổng tỉ lệ phần trăm của các mục còn lại là. Phương án chính xác là:\par
\hspace*{1.5em}\textbf{A.} 46\%\quad\quad \textbf{B.} 28\%\quad\quad \textbf{C.} 18\%\quad\quad \textbf{D.} 38\%\par\vspace{2pt}
\textbf{\color{primary} Câu 11.} Hãy chỉ ra kết quả đúng trong bốn lựa chọn A, B, C, D đối với yêu cầu: Trong các dữ liệu định lượng sau, dữ liệu nào là dữ liệu liên tục:\par
\begin{itemize}\setlength{\itemsep}{1pt}\setlength{\parskip}{0pt}
  \item[] \textbf{A.} Số học sinh giỏi của mỗi lớp khối 8
  \item[] \textbf{B.} Số quyển vở trong cặp của mỗi bạn
  \item[] \textbf{C.} Số bàn thắng ghi được trong mỗi trận đấu
  \item[] \textbf{D.} Chiều cao (tính bằng cm) của cây trồng trong vườn thí nghiệm
\end{itemize}\vspace{2pt}
\textbf{\color{primary} Câu 12.} Câu hỏi ôn tập lý thuyết và kỹ năng cơ bản (Chương V): Để so sánh trực quan hai tập dữ liệu cùng loại (ví dụ số học sinh nam và nữ của 4 lớp 8), ta nên chọn loại biểu đồ nào:\par
\hspace*{1.5em}\textbf{A.} Biểu đồ đoạn thẳng đơn\quad\quad \textbf{B.} Danh sách liệt kê\quad\quad \textbf{C.} Biểu đồ cột kép\quad\quad \textbf{D.} Biểu đồ hình quạt tròn\par\vspace{2pt}
\textbf{\color{primary} Câu 13.} Theo nội dung sách giáo khoa Kết nối tri thức với cuộc sống (Chương V), trong các dữ liệu định lượng sau, dữ liệu nào là dữ liệu rời rạc:\par
\begin{itemize}\setlength{\itemsep}{1pt}\setlength{\parskip}{0pt}
  \item[] \textbf{A.} Thời gian chạy 100m của các vận động viên
  \item[] \textbf{B.} Nhiệt độ trung bình mỗi ngày
  \item[] \textbf{C.} Khối lượng của các quả dưa hấu
  \item[] \textbf{D.} Số thành viên trong mỗi hộ gia đình ở tổ dân phố
\end{itemize}\vspace{2pt}
\textbf{\color{primary} Câu 14.} Tìm đáp số đúng của bài toán cơ bản sau: Bảng thống kê số cây trồng của 3 tổ trong lớp 8A lần lượt là 25 cây, 28 cây và 31 cây. Tổng số cây cả 3 tổ trồng được là:\par
\hspace*{1.5em}\textbf{A.} 79 cây\quad\quad \textbf{B.} 84 cây\quad\quad \textbf{C.} 28 cây\quad\quad \textbf{D.} 89 cây\par\vspace{2pt}
\textbf{\color{primary} Câu 15.} Xét đối tượng toán học trong Chương V: Biểu đồ hình quạt tròn biểu diễn cơ cấu chi tiêu của một gia đình, biết mục Giáo dục chiếm 32\%, mục Ăn uống chiếm 45\%. Tổng tỉ lệ phần trăm của các mục còn lại là. Giá trị hoặc tính chất tương ứng là:\par
\hspace*{1.5em}\textbf{A.} 13\%\quad\quad \textbf{B.} 23\%\quad\quad \textbf{C.} 41\%\quad\quad \textbf{D.} 33\%\par\vspace{2pt}
\textbf{\color{primary} Câu 16.} Hoàn thành phép xác định cơ bản sau đây: Trong các dữ liệu định lượng sau, dữ liệu nào là dữ liệu liên tục:\par
\begin{itemize}\setlength{\itemsep}{1pt}\setlength{\parskip}{0pt}
  \item[] \textbf{A.} Chiều cao (tính bằng cm) của cây trồng trong vườn thí nghiệm
  \item[] \textbf{B.} Số bàn thắng ghi được trong mỗi trận đấu
  \item[] \textbf{C.} Số quyển vở trong cặp của mỗi bạn
  \item[] \textbf{D.} Số học sinh giỏi của mỗi lớp khối 8
\end{itemize}\vspace{2pt}
\textbf{\color{primary} Câu 17.} Trong các khẳng định hoặc giá trị dưới đây, đâu là đáp án đúng cho câu hỏi: Để so sánh trực quan hai tập dữ liệu cùng loại (ví dụ số học sinh nam và nữ của 4 lớp 8), ta nên chọn loại biểu đồ nào?\par
\hspace*{1.5em}\textbf{A.} Biểu đồ đoạn thẳng đơn\quad\quad \textbf{B.} Danh sách liệt kê\quad\quad \textbf{C.} Biểu đồ hình quạt tròn\quad\quad \textbf{D.} Biểu đồ cột kép\par\vspace{2pt}
\textbf{\color{primary} Câu 18.} Bài tập khởi động Chương V: Trong các dữ liệu định lượng sau, dữ liệu nào là dữ liệu rời rạc:\par
\begin{itemize}\setlength{\itemsep}{1pt}\setlength{\parskip}{0pt}
  \item[] \textbf{A.} Thời gian chạy 100m của các vận động viên
  \item[] \textbf{B.} Số thành viên trong mỗi hộ gia đình ở tổ dân phố
  \item[] \textbf{C.} Khối lượng của các quả dưa hấu
  \item[] \textbf{D.} Nhiệt độ trung bình mỗi ngày
\end{itemize}\vspace{2pt}
\textbf{\color{primary} Câu 19.} Quan sát dữ kiện đề bài: Bảng thống kê số cây trồng của 3 tổ trong lớp 8A lần lượt là 30 cây, 33 cây và 36 cây. Tổng số cây cả 3 tổ trồng được là. Hãy chọn phương án đúng:\par
\hspace*{1.5em}\textbf{A.} 99 cây\quad\quad \textbf{B.} 33 cây\quad\quad \textbf{C.} 104 cây\quad\quad \textbf{D.} 94 cây\par\vspace{2pt}
\textbf{\color{primary} Câu 20.} Sử dụng công thức hoặc tính chất đầu bài của Chương V để giải quyết yêu cầu: Biểu đồ hình quạt tròn biểu diễn cơ cấu chi tiêu của một gia đình, biết mục Giáo dục chiếm 22\%, mục Ăn uống chiếm 45\%. Tổng tỉ lệ phần trăm của các mục còn lại là:\par
\hspace*{1.5em}\textbf{A.} 33\%\quad\quad \textbf{B.} 43\%\quad\quad \textbf{C.} 23\%\quad\quad \textbf{D.} 51\%\par\vspace{2pt}
\subsection*{\color{accent} Phần II. Mức độ Trung bình (Thông hiểu -- Câu 21 đến Câu 40)}
\textbf{\color{primary} Câu 21.} Khảo sát 200 học sinh khối 8 về phương tiện đi học: Đi xe đạp chiếm 60\%, Đi bộ chiếm 25\%, còn lại là Bố mẹ đưa đón. Số học sinh được bố mẹ đưa đón là:\par
\hspace*{1.5em}\textbf{A.} 120 học sinh\quad\quad \textbf{B.} 30 học sinh\quad\quad \textbf{C.} 40 học sinh\quad\quad \textbf{D.} 50 học sinh\par\vspace{2pt}
\textbf{\color{primary} Câu 22.} Vận dụng quy tắc trong Chương V (Dữ liệu và biểu đồ) để giải bài toán: Sản lượng lúa của một xã trong hai vụ Chiêm và Mùa lần lượt là 320 tấn và 520 tấn. Tỉ số phần trăm sản lượng vụ Chiêm so với tổng sản lượng cả năm là:\par
\hspace*{1.5em}\textbf{A.} 61.9\%\quad\quad \textbf{B.} 61.5\%\quad\quad \textbf{C.} 50,0\%\quad\quad \textbf{D.} 38.1\%\par\vspace{2pt}
\textbf{\color{primary} Câu 23.} Biến đổi và tính toán theo kiến thức Chương V: Để thu thập dữ liệu về 'Mức tiêu thụ điện trung bình tháng của các hộ dân tại thành phố Đà Nẵng', phương pháp thu thập nào là thu thập gián tiếp. Kết quả thu được là:\par
\begin{itemize}\setlength{\itemsep}{1pt}\setlength{\parskip}{0pt}
  \item[] \textbf{A.} Tra cứu số liệu thống kê công bố trên Cổng thông tin của Công ty Điện lực
  \item[] \textbf{B.} Đi phỏng vấn trực tiếp từng hộ dân
  \item[] \textbf{C.} Phát phiếu hỏi trực tiếp tại khu phố
  \item[] \textbf{D.} Tự ghi chỉ số công tơ điện tại nhà
\end{itemize}\vspace{2pt}
\textbf{\color{primary} Câu 24.} Xét bài toán mức độ thông hiểu sau: Khảo sát 350 học sinh khối 8 về phương tiện đi học: Đi xe đạp chiếm 60\%, Đi bộ chiếm 25\%, còn lại là Bố mẹ đưa đón. Số học sinh được bố mẹ đưa đón là. Đáp án đúng là:\par
\hspace*{1.5em}\textbf{A.} 87 học sinh\quad\quad \textbf{B.} 52 học sinh\quad\quad \textbf{C.} 210 học sinh\quad\quad \textbf{D.} 70 học sinh\par\vspace{2pt}
\textbf{\color{primary} Câu 25.} Để hoàn thành phiếu học tập Chương V, hãy xác định: Sản lượng lúa của một xã trong hai vụ Chiêm và Mùa lần lượt là 380 tấn và 580 tấn. Tỉ số phần trăm sản lượng vụ Chiêm so với tổng sản lượng cả năm là:\par
\hspace*{1.5em}\textbf{A.} 60.4\%\quad\quad \textbf{B.} 65.5\%\quad\quad \textbf{C.} 50,0\%\quad\quad \textbf{D.} 39.6\%\par\vspace{2pt}
\textbf{\color{primary} Câu 26.} Phân tích dữ kiện bài toán: Để thu thập dữ liệu về 'Mức tiêu thụ điện trung bình tháng của các hộ dân tại thành phố Đà Nẵng', phương pháp thu thập nào là thu thập gián tiếp. Giá trị cần tìm bằng:\par
\begin{itemize}\setlength{\itemsep}{1pt}\setlength{\parskip}{0pt}
  \item[] \textbf{A.} Phát phiếu hỏi trực tiếp tại khu phố
  \item[] \textbf{B.} Tra cứu số liệu thống kê công bố trên Cổng thông tin của Công ty Điện lực
  \item[] \textbf{C.} Đi phỏng vấn trực tiếp từng hộ dân
  \item[] \textbf{D.} Tự ghi chỉ số công tơ điện tại nhà
\end{itemize}\vspace{2pt}
\textbf{\color{primary} Câu 27.} Áp dụng tính chất đã học ở Chương V vào bài tập sau: Khảo sát 300 học sinh khối 8 về phương tiện đi học: Đi xe đạp chiếm 60\%, Đi bộ chiếm 25\%, còn lại là Bố mẹ đưa đón. Số học sinh được bố mẹ đưa đón là:\par
\hspace*{1.5em}\textbf{A.} 45 học sinh\quad\quad \textbf{B.} 60 học sinh\quad\quad \textbf{C.} 75 học sinh\quad\quad \textbf{D.} 180 học sinh\par\vspace{2pt}
\textbf{\color{primary} Câu 28.} Thực hiện các bước suy luận hợp lý để giải quyết yêu cầu: Sản lượng lúa của một xã trong hai vụ Chiêm và Mùa lần lượt là 440 tấn và 640 tấn. Tỉ số phần trăm sản lượng vụ Chiêm so với tổng sản lượng cả năm là:\par
\hspace*{1.5em}\textbf{A.} 68.8\%\quad\quad \textbf{B.} 59.3\%\quad\quad \textbf{C.} 40.7\%\quad\quad \textbf{D.} 50,0\%\par\vspace{2pt}
\textbf{\color{primary} Câu 29.} Cho bài toán tự luyện (Dữ liệu và biểu đồ): Để thu thập dữ liệu về 'Mức tiêu thụ điện trung bình tháng của các hộ dân tại thành phố Đà Nẵng', phương pháp thu thập nào là thu thập gián tiếp. Phương án nào sau đây chính xác?\par
\begin{itemize}\setlength{\itemsep}{1pt}\setlength{\parskip}{0pt}
  \item[] \textbf{A.} Tra cứu số liệu thống kê công bố trên Cổng thông tin của Công ty Điện lực
  \item[] \textbf{B.} Đi phỏng vấn trực tiếp từng hộ dân
  \item[] \textbf{C.} Phát phiếu hỏi trực tiếp tại khu phố
  \item[] \textbf{D.} Tự ghi chỉ số công tơ điện tại nhà
\end{itemize}\vspace{2pt}
\textbf{\color{primary} Câu 30.} Kết hợp các phép biến đổi trong Chương V, hãy tìm kết quả của bài toán: Khảo sát 250 học sinh khối 8 về phương tiện đi học: Đi xe đạp chiếm 60\%, Đi bộ chiếm 25\%, còn lại là Bố mẹ đưa đón. Số học sinh được bố mẹ đưa đón là:\par
\hspace*{1.5em}\textbf{A.} 37 học sinh\quad\quad \textbf{B.} 50 học sinh\quad\quad \textbf{C.} 62 học sinh\quad\quad \textbf{D.} 150 học sinh\par\vspace{2pt}
\textbf{\color{primary} Câu 31.} Một bài kiểm tra định kỳ nội dung Chương V nêu câu hỏi: Sản lượng lúa của một xã trong hai vụ Chiêm và Mùa lần lượt là 500 tấn và 700 tấn. Tỉ số phần trăm sản lượng vụ Chiêm so với tổng sản lượng cả năm là. Hãy chọn đáp số đúng:\par
\hspace*{1.5em}\textbf{A.} 58.3\%\quad\quad \textbf{B.} 41.7\%\quad\quad \textbf{C.} 50,0\%\quad\quad \textbf{D.} 71.4\%\par\vspace{2pt}
\textbf{\color{primary} Câu 32.} Tìm lời giải đúng cho bài toán thông hiểu: Để thu thập dữ liệu về 'Mức tiêu thụ điện trung bình tháng của các hộ dân tại thành phố Đà Nẵng', phương pháp thu thập nào là thu thập gián tiếp:\par
\begin{itemize}\setlength{\itemsep}{1pt}\setlength{\parskip}{0pt}
  \item[] \textbf{A.} Tự ghi chỉ số công tơ điện tại nhà
  \item[] \textbf{B.} Phát phiếu hỏi trực tiếp tại khu phố
  \item[] \textbf{C.} Đi phỏng vấn trực tiếp từng hộ dân
  \item[] \textbf{D.} Tra cứu số liệu thống kê công bố trên Cổng thông tin của Công ty Điện lực
\end{itemize}\vspace{2pt}
\textbf{\color{primary} Câu 33.} Sử dụng mối liên hệ giữa các đại lượng trong Chương V: Khảo sát 200 học sinh khối 8 về phương tiện đi học: Đi xe đạp chiếm 60\%, Đi bộ chiếm 25\%, còn lại là Bố mẹ đưa đón. Số học sinh được bố mẹ đưa đón là:\par
\hspace*{1.5em}\textbf{A.} 30 học sinh\quad\quad \textbf{B.} 50 học sinh\quad\quad \textbf{C.} 120 học sinh\quad\quad \textbf{D.} 40 học sinh\par\vspace{2pt}
\textbf{\color{primary} Câu 34.} Hãy tính toán và chọn phương án đúng đối với bài tập: Sản lượng lúa của một xã trong hai vụ Chiêm và Mùa lần lượt là 560 tấn và 760 tấn. Tỉ số phần trăm sản lượng vụ Chiêm so với tổng sản lượng cả năm là:\par
\hspace*{1.5em}\textbf{A.} 50,0\%\quad\quad \textbf{B.} 73.7\%\quad\quad \textbf{C.} 57.6\%\quad\quad \textbf{D.} 42.4\%\par\vspace{2pt}
\textbf{\color{primary} Câu 35.} Dựa trên các tính chất trọng tâm của Chương V, xét bài toán: Để thu thập dữ liệu về 'Mức tiêu thụ điện trung bình tháng của các hộ dân tại thành phố Đà Nẵng', phương pháp thu thập nào là thu thập gián tiếp:\par
\begin{itemize}\setlength{\itemsep}{1pt}\setlength{\parskip}{0pt}
  \item[] \textbf{A.} Đi phỏng vấn trực tiếp từng hộ dân
  \item[] \textbf{B.} Tự ghi chỉ số công tơ điện tại nhà
  \item[] \textbf{C.} Phát phiếu hỏi trực tiếp tại khu phố
  \item[] \textbf{D.} Tra cứu số liệu thống kê công bố trên Cổng thông tin của Công ty Điện lực
\end{itemize}\vspace{2pt}
\textbf{\color{primary} Câu 36.} Thực hành giải toán (Dữ liệu và biểu đồ): Khảo sát 350 học sinh khối 8 về phương tiện đi học: Đi xe đạp chiếm 60\%, Đi bộ chiếm 25\%, còn lại là Bố mẹ đưa đón. Số học sinh được bố mẹ đưa đón là. Đáp án nào dưới đây thỏa mãn?\par
\hspace*{1.5em}\textbf{A.} 70 học sinh\quad\quad \textbf{B.} 52 học sinh\quad\quad \textbf{C.} 87 học sinh\quad\quad \textbf{D.} 210 học sinh\par\vspace{2pt}
\textbf{\color{primary} Câu 37.} Kiểm tra kỹ năng biến đổi và suy luận qua bài toán: Sản lượng lúa của một xã trong hai vụ Chiêm và Mùa lần lượt là 620 tấn và 820 tấn. Tỉ số phần trăm sản lượng vụ Chiêm so với tổng sản lượng cả năm là:\par
\hspace*{1.5em}\textbf{A.} 50,0\%\quad\quad \textbf{B.} 43.1\%\quad\quad \textbf{C.} 75.6\%\quad\quad \textbf{D.} 56.9\%\par\vspace{2pt}
\textbf{\color{primary} Câu 38.} Cho giả thiết thuộc phạm vi Chương V: Để thu thập dữ liệu về 'Mức tiêu thụ điện trung bình tháng của các hộ dân tại thành phố Đà Nẵng', phương pháp thu thập nào là thu thập gián tiếp. Kết luận đúng là:\par
\begin{itemize}\setlength{\itemsep}{1pt}\setlength{\parskip}{0pt}
  \item[] \textbf{A.} Phát phiếu hỏi trực tiếp tại khu phố
  \item[] \textbf{B.} Tự ghi chỉ số công tơ điện tại nhà
  \item[] \textbf{C.} Tra cứu số liệu thống kê công bố trên Cổng thông tin của Công ty Điện lực
  \item[] \textbf{D.} Đi phỏng vấn trực tiếp từng hộ dân
\end{itemize}\vspace{2pt}
\textbf{\color{primary} Câu 39.} Rút gọn hoặc tính giá trị theo yêu cầu sau của Chương V: Khảo sát 300 học sinh khối 8 về phương tiện đi học: Đi xe đạp chiếm 60\%, Đi bộ chiếm 25\%, còn lại là Bố mẹ đưa đón. Số học sinh được bố mẹ đưa đón là:\par
\hspace*{1.5em}\textbf{A.} 75 học sinh\quad\quad \textbf{B.} 180 học sinh\quad\quad \textbf{C.} 60 học sinh\quad\quad \textbf{D.} 45 học sinh\par\vspace{2pt}
\textbf{\color{primary} Câu 40.} Xác định phương án đúng sau khi thực hiện phép tính / suy luận: Sản lượng lúa của một xã trong hai vụ Chiêm và Mùa lần lượt là 680 tấn và 880 tấn. Tỉ số phần trăm sản lượng vụ Chiêm so với tổng sản lượng cả năm là:\par
\hspace*{1.5em}\textbf{A.} 43.6\%\quad\quad \textbf{B.} 50,0\%\quad\quad \textbf{C.} 56.4\%\quad\quad \textbf{D.} 77.3\%\par\vspace{2pt}
\subsection*{\color{accent} Phần III. Mức độ Khó (Vận dụng \& Vận dụng cao -- Câu 41 đến Câu 50)}
\textbf{\color{primary} Câu 41.} Tổng chi tiêu tháng 1 của một gia đình là 12 triệu đồng, trong đó tiền ăn chiếm 40\%. Sang tháng 2, tổng chi tiêu tăng 10\% nhưng tỉ lệ tiền ăn chiếm 45\% tổng chi tiêu tháng 2. Số tiền ăn tháng 2 bằng:\par
\hspace*{1.5em}\textbf{A.} 6.60 triệu đồng\quad\quad \textbf{B.} 5.94 triệu đồng\quad\quad \textbf{C.} 5.28 triệu đồng\quad\quad \textbf{D.} 5.40 triệu đồng\par\vspace{2pt}
\textbf{\color{primary} Câu 42.} Bài toán vận dụng (Chương V – Dữ liệu và biểu đồ): Điểm trung bình môn Toán của tổ 1 (gồm 10 học sinh) là 8,0 và của tổ 2 (gồm 10 học sinh) là 7.2. Điểm trung bình chung của cả 20 học sinh hai tổ là. Kết quả đúng là:\par
\hspace*{1.5em}\textbf{A.} 7.20\quad\quad \textbf{B.} 8.00\quad\quad \textbf{C.} 7.60\quad\quad \textbf{D.} 8.10\par\vspace{2pt}
\textbf{\color{primary} Câu 43.} Vận dụng kiến thức nâng cao trong phạm vi Chương V để giải bài toán: Tổng chi tiêu tháng 1 của một gia đình là 16 triệu đồng, trong đó tiền ăn chiếm 40\%. Sang tháng 2, tổng chi tiêu tăng 10\% nhưng tỉ lệ tiền ăn chiếm 45\% tổng chi tiêu tháng 2. Số tiền ăn tháng 2 bằng:\par
\hspace*{1.5em}\textbf{A.} 7.04 triệu đồng\quad\quad \textbf{B.} 7.20 triệu đồng\quad\quad \textbf{C.} 7.92 triệu đồng\quad\quad \textbf{D.} 8.80 triệu đồng\par\vspace{2pt}
\textbf{\color{primary} Câu 44.} Phân tích và giải bài toán tổng hợp sau của Chương V: Điểm trung bình môn Toán của tổ 1 (gồm 10 học sinh) là 8,0 và của tổ 2 (gồm 10 học sinh) là 7.6. Điểm trung bình chung của cả 20 học sinh hai tổ là:\par
\hspace*{1.5em}\textbf{A.} 7.80\quad\quad \textbf{B.} 8.00\quad\quad \textbf{C.} 7.60\quad\quad \textbf{D.} 8.30\par\vspace{2pt}
\textbf{\color{primary} Câu 45.} Thử thách tư duy trong phạm vi kiến thức Chương V: Tổng chi tiêu tháng 1 của một gia đình là 12 triệu đồng, trong đó tiền ăn chiếm 40\%. Sang tháng 2, tổng chi tiêu tăng 10\% nhưng tỉ lệ tiền ăn chiếm 45\% tổng chi tiêu tháng 2. Số tiền ăn tháng 2 bằng. Đáp số chính xác là:\par
\hspace*{1.5em}\textbf{A.} 5.40 triệu đồng\quad\quad \textbf{B.} 5.28 triệu đồng\quad\quad \textbf{C.} 5.94 triệu đồng\quad\quad \textbf{D.} 6.60 triệu đồng\par\vspace{2pt}
\textbf{\color{primary} Câu 46.} Sử dụng phương pháp giải toán của Chương V để tìm đáp án cho bài toán: Điểm trung bình môn Toán của tổ 1 (gồm 10 học sinh) là 8,0 và của tổ 2 (gồm 10 học sinh) là 7.0. Điểm trung bình chung của cả 20 học sinh hai tổ là:\par
\hspace*{1.5em}\textbf{A.} 7.50 (1)\quad\quad \textbf{B.} 8.00\quad\quad \textbf{C.} 7.00\quad\quad \textbf{D.} 7.50\par\vspace{2pt}
\textbf{\color{primary} Câu 47.} Xét bài toán vận dụng thực tiễn / nâng cao thuộc Chương V: Tổng chi tiêu tháng 1 của một gia đình là 16 triệu đồng, trong đó tiền ăn chiếm 40\%. Sang tháng 2, tổng chi tiêu tăng 10\% nhưng tỉ lệ tiền ăn chiếm 45\% tổng chi tiêu tháng 2. Số tiền ăn tháng 2 bằng:\par
\hspace*{1.5em}\textbf{A.} 7.20 triệu đồng\quad\quad \textbf{B.} 8.80 triệu đồng\quad\quad \textbf{C.} 7.92 triệu đồng\quad\quad \textbf{D.} 7.04 triệu đồng\par\vspace{2pt}
\textbf{\color{primary} Câu 48.} Tìm kết quả cuối cùng của bài toán phân loại học sinh khá giỏi (Chương V): Điểm trung bình môn Toán của tổ 1 (gồm 10 học sinh) là 8,0 và của tổ 2 (gồm 10 học sinh) là 7.4. Điểm trung bình chung của cả 20 học sinh hai tổ là:\par
\hspace*{1.5em}\textbf{A.} 7.70\quad\quad \textbf{B.} 8.00\quad\quad \textbf{C.} 8.20\quad\quad \textbf{D.} 7.40\par\vspace{2pt}
\textbf{\color{primary} Câu 49.} Kết hợp các tính chất đã học trong Chương V, hãy giải quyết bài toán: Tổng chi tiêu tháng 1 của một gia đình là 12 triệu đồng, trong đó tiền ăn chiếm 40\%. Sang tháng 2, tổng chi tiêu tăng 10\% nhưng tỉ lệ tiền ăn chiếm 45\% tổng chi tiêu tháng 2. Số tiền ăn tháng 2 bằng:\par
\hspace*{1.5em}\textbf{A.} 6.60 triệu đồng\quad\quad \textbf{B.} 5.94 triệu đồng\quad\quad \textbf{C.} 5.40 triệu đồng\quad\quad \textbf{D.} 5.28 triệu đồng\par\vspace{2pt}
\textbf{\color{primary} Câu 50.} Câu hỏi vận dụng cao cuối Chương V: Điểm trung bình môn Toán của tổ 1 (gồm 10 học sinh) là 8,0 và của tổ 2 (gồm 10 học sinh) là 7.8. Điểm trung bình chung của cả 20 học sinh hai tổ là. Phương án đúng là:\par
\hspace*{1.5em}\textbf{A.} 7.90\quad\quad \textbf{B.} 7.80\quad\quad \textbf{C.} 8.00\quad\quad \textbf{D.} 8.40\par\vspace{2pt}
\subsection*{\color{primary} Bảng đáp án Chương V: Dữ liệu và biểu đồ}
\begin{center}
\renewcommand{\arraystretch}{1.25}
\begin{tabularx}{\textwidth}{*{10}{>{\centering\arraybackslash}X}}
\toprule
\textbf{\color{primary} Cột 1} & \textbf{\color{primary} Cột 2} & \textbf{\color{primary} Cột 3} & \textbf{\color{primary} Cột 4} & \textbf{\color{primary} Cột 5} & \textbf{\color{primary} Cột 6} & \textbf{\color{primary} Cột 7} & \textbf{\color{primary} Cột 8} & \textbf{\color{primary} Cột 9} & \textbf{\color{primary} Cột 10} \\
\midrule
\textbf{1.D} & \textbf{2.C} & \textbf{3.B} & \textbf{4.A} & \textbf{5.D} & \textbf{6.C} & \textbf{7.B} & \textbf{8.D} & \textbf{9.A} & \textbf{10.B} \\
\rowcolor{rowtint} \textbf{11.D} & \textbf{12.C} & \textbf{13.D} & \textbf{14.B} & \textbf{15.B} & \textbf{16.A} & \textbf{17.D} & \textbf{18.B} & \textbf{19.A} & \textbf{20.A} \\
\textbf{21.B} & \textbf{22.D} & \textbf{23.A} & \textbf{24.B} & \textbf{25.D} & \textbf{26.B} & \textbf{27.A} & \textbf{28.C} & \textbf{29.A} & \textbf{30.A} \\
\rowcolor{rowtint} \textbf{31.B} & \textbf{32.D} & \textbf{33.A} & \textbf{34.D} & \textbf{35.D} & \textbf{36.B} & \textbf{37.B} & \textbf{38.C} & \textbf{39.D} & \textbf{40.A} \\
\textbf{41.B} & \textbf{42.C} & \textbf{43.C} & \textbf{44.A} & \textbf{45.C} & \textbf{46.D} & \textbf{47.C} & \textbf{48.A} & \textbf{49.B} & \textbf{50.A} \\
\bottomrule
\end{tabularx}
\end{center}\vspace{10pt}
\section*{\color{primary} CHƯƠNG VI: PHÂN THỨC ĐẠI SỐ}
\textit{Trọng tâm kiến thức SGK: Phân thức đại số, tính chất cơ bản của phân thức, rút gọn phân thức, quy đồng mẫu thức và các phép cộng, trừ, nhân, chia phân thức đại số.}\par\vspace{4pt}
\subsection*{\color{accent} Phần I. Mức độ Dễ (Nhận biết -- Câu 1 đến Câu 20)}
\textbf{\color{primary} Câu 1.} Điều kiện xác định của phân thức đại số (2x + 1)/(x - 2) là:\par
\hspace*{1.5em}\textbf{A.} x $\ne$ -1/2\quad\quad \textbf{B.} x $\ne$ -2\quad\quad \textbf{C.} x $\ne$ 2\quad\quad \textbf{D.} x = 2\par\vspace{2pt}
\textbf{\color{primary} Câu 2.} Thực hiện yêu cầu nhận biết sau: Rút gọn phân thức 6x$^2$y / (9xy$^2$) với x $\ne$ 0, y $\ne$ 0 ta được.\par
\hspace*{1.5em}\textbf{A.} 2y / (3x)\quad\quad \textbf{B.} 2x / (3y)\quad\quad \textbf{C.} 3x / (2y)\quad\quad \textbf{D.} 2x$^2$ / (3y$^2$)\par\vspace{2pt}
\textbf{\color{primary} Câu 3.} Trong bài học thuộc Chương VI, xét nội dung: Phân thức đối của phân thức (x - 5)/(x$^2$ + 1) là. Đáp án đúng là:\par
\hspace*{1.5em}\textbf{A.} (-x - 5)/(x$^2$ + 1)\quad\quad \textbf{B.} (x$^2$ + 1)/(x - 5)\quad\quad \textbf{C.} (5 - x)/(x$^2$ + 1)\quad\quad \textbf{D.} (x + 5)/(x$^2$ + 1)\par\vspace{2pt}
\textbf{\color{primary} Câu 4.} Dựa vào kiến thức chuẩn SGK của Chương VI (Phân thức đại số), điều kiện xác định của phân thức đại số (2x + 1)/(x - 5) là:\par
\hspace*{1.5em}\textbf{A.} x $\ne$ -1/2\quad\quad \textbf{B.} x $\ne$ -5\quad\quad \textbf{C.} x $\ne$ 5\quad\quad \textbf{D.} x = 5\par\vspace{2pt}
\textbf{\color{primary} Câu 5.} Chọn kết quả đúng nhất khi xét: Rút gọn phân thức 12x$^2$y / (18xy$^2$) với x $\ne$ 0, y $\ne$ 0 ta được:\par
\hspace*{1.5em}\textbf{A.} 2y / (3x)\quad\quad \textbf{B.} 2x$^2$ / (3y$^2$)\quad\quad \textbf{C.} 2x / (3y)\quad\quad \textbf{D.} 3x / (2y)\par\vspace{2pt}
\textbf{\color{primary} Câu 6.} Khi làm bài tập cơ bản về phân thức đại số, một học sinh gặp câu hỏi: Phân thức đối của phân thức (x - 8)/(x$^2$ + 1) là. Kết quả nào dưới đây đúng?\par
\hspace*{1.5em}\textbf{A.} (x + 8)/(x$^2$ + 1)\quad\quad \textbf{B.} (8 - x)/(x$^2$ + 1)\quad\quad \textbf{C.} (x$^2$ + 1)/(x - 8)\quad\quad \textbf{D.} (-x - 8)/(x$^2$ + 1)\par\vspace{2pt}
\textbf{\color{primary} Câu 7.} Xác định giá trị hoặc khẳng định đúng cho bài toán sau: Điều kiện xác định của phân thức đại số (2x + 1)/(x - 8) là:\par
\hspace*{1.5em}\textbf{A.} x = 8\quad\quad \textbf{B.} x $\ne$ 8\quad\quad \textbf{C.} x $\ne$ -1/2\quad\quad \textbf{D.} x $\ne$ -8\par\vspace{2pt}
\textbf{\color{primary} Câu 8.} Áp dụng trực tiếp định nghĩa và quy tắc trong Chương VI: Rút gọn phân thức 18x$^2$y / (27xy$^2$) với x $\ne$ 0, y $\ne$ 0 ta được:\par
\hspace*{1.5em}\textbf{A.} 2x$^2$ / (3y$^2$)\quad\quad \textbf{B.} 3x / (2y)\quad\quad \textbf{C.} 2y / (3x)\quad\quad \textbf{D.} 2x / (3y)\par\vspace{2pt}
\textbf{\color{primary} Câu 9.} Kiểm tra nhanh kiến thức nền tảng: Phân thức đối của phân thức (x - 11)/(x$^2$ + 1) là:\par
\hspace*{1.5em}\textbf{A.} (11 - x)/(x$^2$ + 1)\quad\quad \textbf{B.} (x$^2$ + 1)/(x - 11)\quad\quad \textbf{C.} (x + 11)/(x$^2$ + 1)\quad\quad \textbf{D.} (-x - 11)/(x$^2$ + 1)\par\vspace{2pt}
\textbf{\color{primary} Câu 10.} Cho dữ kiện nhận biết sau đây: Điều kiện xác định của phân thức đại số (2x + 1)/(x - 11) là. Phương án chính xác là:\par
\hspace*{1.5em}\textbf{A.} x = 11\quad\quad \textbf{B.} x $\ne$ -11\quad\quad \textbf{C.} x $\ne$ 11\quad\quad \textbf{D.} x $\ne$ -1/2\par\vspace{2pt}
\textbf{\color{primary} Câu 11.} Hãy chỉ ra kết quả đúng trong bốn lựa chọn A, B, C, D đối với yêu cầu: Rút gọn phân thức 24x$^2$y / (36xy$^2$) với x $\ne$ 0, y $\ne$ 0 ta được:\par
\hspace*{1.5em}\textbf{A.} 2x$^2$ / (3y$^2$)\quad\quad \textbf{B.} 3x / (2y)\quad\quad \textbf{C.} 2x / (3y)\quad\quad \textbf{D.} 2y / (3x)\par\vspace{2pt}
\textbf{\color{primary} Câu 12.} Câu hỏi ôn tập lý thuyết và kỹ năng cơ bản (Chương VI): Phân thức đối của phân thức (x - 14)/(x$^2$ + 1) là:\par
\hspace*{1.5em}\textbf{A.} (x$^2$ + 1)/(x - 14)\quad\quad \textbf{B.} (14 - x)/(x$^2$ + 1)\quad\quad \textbf{C.} (-x - 14)/(x$^2$ + 1)\quad\quad \textbf{D.} (x + 14)/(x$^2$ + 1)\par\vspace{2pt}
\textbf{\color{primary} Câu 13.} Theo nội dung sách giáo khoa Kết nối tri thức với cuộc sống (Chương VI), điều kiện xác định của phân thức đại số (2x + 1)/(x - 14) là:\par
\hspace*{1.5em}\textbf{A.} x = 14\quad\quad \textbf{B.} x $\ne$ 14\quad\quad \textbf{C.} x $\ne$ -1/2\quad\quad \textbf{D.} x $\ne$ -14\par\vspace{2pt}
\textbf{\color{primary} Câu 14.} Tìm đáp số đúng của bài toán cơ bản sau: Rút gọn phân thức 30x$^2$y / (45xy$^2$) với x $\ne$ 0, y $\ne$ 0 ta được:\par
\hspace*{1.5em}\textbf{A.} 2x / (3y)\quad\quad \textbf{B.} 2x$^2$ / (3y$^2$)\quad\quad \textbf{C.} 2y / (3x)\quad\quad \textbf{D.} 3x / (2y)\par\vspace{2pt}
\textbf{\color{primary} Câu 15.} Xét đối tượng toán học trong Chương VI: Phân thức đối của phân thức (x - 17)/(x$^2$ + 1) là. Giá trị hoặc tính chất tương ứng là:\par
\hspace*{1.5em}\textbf{A.} (x + 17)/(x$^2$ + 1)\quad\quad \textbf{B.} (x$^2$ + 1)/(x - 17)\quad\quad \textbf{C.} (17 - x)/(x$^2$ + 1)\quad\quad \textbf{D.} (-x - 17)/(x$^2$ + 1)\par\vspace{2pt}
\textbf{\color{primary} Câu 16.} Hoàn thành phép xác định cơ bản sau đây: Điều kiện xác định của phân thức đại số (2x + 1)/(x - 17) là:\par
\hspace*{1.5em}\textbf{A.} x $\ne$ -17\quad\quad \textbf{B.} x $\ne$ 17\quad\quad \textbf{C.} x $\ne$ -1/2\quad\quad \textbf{D.} x = 17\par\vspace{2pt}
\textbf{\color{primary} Câu 17.} Trong các khẳng định hoặc giá trị dưới đây, đâu là đáp án đúng cho câu hỏi: Rút gọn phân thức 36x$^2$y / (54xy$^2$) với x $\ne$ 0, y $\ne$ 0 ta được?\par
\hspace*{1.5em}\textbf{A.} 3x / (2y)\quad\quad \textbf{B.} 2y / (3x)\quad\quad \textbf{C.} 2x / (3y)\quad\quad \textbf{D.} 2x$^2$ / (3y$^2$)\par\vspace{2pt}
\textbf{\color{primary} Câu 18.} Bài tập khởi động Chương VI: Phân thức đối của phân thức (x - 20)/(x$^2$ + 1) là:\par
\hspace*{1.5em}\textbf{A.} (20 - x)/(x$^2$ + 1)\quad\quad \textbf{B.} (x + 20)/(x$^2$ + 1)\quad\quad \textbf{C.} (x$^2$ + 1)/(x - 20)\quad\quad \textbf{D.} (-x - 20)/(x$^2$ + 1)\par\vspace{2pt}
\textbf{\color{primary} Câu 19.} Quan sát dữ kiện đề bài: Điều kiện xác định của phân thức đại số (2x + 1)/(x - 20) là. Hãy chọn phương án đúng:\par
\hspace*{1.5em}\textbf{A.} x $\ne$ -20\quad\quad \textbf{B.} x $\ne$ -1/2\quad\quad \textbf{C.} x = 20\quad\quad \textbf{D.} x $\ne$ 20\par\vspace{2pt}
\textbf{\color{primary} Câu 20.} Sử dụng công thức hoặc tính chất đầu bài của Chương VI để giải quyết yêu cầu: Rút gọn phân thức 42x$^2$y / (63xy$^2$) với x $\ne$ 0, y $\ne$ 0 ta được:\par
\hspace*{1.5em}\textbf{A.} 2x$^2$ / (3y$^2$)\quad\quad \textbf{B.} 2y / (3x)\quad\quad \textbf{C.} 3x / (2y)\quad\quad \textbf{D.} 2x / (3y)\par\vspace{2pt}
\subsection*{\color{accent} Phần II. Mức độ Trung bình (Thông hiểu -- Câu 21 đến Câu 40)}
\textbf{\color{primary} Câu 21.} Rút gọn phân thức (x$^2$ - 4) / (x + 2) với x $\ne$ -2 ta được:\par
\hspace*{1.5em}\textbf{A.} x + 2\quad\quad \textbf{B.} 1\quad\quad \textbf{C.} x - 2\quad\quad \textbf{D.} x - 4\par\vspace{2pt}
\textbf{\color{primary} Câu 22.} Vận dụng quy tắc trong Chương VI (Phân thức đại số) để giải bài toán: Thực hiện phép cộng hai phân thức cùng mẫu: 2x/(x + 4) + 8/(x + 4) với x $\ne$ -4:\par
\hspace*{1.5em}\textbf{A.} 2\quad\quad \textbf{B.} 2/(x + 4)\quad\quad \textbf{C.} 2x + 8\quad\quad \textbf{D.} 2x\par\vspace{2pt}
\textbf{\color{primary} Câu 23.} Biến đổi và tính toán theo kiến thức Chương VI: Tính giá trị của phân thức A = (x$^2$ - 4)/(x - 2) tại x = 12. Kết quả thu được là:\par
\hspace*{1.5em}\textbf{A.} 140\quad\quad \textbf{B.} 14\quad\quad \textbf{C.} 12\quad\quad \textbf{D.} 10\par\vspace{2pt}
\textbf{\color{primary} Câu 24.} Xét bài toán mức độ thông hiểu sau: Rút gọn phân thức (x$^2$ - 25) / (x + 5) với x $\ne$ -5 ta được. Đáp án đúng là:\par
\hspace*{1.5em}\textbf{A.} x - 25\quad\quad \textbf{B.} x - 5\quad\quad \textbf{C.} 1\quad\quad \textbf{D.} x + 5\par\vspace{2pt}
\textbf{\color{primary} Câu 25.} Để hoàn thành phiếu học tập Chương VI, hãy xác định: Thực hiện phép cộng hai phân thức cùng mẫu: 2x/(x + 7) + 14/(x + 7) với x $\ne$ -7:\par
\hspace*{1.5em}\textbf{A.} 2/(x + 7)\quad\quad \textbf{B.} 2x\quad\quad \textbf{C.} 2x + 14\quad\quad \textbf{D.} 2\par\vspace{2pt}
\textbf{\color{primary} Câu 26.} Phân tích dữ kiện bài toán: Tính giá trị của phân thức A = (x$^2$ - 4)/(x - 2) tại x = 15. Giá trị cần tìm bằng:\par
\hspace*{1.5em}\textbf{A.} 221\quad\quad \textbf{B.} 13\quad\quad \textbf{C.} 17\quad\quad \textbf{D.} 15\par\vspace{2pt}
\textbf{\color{primary} Câu 27.} Áp dụng tính chất đã học ở Chương VI vào bài tập sau: Rút gọn phân thức (x$^2$ - 64) / (x + 8) với x $\ne$ -8 ta được:\par
\hspace*{1.5em}\textbf{A.} x - 8\quad\quad \textbf{B.} x + 8\quad\quad \textbf{C.} 1\quad\quad \textbf{D.} x - 64\par\vspace{2pt}
\textbf{\color{primary} Câu 28.} Thực hiện các bước suy luận hợp lý để giải quyết yêu cầu: Thực hiện phép cộng hai phân thức cùng mẫu: 2x/(x + 10) + 20/(x + 10) với x $\ne$ -10:\par
\hspace*{1.5em}\textbf{A.} 2x + 20\quad\quad \textbf{B.} 2/(x + 10)\quad\quad \textbf{C.} 2x\quad\quad \textbf{D.} 2\par\vspace{2pt}
\textbf{\color{primary} Câu 29.} Cho bài toán tự luyện (Phân thức đại số): Tính giá trị của phân thức A = (x$^2$ - 4)/(x - 2) tại x = 18. Phương án nào sau đây chính xác?\par
\hspace*{1.5em}\textbf{A.} 320\quad\quad \textbf{B.} 16\quad\quad \textbf{C.} 18\quad\quad \textbf{D.} 20\par\vspace{2pt}
\textbf{\color{primary} Câu 30.} Kết hợp các phép biến đổi trong Chương VI, hãy tìm kết quả của bài toán: Rút gọn phân thức (x$^2$ - 121) / (x + 11) với x $\ne$ -11 ta được:\par
\hspace*{1.5em}\textbf{A.} x - 121\quad\quad \textbf{B.} 1\quad\quad \textbf{C.} x + 11\quad\quad \textbf{D.} x - 11\par\vspace{2pt}
\textbf{\color{primary} Câu 31.} Một bài kiểm tra định kỳ nội dung Chương VI nêu câu hỏi: Thực hiện phép cộng hai phân thức cùng mẫu: 2x/(x + 13) + 26/(x + 13) với x $\ne$ -13. Hãy chọn đáp số đúng:\par
\hspace*{1.5em}\textbf{A.} 2\quad\quad \textbf{B.} 2x\quad\quad \textbf{C.} 2x + 26\quad\quad \textbf{D.} 2/(x + 13)\par\vspace{2pt}
\textbf{\color{primary} Câu 32.} Tìm lời giải đúng cho bài toán thông hiểu: Tính giá trị của phân thức A = (x$^2$ - 4)/(x - 2) tại x = 21:\par
\hspace*{1.5em}\textbf{A.} 23\quad\quad \textbf{B.} 19\quad\quad \textbf{C.} 21\quad\quad \textbf{D.} 437\par\vspace{2pt}
\textbf{\color{primary} Câu 33.} Sử dụng mối liên hệ giữa các đại lượng trong Chương VI: Rút gọn phân thức (x$^2$ - 196) / (x + 14) với x $\ne$ -14 ta được:\par
\hspace*{1.5em}\textbf{A.} x - 14\quad\quad \textbf{B.} x - 196\quad\quad \textbf{C.} 1\quad\quad \textbf{D.} x + 14\par\vspace{2pt}
\textbf{\color{primary} Câu 34.} Hãy tính toán và chọn phương án đúng đối với bài tập: Thực hiện phép cộng hai phân thức cùng mẫu: 2x/(x + 16) + 32/(x + 16) với x $\ne$ -16:\par
\hspace*{1.5em}\textbf{A.} 2\quad\quad \textbf{B.} 2x + 32\quad\quad \textbf{C.} 2/(x + 16)\quad\quad \textbf{D.} 2x\par\vspace{2pt}
\textbf{\color{primary} Câu 35.} Dựa trên các tính chất trọng tâm của Chương VI, xét bài toán: Tính giá trị của phân thức A = (x$^2$ - 4)/(x - 2) tại x = 24:\par
\hspace*{1.5em}\textbf{A.} 572\quad\quad \textbf{B.} 22\quad\quad \textbf{C.} 24\quad\quad \textbf{D.} 26\par\vspace{2pt}
\textbf{\color{primary} Câu 36.} Thực hành giải toán (Phân thức đại số): Rút gọn phân thức (x$^2$ - 289) / (x + 17) với x $\ne$ -17 ta được. Đáp án nào dưới đây thỏa mãn?\par
\hspace*{1.5em}\textbf{A.} x - 17\quad\quad \textbf{B.} x - 289\quad\quad \textbf{C.} 1\quad\quad \textbf{D.} x + 17\par\vspace{2pt}
\textbf{\color{primary} Câu 37.} Kiểm tra kỹ năng biến đổi và suy luận qua bài toán: Thực hiện phép cộng hai phân thức cùng mẫu: 2x/(x + 19) + 38/(x + 19) với x $\ne$ -19:\par
\hspace*{1.5em}\textbf{A.} 2/(x + 19)\quad\quad \textbf{B.} 2x + 38\quad\quad \textbf{C.} 2x\quad\quad \textbf{D.} 2\par\vspace{2pt}
\textbf{\color{primary} Câu 38.} Cho giả thiết thuộc phạm vi Chương VI: Tính giá trị của phân thức A = (x$^2$ - 4)/(x - 2) tại x = 27. Kết luận đúng là:\par
\hspace*{1.5em}\textbf{A.} 29\quad\quad \textbf{B.} 27\quad\quad \textbf{C.} 725\quad\quad \textbf{D.} 25\par\vspace{2pt}
\textbf{\color{primary} Câu 39.} Rút gọn hoặc tính giá trị theo yêu cầu sau của Chương VI: Rút gọn phân thức (x$^2$ - 400) / (x + 20) với x $\ne$ -20 ta được:\par
\hspace*{1.5em}\textbf{A.} 1\quad\quad \textbf{B.} x - 20\quad\quad \textbf{C.} x + 20\quad\quad \textbf{D.} x - 400\par\vspace{2pt}
\textbf{\color{primary} Câu 40.} Xác định phương án đúng sau khi thực hiện phép tính / suy luận: Thực hiện phép cộng hai phân thức cùng mẫu: 2x/(x + 22) + 44/(x + 22) với x $\ne$ -22:\par
\hspace*{1.5em}\textbf{A.} 2x\quad\quad \textbf{B.} 2\quad\quad \textbf{C.} 2x + 44\quad\quad \textbf{D.} 2/(x + 22)\par\vspace{2pt}
\subsection*{\color{accent} Phần III. Mức độ Khó (Vận dụng \& Vận dụng cao -- Câu 41 đến Câu 50)}
\textbf{\color{primary} Câu 41.} Tìm tất cả các giá trị nguyên của x để phân thức P = (2x + 7)/(x + 2) nhận giá trị nguyên:\par
\hspace*{1.5em}\textbf{A.} x $\in$ \{1; 3; -1; -3\}\quad\quad \textbf{B.} x $\in$ \{-3; -1\}\quad\quad \textbf{C.} x = -2\quad\quad \textbf{D.} x $\in$ \{-5; -3; -1; 1\}\par\vspace{2pt}
\textbf{\color{primary} Câu 42.} Bài toán vận dụng (Chương VI – Phân thức đại số): Cho x + 1/x = 4 (với x $\ne$ 0). Giá trị của biểu thức M = x$^2$ + 1/x$^2$ bằng. Kết quả đúng là:\par
\hspace*{1.5em}\textbf{A.} 14\quad\quad \textbf{B.} 16\quad\quad \textbf{C.} 18\quad\quad \textbf{D.} 8\par\vspace{2pt}
\textbf{\color{primary} Câu 43.} Vận dụng kiến thức nâng cao trong phạm vi Chương VI để giải bài toán: Tìm tất cả các giá trị nguyên của x để phân thức P = (2x + 11)/(x + 4) nhận giá trị nguyên:\par
\hspace*{1.5em}\textbf{A.} x $\in$ \{1; 3; -1; -3\}\quad\quad \textbf{B.} x = -4\quad\quad \textbf{C.} x $\in$ \{-7; -5; -3; -1\}\quad\quad \textbf{D.} x $\in$ \{-5; -3\}\par\vspace{2pt}
\textbf{\color{primary} Câu 44.} Phân tích và giải bài toán tổng hợp sau của Chương VI: Cho x + 1/x = 6 (với x $\ne$ 0). Giá trị của biểu thức M = x$^2$ + 1/x$^2$ bằng:\par
\hspace*{1.5em}\textbf{A.} 38\quad\quad \textbf{B.} 36\quad\quad \textbf{C.} 12\quad\quad \textbf{D.} 34\par\vspace{2pt}
\textbf{\color{primary} Câu 45.} Thử thách tư duy trong phạm vi kiến thức Chương VI: Tìm tất cả các giá trị nguyên của x để phân thức P = (2x + 15)/(x + 6) nhận giá trị nguyên. Đáp số chính xác là:\par
\hspace*{1.5em}\textbf{A.} x $\in$ \{-7; -5\}\quad\quad \textbf{B.} x $\in$ \{1; 3; -1; -3\}\quad\quad \textbf{C.} x $\in$ \{-9; -7; -5; -3\}\quad\quad \textbf{D.} x = -6\par\vspace{2pt}
\textbf{\color{primary} Câu 46.} Sử dụng phương pháp giải toán của Chương VI để tìm đáp án cho bài toán: Cho x + 1/x = 8 (với x $\ne$ 0). Giá trị của biểu thức M = x$^2$ + 1/x$^2$ bằng:\par
\hspace*{1.5em}\textbf{A.} 64\quad\quad \textbf{B.} 16\quad\quad \textbf{C.} 66\quad\quad \textbf{D.} 62\par\vspace{2pt}
\textbf{\color{primary} Câu 47.} Xét bài toán vận dụng thực tiễn / nâng cao thuộc Chương VI: Tìm tất cả các giá trị nguyên của x để phân thức P = (2x + 19)/(x + 8) nhận giá trị nguyên:\par
\hspace*{1.5em}\textbf{A.} x $\in$ \{1; 3; -1; -3\}\quad\quad \textbf{B.} x = -8\quad\quad \textbf{C.} x $\in$ \{-11; -9; -7; -5\}\quad\quad \textbf{D.} x $\in$ \{-9; -7\}\par\vspace{2pt}
\textbf{\color{primary} Câu 48.} Tìm kết quả cuối cùng của bài toán phân loại học sinh khá giỏi (Chương VI): Cho x + 1/x = 10 (với x $\ne$ 0). Giá trị của biểu thức M = x$^2$ + 1/x$^2$ bằng:\par
\hspace*{1.5em}\textbf{A.} 102\quad\quad \textbf{B.} 100\quad\quad \textbf{C.} 98\quad\quad \textbf{D.} 20\par\vspace{2pt}
\textbf{\color{primary} Câu 49.} Kết hợp các tính chất đã học trong Chương VI, hãy giải quyết bài toán: Tìm tất cả các giá trị nguyên của x để phân thức P = (2x + 23)/(x + 10) nhận giá trị nguyên:\par
\hspace*{1.5em}\textbf{A.} x $\in$ \{-13; -11; -9; -7\}\quad\quad \textbf{B.} x = -10\quad\quad \textbf{C.} x $\in$ \{1; 3; -1; -3\}\quad\quad \textbf{D.} x $\in$ \{-11; -9\}\par\vspace{2pt}
\textbf{\color{primary} Câu 50.} Câu hỏi vận dụng cao cuối Chương VI: Cho x + 1/x = 12 (với x $\ne$ 0). Giá trị của biểu thức M = x$^2$ + 1/x$^2$ bằng. Phương án đúng là:\par
\hspace*{1.5em}\textbf{A.} 24\quad\quad \textbf{B.} 144\quad\quad \textbf{C.} 146\quad\quad \textbf{D.} 142\par\vspace{2pt}
\subsection*{\color{primary} Bảng đáp án Chương VI: Phân thức đại số}
\begin{center}
\renewcommand{\arraystretch}{1.25}
\begin{tabularx}{\textwidth}{*{10}{>{\centering\arraybackslash}X}}
\toprule
\textbf{\color{primary} Cột 1} & \textbf{\color{primary} Cột 2} & \textbf{\color{primary} Cột 3} & \textbf{\color{primary} Cột 4} & \textbf{\color{primary} Cột 5} & \textbf{\color{primary} Cột 6} & \textbf{\color{primary} Cột 7} & \textbf{\color{primary} Cột 8} & \textbf{\color{primary} Cột 9} & \textbf{\color{primary} Cột 10} \\
\midrule
\textbf{1.C} & \textbf{2.B} & \textbf{3.C} & \textbf{4.C} & \textbf{5.C} & \textbf{6.B} & \textbf{7.B} & \textbf{8.D} & \textbf{9.A} & \textbf{10.C} \\
\rowcolor{rowtint} \textbf{11.C} & \textbf{12.B} & \textbf{13.B} & \textbf{14.A} & \textbf{15.C} & \textbf{16.B} & \textbf{17.C} & \textbf{18.A} & \textbf{19.D} & \textbf{20.D} \\
\textbf{21.C} & \textbf{22.A} & \textbf{23.B} & \textbf{24.B} & \textbf{25.D} & \textbf{26.C} & \textbf{27.A} & \textbf{28.D} & \textbf{29.D} & \textbf{30.D} \\
\rowcolor{rowtint} \textbf{31.A} & \textbf{32.A} & \textbf{33.A} & \textbf{34.A} & \textbf{35.D} & \textbf{36.A} & \textbf{37.D} & \textbf{38.A} & \textbf{39.B} & \textbf{40.B} \\
\textbf{41.D} & \textbf{42.A} & \textbf{43.C} & \textbf{44.D} & \textbf{45.C} & \textbf{46.D} & \textbf{47.C} & \textbf{48.C} & \textbf{49.A} & \textbf{50.D} \\
\bottomrule
\end{tabularx}
\end{center}\vspace{10pt}
\section*{\color{primary} CHƯƠNG VII: PHƯƠNG TRÌNH BẬC NHẤT VÀ HÀM SỐ BẬC NHẤT}
\textit{Trọng tâm kiến thức SGK: Phương trình bậc nhất một ẩn, giải bài toán bằng cách lập phương trình, khái niệm hàm số và đồ thị, hàm số bậc nhất y = ax + b (a $\ne$ 0) và hệ số góc của đường thẳng.}\par\vspace{4pt}
\subsection*{\color{accent} Phần I. Mức độ Dễ (Nhận biết -- Câu 1 đến Câu 20)}
\textbf{\color{primary} Câu 1.} Nghiệm của phương trình bậc nhất 3x - 6 = 0 là:\par
\hspace*{1.5em}\textbf{A.} x = 3\quad\quad \textbf{B.} x = -2\quad\quad \textbf{C.} x = 6\quad\quad \textbf{D.} x = 2\par\vspace{2pt}
\textbf{\color{primary} Câu 2.} Thực hiện yêu cầu nhận biết sau: Hệ số góc của đường thẳng y = -3x + 7 là.\par
\hspace*{1.5em}\textbf{A.} a = 7\quad\quad \textbf{B.} a = -7\quad\quad \textbf{C.} a = 3\quad\quad \textbf{D.} a = -3\par\vspace{2pt}
\textbf{\color{primary} Câu 3.} Trong bài học thuộc Chương VII, xét nội dung: Cho hàm số f(x) = 2x - 5. Giá trị của f(4) bằng. Đáp án đúng là:\par
\hspace*{1.5em}\textbf{A.} 13\quad\quad \textbf{B.} 3\quad\quad \textbf{C.} -3\quad\quad \textbf{D.} 6\par\vspace{2pt}
\textbf{\color{primary} Câu 4.} Dựa vào kiến thức chuẩn SGK của Chương VII (Phương trình bậc nhất và hàm số bậc nhất), nghiệm của phương trình bậc nhất 3x - 15 = 0 là:\par
\hspace*{1.5em}\textbf{A.} x = 5\quad\quad \textbf{B.} x = 3\quad\quad \textbf{C.} x = -5\quad\quad \textbf{D.} x = 15\par\vspace{2pt}
\textbf{\color{primary} Câu 5.} Chọn kết quả đúng nhất khi xét: Hệ số góc của đường thẳng y = -6x + 13 là:\par
\hspace*{1.5em}\textbf{A.} a = 6\quad\quad \textbf{B.} a = 13\quad\quad \textbf{C.} a = -6\quad\quad \textbf{D.} a = -13\par\vspace{2pt}
\textbf{\color{primary} Câu 6.} Khi làm bài tập cơ bản về phương trình bậc nhất và hàm số bậc nhất, một học sinh gặp câu hỏi: Cho hàm số f(x) = 2x - 8. Giá trị của f(4) bằng. Kết quả nào dưới đây đúng?\par
\hspace*{1.5em}\textbf{A.} 0\quad\quad \textbf{B.} 16\quad\quad \textbf{C.} 3\quad\quad \textbf{D.} -0\par\vspace{2pt}
\textbf{\color{primary} Câu 7.} Xác định giá trị hoặc khẳng định đúng cho bài toán sau: Nghiệm của phương trình bậc nhất 3x - 24 = 0 là:\par
\hspace*{1.5em}\textbf{A.} x = 24\quad\quad \textbf{B.} x = 3\quad\quad \textbf{C.} x = 8\quad\quad \textbf{D.} x = -8\par\vspace{2pt}
\textbf{\color{primary} Câu 8.} Áp dụng trực tiếp định nghĩa và quy tắc trong Chương VII: Hệ số góc của đường thẳng y = -9x + 19 là:\par
\hspace*{1.5em}\textbf{A.} a = 19\quad\quad \textbf{B.} a = -19\quad\quad \textbf{C.} a = -9\quad\quad \textbf{D.} a = 9\par\vspace{2pt}
\textbf{\color{primary} Câu 9.} Kiểm tra nhanh kiến thức nền tảng: Cho hàm số f(x) = 2x - 11. Giá trị của f(4) bằng:\par
\hspace*{1.5em}\textbf{A.} --3\quad\quad \textbf{B.} 0\quad\quad \textbf{C.} 19\quad\quad \textbf{D.} -3\par\vspace{2pt}
\textbf{\color{primary} Câu 10.} Cho dữ kiện nhận biết sau đây: Nghiệm của phương trình bậc nhất 3x - 33 = 0 là. Phương án chính xác là:\par
\hspace*{1.5em}\textbf{A.} x = 11\quad\quad \textbf{B.} x = -11\quad\quad \textbf{C.} x = 33\quad\quad \textbf{D.} x = 3\par\vspace{2pt}
\textbf{\color{primary} Câu 11.} Hãy chỉ ra kết quả đúng trong bốn lựa chọn A, B, C, D đối với yêu cầu: Hệ số góc của đường thẳng y = -12x + 25 là:\par
\hspace*{1.5em}\textbf{A.} a = -25\quad\quad \textbf{B.} a = 25\quad\quad \textbf{C.} a = -12\quad\quad \textbf{D.} a = 12\par\vspace{2pt}
\textbf{\color{primary} Câu 12.} Câu hỏi ôn tập lý thuyết và kỹ năng cơ bản (Chương VII): Cho hàm số f(x) = 2x - 14. Giá trị của f(4) bằng:\par
\hspace*{1.5em}\textbf{A.} 22\quad\quad \textbf{B.} -6\quad\quad \textbf{C.} --6\quad\quad \textbf{D.} -3\par\vspace{2pt}
\textbf{\color{primary} Câu 13.} Theo nội dung sách giáo khoa Kết nối tri thức với cuộc sống (Chương VII), nghiệm của phương trình bậc nhất 3x - 42 = 0 là:\par
\hspace*{1.5em}\textbf{A.} x = 14\quad\quad \textbf{B.} x = 3\quad\quad \textbf{C.} x = -14\quad\quad \textbf{D.} x = 42\par\vspace{2pt}
\textbf{\color{primary} Câu 14.} Tìm đáp số đúng của bài toán cơ bản sau: Hệ số góc của đường thẳng y = -15x + 31 là:\par
\hspace*{1.5em}\textbf{A.} a = 31\quad\quad \textbf{B.} a = -31\quad\quad \textbf{C.} a = -15\quad\quad \textbf{D.} a = 15\par\vspace{2pt}
\textbf{\color{primary} Câu 15.} Xét đối tượng toán học trong Chương VII: Cho hàm số f(x) = 2x - 17. Giá trị của f(4) bằng. Giá trị hoặc tính chất tương ứng là:\par
\hspace*{1.5em}\textbf{A.} -9\quad\quad \textbf{B.} 25\quad\quad \textbf{C.} -6\quad\quad \textbf{D.} --9\par\vspace{2pt}
\textbf{\color{primary} Câu 16.} Hoàn thành phép xác định cơ bản sau đây: Nghiệm của phương trình bậc nhất 3x - 51 = 0 là:\par
\hspace*{1.5em}\textbf{A.} x = -17\quad\quad \textbf{B.} x = 51\quad\quad \textbf{C.} x = 17\quad\quad \textbf{D.} x = 3\par\vspace{2pt}
\textbf{\color{primary} Câu 17.} Trong các khẳng định hoặc giá trị dưới đây, đâu là đáp án đúng cho câu hỏi: Hệ số góc của đường thẳng y = -18x + 37 là?\par
\hspace*{1.5em}\textbf{A.} a = 37\quad\quad \textbf{B.} a = -37\quad\quad \textbf{C.} a = 18\quad\quad \textbf{D.} a = -18\par\vspace{2pt}
\textbf{\color{primary} Câu 18.} Bài tập khởi động Chương VII: Cho hàm số f(x) = 2x - 20. Giá trị của f(4) bằng:\par
\hspace*{1.5em}\textbf{A.} 28\quad\quad \textbf{B.} -12\quad\quad \textbf{C.} -9\quad\quad \textbf{D.} --12\par\vspace{2pt}
\textbf{\color{primary} Câu 19.} Quan sát dữ kiện đề bài: Nghiệm của phương trình bậc nhất 3x - 60 = 0 là. Hãy chọn phương án đúng:\par
\hspace*{1.5em}\textbf{A.} x = -20\quad\quad \textbf{B.} x = 20\quad\quad \textbf{C.} x = 60\quad\quad \textbf{D.} x = 3\par\vspace{2pt}
\textbf{\color{primary} Câu 20.} Sử dụng công thức hoặc tính chất đầu bài của Chương VII để giải quyết yêu cầu: Hệ số góc của đường thẳng y = -21x + 43 là:\par
\hspace*{1.5em}\textbf{A.} a = -43\quad\quad \textbf{B.} a = 43\quad\quad \textbf{C.} a = 21\quad\quad \textbf{D.} a = -21\par\vspace{2pt}
\subsection*{\color{accent} Phần II. Mức độ Trung bình (Thông hiểu -- Câu 21 đến Câu 40)}
\textbf{\color{primary} Câu 21.} Tìm giá trị của tham số m để hai đường thẳng (d$_1$): y = (m - 1)x + 3 và (d$_2$): y = 2x - 5 song song với nhau:\par
\hspace*{1.5em}\textbf{A.} m = 2\quad\quad \textbf{B.} m = 3\quad\quad \textbf{C.} m = -3\quad\quad \textbf{D.} m = -1\par\vspace{2pt}
\textbf{\color{primary} Câu 22.} Vận dụng quy tắc trong Chương VII (Phương trình bậc nhất và hàm số bậc nhất) để giải bài toán: Giải phương trình: 5(x - 2) = 3x + 6.:\par
\hspace*{1.5em}\textbf{A.} x = 5\quad\quad \textbf{B.} x = 8\quad\quad \textbf{C.} x = -8\quad\quad \textbf{D.} x = 6\par\vspace{2pt}
\textbf{\color{primary} Câu 23.} Biến đổi và tính toán theo kiến thức Chương VII: Đường thẳng y = 2x + b đi qua điểm A(1; 8) khi hệ số b bằng. Kết quả thu được là:\par
\hspace*{1.5em}\textbf{A.} b = 8\quad\quad \textbf{B.} b = 6\quad\quad \textbf{C.} b = -6\quad\quad \textbf{D.} b = 10\par\vspace{2pt}
\textbf{\color{primary} Câu 24.} Xét bài toán mức độ thông hiểu sau: Tìm giá trị của tham số m để hai đường thẳng (d$_1$): y = (m - 4)x + 3 và (d$_2$): y = 2x - 5 song song với nhau. Đáp án đúng là:\par
\hspace*{1.5em}\textbf{A.} m = 6 (1)\quad\quad \textbf{B.} m = 2\quad\quad \textbf{C.} m = 6\quad\quad \textbf{D.} m = -6\par\vspace{2pt}
\textbf{\color{primary} Câu 25.} Để hoàn thành phiếu học tập Chương VII, hãy xác định: Giải phương trình: 5(x - 5) = 3x + 9.:\par
\hspace*{1.5em}\textbf{A.} x = 17\quad\quad \textbf{B.} x = 9\quad\quad \textbf{C.} x = -17\quad\quad \textbf{D.} x = 11\par\vspace{2pt}
\textbf{\color{primary} Câu 26.} Phân tích dữ kiện bài toán: Đường thẳng y = 2x + b đi qua điểm A(1; 11) khi hệ số b bằng. Giá trị cần tìm bằng:\par
\hspace*{1.5em}\textbf{A.} b = -9\quad\quad \textbf{B.} b = 9\quad\quad \textbf{C.} b = 13\quad\quad \textbf{D.} b = 11\par\vspace{2pt}
\textbf{\color{primary} Câu 27.} Áp dụng tính chất đã học ở Chương VII vào bài tập sau: Tìm giá trị của tham số m để hai đường thẳng (d$_1$): y = (m - 7)x + 3 và (d$_2$): y = 2x - 5 song song với nhau:\par
\hspace*{1.5em}\textbf{A.} m = -9\quad\quad \textbf{B.} m = 5\quad\quad \textbf{C.} m = 2\quad\quad \textbf{D.} m = 9\par\vspace{2pt}
\textbf{\color{primary} Câu 28.} Thực hiện các bước suy luận hợp lý để giải quyết yêu cầu: Giải phương trình: 5(x - 8) = 3x + 12.:\par
\hspace*{1.5em}\textbf{A.} x = -26\quad\quad \textbf{B.} x = 26\quad\quad \textbf{C.} x = 17\quad\quad \textbf{D.} x = 12\par\vspace{2pt}
\textbf{\color{primary} Câu 29.} Cho bài toán tự luyện (Phương trình bậc nhất và hàm số bậc nhất): Đường thẳng y = 2x + b đi qua điểm A(1; 14) khi hệ số b bằng. Phương án nào sau đây chính xác?\par
\hspace*{1.5em}\textbf{A.} b = 16\quad\quad \textbf{B.} b = -12\quad\quad \textbf{C.} b = 14\quad\quad \textbf{D.} b = 12\par\vspace{2pt}
\textbf{\color{primary} Câu 30.} Kết hợp các phép biến đổi trong Chương VII, hãy tìm kết quả của bài toán: Tìm giá trị của tham số m để hai đường thẳng (d$_1$): y = (m - 10)x + 3 và (d$_2$): y = 2x - 5 song song với nhau:\par
\hspace*{1.5em}\textbf{A.} m = 8\quad\quad \textbf{B.} m = -12\quad\quad \textbf{C.} m = 2\quad\quad \textbf{D.} m = 12\par\vspace{2pt}
\textbf{\color{primary} Câu 31.} Một bài kiểm tra định kỳ nội dung Chương VII nêu câu hỏi: Giải phương trình: 5(x - 11) = 3x + 15.. Hãy chọn đáp số đúng:\par
\hspace*{1.5em}\textbf{A.} x = 23\quad\quad \textbf{B.} x = -35\quad\quad \textbf{C.} x = 35\quad\quad \textbf{D.} x = 15\par\vspace{2pt}
\textbf{\color{primary} Câu 32.} Tìm lời giải đúng cho bài toán thông hiểu: Đường thẳng y = 2x + b đi qua điểm A(1; 17) khi hệ số b bằng:\par
\hspace*{1.5em}\textbf{A.} b = -15\quad\quad \textbf{B.} b = 15\quad\quad \textbf{C.} b = 19\quad\quad \textbf{D.} b = 17\par\vspace{2pt}
\textbf{\color{primary} Câu 33.} Sử dụng mối liên hệ giữa các đại lượng trong Chương VII: Tìm giá trị của tham số m để hai đường thẳng (d$_1$): y = (m - 13)x + 3 và (d$_2$): y = 2x - 5 song song với nhau:\par
\hspace*{1.5em}\textbf{A.} m = -15\quad\quad \textbf{B.} m = 15\quad\quad \textbf{C.} m = 2\quad\quad \textbf{D.} m = 11\par\vspace{2pt}
\textbf{\color{primary} Câu 34.} Hãy tính toán và chọn phương án đúng đối với bài tập: Giải phương trình: 5(x - 14) = 3x + 18.:\par
\hspace*{1.5em}\textbf{A.} x = 18\quad\quad \textbf{B.} x = -44\quad\quad \textbf{C.} x = 44\quad\quad \textbf{D.} x = 29\par\vspace{2pt}
\textbf{\color{primary} Câu 35.} Dựa trên các tính chất trọng tâm của Chương VII, xét bài toán: Đường thẳng y = 2x + b đi qua điểm A(1; 20) khi hệ số b bằng:\par
\hspace*{1.5em}\textbf{A.} b = 18\quad\quad \textbf{B.} b = 20\quad\quad \textbf{C.} b = -18\quad\quad \textbf{D.} b = 22\par\vspace{2pt}
\textbf{\color{primary} Câu 36.} Thực hành giải toán (Phương trình bậc nhất và hàm số bậc nhất): Tìm giá trị của tham số m để hai đường thẳng (d$_1$): y = (m - 16)x + 3 và (d$_2$): y = 2x - 5 song song với nhau. Đáp án nào dưới đây thỏa mãn?\par
\hspace*{1.5em}\textbf{A.} m = -18\quad\quad \textbf{B.} m = 2\quad\quad \textbf{C.} m = 18\quad\quad \textbf{D.} m = 14\par\vspace{2pt}
\textbf{\color{primary} Câu 37.} Kiểm tra kỹ năng biến đổi và suy luận qua bài toán: Giải phương trình: 5(x - 17) = 3x + 21.:\par
\hspace*{1.5em}\textbf{A.} x = 53\quad\quad \textbf{B.} x = 21\quad\quad \textbf{C.} x = -53\quad\quad \textbf{D.} x = 35\par\vspace{2pt}
\textbf{\color{primary} Câu 38.} Cho giả thiết thuộc phạm vi Chương VII: Đường thẳng y = 2x + b đi qua điểm A(1; 23) khi hệ số b bằng. Kết luận đúng là:\par
\hspace*{1.5em}\textbf{A.} b = 25\quad\quad \textbf{B.} b = 23\quad\quad \textbf{C.} b = -21\quad\quad \textbf{D.} b = 21\par\vspace{2pt}
\textbf{\color{primary} Câu 39.} Rút gọn hoặc tính giá trị theo yêu cầu sau của Chương VII: Tìm giá trị của tham số m để hai đường thẳng (d$_1$): y = (m - 19)x + 3 và (d$_2$): y = 2x - 5 song song với nhau:\par
\hspace*{1.5em}\textbf{A.} m = 17\quad\quad \textbf{B.} m = -21\quad\quad \textbf{C.} m = 21\quad\quad \textbf{D.} m = 2\par\vspace{2pt}
\textbf{\color{primary} Câu 40.} Xác định phương án đúng sau khi thực hiện phép tính / suy luận: Giải phương trình: 5(x - 20) = 3x + 24.:\par
\hspace*{1.5em}\textbf{A.} x = 24\quad\quad \textbf{B.} x = -62\quad\quad \textbf{C.} x = 62\quad\quad \textbf{D.} x = 41\par\vspace{2pt}
\subsection*{\color{accent} Phần III. Mức độ Khó (Vận dụng \& Vận dụng cao -- Câu 41 đến Câu 50)}
\textbf{\color{primary} Câu 41.} Một ô tô đi từ A đến B với vận tốc 50 km/h rồi quay về từ B đến A với vận tốc 40 km/h. Tổng thời gian cả đi và về là 4.5 giờ. Độ dài quãng đường AB là:\par
\hspace*{1.5em}\textbf{A.} 80 km\quad\quad \textbf{B.} 90 km\quad\quad \textbf{C.} 100 km\quad\quad \textbf{D.} 120 km\par\vspace{2pt}
\textbf{\color{primary} Câu 42.} Bài toán vận dụng (Chương VII – Phương trình bậc nhất và hàm số bậc nhất): Tìm tọa độ giao điểm của hai đường thẳng d$_1$: y = 3x - 3 và d$_2$: y = x + 5. Kết quả đúng là:\par
\hspace*{1.5em}\textbf{A.} (9; 4)\quad\quad \textbf{B.} (-4; 1)\quad\quad \textbf{C.} (4; 9)\quad\quad \textbf{D.} (2; 7)\par\vspace{2pt}
\textbf{\color{primary} Câu 43.} Vận dụng kiến thức nâng cao trong phạm vi Chương VII để giải bài toán: Một ô tô đi từ A đến B với vận tốc 50 km/h rồi quay về từ B đến A với vận tốc 40 km/h. Tổng thời gian cả đi và về là 4.5 giờ. Độ dài quãng đường AB là:\par
\hspace*{1.5em}\textbf{A.} 80 km\quad\quad \textbf{B.} 100 km\quad\quad \textbf{C.} 120 km\quad\quad \textbf{D.} 90 km\par\vspace{2pt}
\textbf{\color{primary} Câu 44.} Phân tích và giải bài toán tổng hợp sau của Chương VII: Tìm tọa độ giao điểm của hai đường thẳng d$_1$: y = 3x - 5 và d$_2$: y = x + 7:\par
\hspace*{1.5em}\textbf{A.} (4; 11)\quad\quad \textbf{B.} (13; 6)\quad\quad \textbf{C.} (-6; 1)\quad\quad \textbf{D.} (6; 13)\par\vspace{2pt}
\textbf{\color{primary} Câu 45.} Thử thách tư duy trong phạm vi kiến thức Chương VII: Một ô tô đi từ A đến B với vận tốc 50 km/h rồi quay về từ B đến A với vận tốc 40 km/h. Tổng thời gian cả đi và về là 4.5 giờ. Độ dài quãng đường AB là. Đáp số chính xác là:\par
\hspace*{1.5em}\textbf{A.} 100 km\quad\quad \textbf{B.} 90 km\quad\quad \textbf{C.} 120 km\quad\quad \textbf{D.} 80 km\par\vspace{2pt}
\textbf{\color{primary} Câu 46.} Sử dụng phương pháp giải toán của Chương VII để tìm đáp án cho bài toán: Tìm tọa độ giao điểm của hai đường thẳng d$_1$: y = 3x - 7 và d$_2$: y = x + 9:\par
\hspace*{1.5em}\textbf{A.} (17; 8)\quad\quad \textbf{B.} (-8; 1)\quad\quad \textbf{C.} (8; 17)\quad\quad \textbf{D.} (6; 15)\par\vspace{2pt}
\textbf{\color{primary} Câu 47.} Xét bài toán vận dụng thực tiễn / nâng cao thuộc Chương VII: Một ô tô đi từ A đến B với vận tốc 50 km/h rồi quay về từ B đến A với vận tốc 40 km/h. Tổng thời gian cả đi và về là 4.5 giờ. Độ dài quãng đường AB là:\par
\hspace*{1.5em}\textbf{A.} 120 km\quad\quad \textbf{B.} 100 km\quad\quad \textbf{C.} 90 km\quad\quad \textbf{D.} 80 km\par\vspace{2pt}
\textbf{\color{primary} Câu 48.} Tìm kết quả cuối cùng của bài toán phân loại học sinh khá giỏi (Chương VII): Tìm tọa độ giao điểm của hai đường thẳng d$_1$: y = 3x - 9 và d$_2$: y = x + 11:\par
\hspace*{1.5em}\textbf{A.} (21; 10)\quad\quad \textbf{B.} (8; 19)\quad\quad \textbf{C.} (-10; 1)\quad\quad \textbf{D.} (10; 21)\par\vspace{2pt}
\textbf{\color{primary} Câu 49.} Kết hợp các tính chất đã học trong Chương VII, hãy giải quyết bài toán: Một ô tô đi từ A đến B với vận tốc 50 km/h rồi quay về từ B đến A với vận tốc 40 km/h. Tổng thời gian cả đi và về là 4.5 giờ. Độ dài quãng đường AB là:\par
\hspace*{1.5em}\textbf{A.} 90 km\quad\quad \textbf{B.} 100 km\quad\quad \textbf{C.} 120 km\quad\quad \textbf{D.} 80 km\par\vspace{2pt}
\textbf{\color{primary} Câu 50.} Câu hỏi vận dụng cao cuối Chương VII: Tìm tọa độ giao điểm của hai đường thẳng d$_1$: y = 3x - 11 và d$_2$: y = x + 13. Phương án đúng là:\par
\hspace*{1.5em}\textbf{A.} (25; 12)\quad\quad \textbf{B.} (10; 23)\quad\quad \textbf{C.} (-12; 1)\quad\quad \textbf{D.} (12; 25)\par\vspace{2pt}
\subsection*{\color{primary} Bảng đáp án Chương VII: Phương trình bậc nhất và hàm số bậc nhất}
\begin{center}
\renewcommand{\arraystretch}{1.25}
\begin{tabularx}{\textwidth}{*{10}{>{\centering\arraybackslash}X}}
\toprule
\textbf{\color{primary} Cột 1} & \textbf{\color{primary} Cột 2} & \textbf{\color{primary} Cột 3} & \textbf{\color{primary} Cột 4} & \textbf{\color{primary} Cột 5} & \textbf{\color{primary} Cột 6} & \textbf{\color{primary} Cột 7} & \textbf{\color{primary} Cột 8} & \textbf{\color{primary} Cột 9} & \textbf{\color{primary} Cột 10} \\
\midrule
\textbf{1.D} & \textbf{2.D} & \textbf{3.B} & \textbf{4.A} & \textbf{5.C} & \textbf{6.A} & \textbf{7.C} & \textbf{8.C} & \textbf{9.D} & \textbf{10.A} \\
\rowcolor{rowtint} \textbf{11.C} & \textbf{12.B} & \textbf{13.A} & \textbf{14.C} & \textbf{15.A} & \textbf{16.C} & \textbf{17.D} & \textbf{18.B} & \textbf{19.B} & \textbf{20.D} \\
\textbf{21.B} & \textbf{22.B} & \textbf{23.B} & \textbf{24.C} & \textbf{25.A} & \textbf{26.B} & \textbf{27.D} & \textbf{28.B} & \textbf{29.D} & \textbf{30.D} \\
\rowcolor{rowtint} \textbf{31.C} & \textbf{32.B} & \textbf{33.B} & \textbf{34.C} & \textbf{35.A} & \textbf{36.C} & \textbf{37.A} & \textbf{38.D} & \textbf{39.C} & \textbf{40.C} \\
\textbf{41.C} & \textbf{42.C} & \textbf{43.B} & \textbf{44.D} & \textbf{45.A} & \textbf{46.C} & \textbf{47.B} & \textbf{48.D} & \textbf{49.B} & \textbf{50.D} \\
\bottomrule
\end{tabularx}
\end{center}\vspace{10pt}
\section*{\color{primary} CHƯƠNG VIII: MỞ ĐẦU VỀ TÍNH XÁC SUẤT CỦA BIẾN CỐ}
\textit{Trọng tâm kiến thức SGK: Kết quả có thể và kết quả thuận lợi, cách tính xác suất của biến cố bằng tỉ số, mối liên hệ giữa xác suất thực nghiệm với xác suất và ứng dụng.}\par\vspace{4pt}
\subsection*{\color{accent} Phần I. Mức độ Dễ (Nhận biết -- Câu 1 đến Câu 20)}
\textbf{\color{primary} Câu 1.} Một hộp có 10 quả bóng được đánh số từ 1 đến 10. Lấy ngẫu nhiên 1 quả. Số kết quả thuận lợi cho biến cố 'Lấy được quả bóng ghi số nhỏ hơn 5' là:\par
\hspace*{1.5em}\textbf{A.} 4 kết quả\quad\quad \textbf{B.} 10 kết quả\quad\quad \textbf{C.} 5 kết quả\quad\quad \textbf{D.} 3 kết quả\par\vspace{2pt}
\textbf{\color{primary} Câu 2.} Thực hiện yêu cầu nhận biết sau: Nếu một biến cố E có k kết quả thuận lợi trong tổng số n kết quả đồng khả năng thì xác suất P(E) được tính bằng công thức.\par
\hspace*{1.5em}\textbf{A.} P(E) = n - k\quad\quad \textbf{B.} P(E) = k / n\quad\quad \textbf{C.} P(E) = k $\cdot$ n\quad\quad \textbf{D.} P(E) = n / k\par\vspace{2pt}
\textbf{\color{primary} Câu 3.} Trong bài học thuộc Chương VIII, xét nội dung: Khi số phép thử n ngày càng lớn thì xác suất thực nghiệm của một biến cố sẽ như thế nào. Đáp án đúng là:\par
\begin{itemize}\setlength{\itemsep}{1pt}\setlength{\parskip}{0pt}
  \item[] \textbf{A.} Luôn bằng 1
  \item[] \textbf{B.} Ngày càng gần với xác suất lý thuyết của biến cố đó
  \item[] \textbf{C.} Luôn bằng 0
  \item[] \textbf{D.} Tăng lên vô hạn
\end{itemize}\vspace{2pt}
\textbf{\color{primary} Câu 4.} Dựa vào kiến thức chuẩn SGK của Chương VIII (Mở đầu về tính xác suất của biến cố), một hộp có 16 quả bóng được đánh số từ 1 đến 16. Lấy ngẫu nhiên 1 quả. Số kết quả thuận lợi cho biến cố 'Lấy được quả bóng ghi số nhỏ hơn 5' là:\par
\hspace*{1.5em}\textbf{A.} 3 kết quả\quad\quad \textbf{B.} 16 kết quả\quad\quad \textbf{C.} 5 kết quả\quad\quad \textbf{D.} 4 kết quả\par\vspace{2pt}
\textbf{\color{primary} Câu 5.} Chọn kết quả đúng nhất khi xét: Nếu một biến cố E có k kết quả thuận lợi trong tổng số n kết quả đồng khả năng thì xác suất P(E) được tính bằng công thức:\par
\hspace*{1.5em}\textbf{A.} P(E) = n - k\quad\quad \textbf{B.} P(E) = k $\cdot$ n\quad\quad \textbf{C.} P(E) = n / k\quad\quad \textbf{D.} P(E) = k / n\par\vspace{2pt}
\textbf{\color{primary} Câu 6.} Khi làm bài tập cơ bản về mở đầu về tính xác suất của biến cố, một học sinh gặp câu hỏi: Khi số phép thử n ngày càng lớn thì xác suất thực nghiệm của một biến cố sẽ như thế nào. Kết quả nào dưới đây đúng?\par
\begin{itemize}\setlength{\itemsep}{1pt}\setlength{\parskip}{0pt}
  \item[] \textbf{A.} Ngày càng gần với xác suất lý thuyết của biến cố đó
  \item[] \textbf{B.} Tăng lên vô hạn
  \item[] \textbf{C.} Luôn bằng 1
  \item[] \textbf{D.} Luôn bằng 0
\end{itemize}\vspace{2pt}
\textbf{\color{primary} Câu 7.} Xác định giá trị hoặc khẳng định đúng cho bài toán sau: Một hộp có 22 quả bóng được đánh số từ 1 đến 22. Lấy ngẫu nhiên 1 quả. Số kết quả thuận lợi cho biến cố 'Lấy được quả bóng ghi số nhỏ hơn 5' là:\par
\hspace*{1.5em}\textbf{A.} 5 kết quả\quad\quad \textbf{B.} 3 kết quả\quad\quad \textbf{C.} 4 kết quả\quad\quad \textbf{D.} 22 kết quả\par\vspace{2pt}
\textbf{\color{primary} Câu 8.} Áp dụng trực tiếp định nghĩa và quy tắc trong Chương VIII: Nếu một biến cố E có k kết quả thuận lợi trong tổng số n kết quả đồng khả năng thì xác suất P(E) được tính bằng công thức:\par
\hspace*{1.5em}\textbf{A.} P(E) = k $\cdot$ n\quad\quad \textbf{B.} P(E) = n - k\quad\quad \textbf{C.} P(E) = n / k\quad\quad \textbf{D.} P(E) = k / n\par\vspace{2pt}
\textbf{\color{primary} Câu 9.} Kiểm tra nhanh kiến thức nền tảng: Khi số phép thử n ngày càng lớn thì xác suất thực nghiệm của một biến cố sẽ như thế nào?\par
\begin{itemize}\setlength{\itemsep}{1pt}\setlength{\parskip}{0pt}
  \item[] \textbf{A.} Tăng lên vô hạn
  \item[] \textbf{B.} Luôn bằng 1
  \item[] \textbf{C.} Ngày càng gần với xác suất lý thuyết của biến cố đó
  \item[] \textbf{D.} Luôn bằng 0
\end{itemize}\vspace{2pt}
\textbf{\color{primary} Câu 10.} Cho dữ kiện nhận biết sau đây: Một hộp có 28 quả bóng được đánh số từ 1 đến 28. Lấy ngẫu nhiên 1 quả. Số kết quả thuận lợi cho biến cố 'Lấy được quả bóng ghi số nhỏ hơn 5' là. Phương án chính xác là:\par
\hspace*{1.5em}\textbf{A.} 4 kết quả\quad\quad \textbf{B.} 28 kết quả\quad\quad \textbf{C.} 3 kết quả\quad\quad \textbf{D.} 5 kết quả\par\vspace{2pt}
\textbf{\color{primary} Câu 11.} Hãy chỉ ra kết quả đúng trong bốn lựa chọn A, B, C, D đối với yêu cầu: Nếu một biến cố E có k kết quả thuận lợi trong tổng số n kết quả đồng khả năng thì xác suất P(E) được tính bằng công thức:\par
\hspace*{1.5em}\textbf{A.} P(E) = n / k\quad\quad \textbf{B.} P(E) = n - k\quad\quad \textbf{C.} P(E) = k $\cdot$ n\quad\quad \textbf{D.} P(E) = k / n\par\vspace{2pt}
\textbf{\color{primary} Câu 12.} Câu hỏi ôn tập lý thuyết và kỹ năng cơ bản (Chương VIII): Khi số phép thử n ngày càng lớn thì xác suất thực nghiệm của một biến cố sẽ như thế nào:\par
\begin{itemize}\setlength{\itemsep}{1pt}\setlength{\parskip}{0pt}
  \item[] \textbf{A.} Luôn bằng 1
  \item[] \textbf{B.} Luôn bằng 0
  \item[] \textbf{C.} Tăng lên vô hạn
  \item[] \textbf{D.} Ngày càng gần với xác suất lý thuyết của biến cố đó
\end{itemize}\vspace{2pt}
\textbf{\color{primary} Câu 13.} Theo nội dung sách giáo khoa Kết nối tri thức với cuộc sống (Chương VIII), một hộp có 34 quả bóng được đánh số từ 1 đến 34. Lấy ngẫu nhiên 1 quả. Số kết quả thuận lợi cho biến cố 'Lấy được quả bóng ghi số nhỏ hơn 5' là:\par
\hspace*{1.5em}\textbf{A.} 5 kết quả\quad\quad \textbf{B.} 3 kết quả\quad\quad \textbf{C.} 4 kết quả\quad\quad \textbf{D.} 34 kết quả\par\vspace{2pt}
\textbf{\color{primary} Câu 14.} Tìm đáp số đúng của bài toán cơ bản sau: Nếu một biến cố E có k kết quả thuận lợi trong tổng số n kết quả đồng khả năng thì xác suất P(E) được tính bằng công thức:\par
\hspace*{1.5em}\textbf{A.} P(E) = n / k\quad\quad \textbf{B.} P(E) = n - k\quad\quad \textbf{C.} P(E) = k $\cdot$ n\quad\quad \textbf{D.} P(E) = k / n\par\vspace{2pt}
\textbf{\color{primary} Câu 15.} Xét đối tượng toán học trong Chương VIII: Khi số phép thử n ngày càng lớn thì xác suất thực nghiệm của một biến cố sẽ như thế nào. Giá trị hoặc tính chất tương ứng là:\par
\begin{itemize}\setlength{\itemsep}{1pt}\setlength{\parskip}{0pt}
  \item[] \textbf{A.} Luôn bằng 1
  \item[] \textbf{B.} Luôn bằng 0
  \item[] \textbf{C.} Ngày càng gần với xác suất lý thuyết của biến cố đó
  \item[] \textbf{D.} Tăng lên vô hạn
\end{itemize}\vspace{2pt}
\textbf{\color{primary} Câu 16.} Hoàn thành phép xác định cơ bản sau đây: Một hộp có 40 quả bóng được đánh số từ 1 đến 40. Lấy ngẫu nhiên 1 quả. Số kết quả thuận lợi cho biến cố 'Lấy được quả bóng ghi số nhỏ hơn 5' là:\par
\hspace*{1.5em}\textbf{A.} 4 kết quả\quad\quad \textbf{B.} 40 kết quả\quad\quad \textbf{C.} 5 kết quả\quad\quad \textbf{D.} 3 kết quả\par\vspace{2pt}
\textbf{\color{primary} Câu 17.} Trong các khẳng định hoặc giá trị dưới đây, đâu là đáp án đúng cho câu hỏi: Nếu một biến cố E có k kết quả thuận lợi trong tổng số n kết quả đồng khả năng thì xác suất P(E) được tính bằng công thức?\par
\hspace*{1.5em}\textbf{A.} P(E) = k / n\quad\quad \textbf{B.} P(E) = n - k\quad\quad \textbf{C.} P(E) = k $\cdot$ n\quad\quad \textbf{D.} P(E) = n / k\par\vspace{2pt}
\textbf{\color{primary} Câu 18.} Bài tập khởi động Chương VIII: Khi số phép thử n ngày càng lớn thì xác suất thực nghiệm của một biến cố sẽ như thế nào:\par
\begin{itemize}\setlength{\itemsep}{1pt}\setlength{\parskip}{0pt}
  \item[] \textbf{A.} Luôn bằng 1
  \item[] \textbf{B.} Luôn bằng 0
  \item[] \textbf{C.} Ngày càng gần với xác suất lý thuyết của biến cố đó
  \item[] \textbf{D.} Tăng lên vô hạn
\end{itemize}\vspace{2pt}
\textbf{\color{primary} Câu 19.} Quan sát dữ kiện đề bài: Một hộp có 46 quả bóng được đánh số từ 1 đến 46. Lấy ngẫu nhiên 1 quả. Số kết quả thuận lợi cho biến cố 'Lấy được quả bóng ghi số nhỏ hơn 5' là. Hãy chọn phương án đúng:\par
\hspace*{1.5em}\textbf{A.} 46 kết quả\quad\quad \textbf{B.} 4 kết quả\quad\quad \textbf{C.} 3 kết quả\quad\quad \textbf{D.} 5 kết quả\par\vspace{2pt}
\textbf{\color{primary} Câu 20.} Sử dụng công thức hoặc tính chất đầu bài của Chương VIII để giải quyết yêu cầu: Nếu một biến cố E có k kết quả thuận lợi trong tổng số n kết quả đồng khả năng thì xác suất P(E) được tính bằng công thức:\par
\hspace*{1.5em}\textbf{A.} P(E) = n / k\quad\quad \textbf{B.} P(E) = n - k\quad\quad \textbf{C.} P(E) = k $\cdot$ n\quad\quad \textbf{D.} P(E) = k / n\par\vspace{2pt}
\subsection*{\color{accent} Phần II. Mức độ Trung bình (Thông hiểu -- Câu 21 đến Câu 40)}
\textbf{\color{primary} Câu 21.} Một hộp có 4 bi trắng, 5 bi xanh và 6 bi vàng. Lấy ngẫu nhiên 1 viên bi. Xác suất để lấy được viên bi KHÔNG phải màu trắng là:\par
\hspace*{1.5em}\textbf{A.} 11/15\quad\quad \textbf{B.} 6/15\quad\quad \textbf{C.} 5/15\quad\quad \textbf{D.} 4/15\par\vspace{2pt}
\textbf{\color{primary} Câu 22.} Vận dụng quy tắc trong Chương VIII (Mở đầu về tính xác suất của biến cố) để giải bài toán: Viết ngẫu nhiên một số tự nhiên có một chữ số (từ 0 đến 9). Xác suất của biến cố 'Số viết ra là ước của 6' bằng:\par
\hspace*{1.5em}\textbf{A.} 3/10\quad\quad \textbf{B.} 1/5\quad\quad \textbf{C.} 2/5\quad\quad \textbf{D.} 1/2\par\vspace{2pt}
\textbf{\color{primary} Câu 23.} Biến đổi và tính toán theo kiến thức Chương VIII: Trong một lô 700 chiếc bóng đèn có 15 bóng bị lỗi. Chọn ngẫu nhiên 1 bóng đèn. Xác suất chọn phải bóng đèn bị lỗi là. Kết quả thu được là:\par
\hspace*{1.5em}\textbf{A.} 685/700\quad\quad \textbf{B.} 15/700\quad\quad \textbf{C.} 15/100\quad\quad \textbf{D.} 1/15\par\vspace{2pt}
\textbf{\color{primary} Câu 24.} Xét bài toán mức độ thông hiểu sau: Một hộp có 7 bi trắng, 8 bi xanh và 9 bi vàng. Lấy ngẫu nhiên 1 viên bi. Xác suất để lấy được viên bi KHÔNG phải màu trắng là. Đáp án đúng là:\par
\hspace*{1.5em}\textbf{A.} 17/24\quad\quad \textbf{B.} 9/24\quad\quad \textbf{C.} 7/24\quad\quad \textbf{D.} 8/24\par\vspace{2pt}
\textbf{\color{primary} Câu 25.} Để hoàn thành phiếu học tập Chương VIII, hãy xác định: Viết ngẫu nhiên một số tự nhiên có một chữ số (từ 0 đến 9). Xác suất của biến cố 'Số viết ra là ước của 6' bằng:\par
\hspace*{1.5em}\textbf{A.} 2/5\quad\quad \textbf{B.} 3/10\quad\quad \textbf{C.} 1/2\quad\quad \textbf{D.} 1/5\par\vspace{2pt}
\textbf{\color{primary} Câu 26.} Phân tích dữ kiện bài toán: Trong một lô 600 chiếc bóng đèn có 15 bóng bị lỗi. Chọn ngẫu nhiên 1 bóng đèn. Xác suất chọn phải bóng đèn bị lỗi là. Giá trị cần tìm bằng:\par
\hspace*{1.5em}\textbf{A.} 585/600\quad\quad \textbf{B.} 15/600\quad\quad \textbf{C.} 1/15\quad\quad \textbf{D.} 15/100\par\vspace{2pt}
\textbf{\color{primary} Câu 27.} Áp dụng tính chất đã học ở Chương VIII vào bài tập sau: Một hộp có 10 bi trắng, 11 bi xanh và 12 bi vàng. Lấy ngẫu nhiên 1 viên bi. Xác suất để lấy được viên bi KHÔNG phải màu trắng là:\par
\hspace*{1.5em}\textbf{A.} 10/33\quad\quad \textbf{B.} 12/33\quad\quad \textbf{C.} 11/33\quad\quad \textbf{D.} 23/33\par\vspace{2pt}
\textbf{\color{primary} Câu 28.} Thực hiện các bước suy luận hợp lý để giải quyết yêu cầu: Viết ngẫu nhiên một số tự nhiên có một chữ số (từ 0 đến 9). Xác suất của biến cố 'Số viết ra là ước của 6' bằng:\par
\hspace*{1.5em}\textbf{A.} 2/5\quad\quad \textbf{B.} 1/5\quad\quad \textbf{C.} 1/2\quad\quad \textbf{D.} 3/10\par\vspace{2pt}
\textbf{\color{primary} Câu 29.} Cho bài toán tự luyện (Mở đầu về tính xác suất của biến cố): Trong một lô 500 chiếc bóng đèn có 15 bóng bị lỗi. Chọn ngẫu nhiên 1 bóng đèn. Xác suất chọn phải bóng đèn bị lỗi là. Phương án nào sau đây chính xác?\par
\hspace*{1.5em}\textbf{A.} 15/100\quad\quad \textbf{B.} 1/15\quad\quad \textbf{C.} 15/500\quad\quad \textbf{D.} 485/500\par\vspace{2pt}
\textbf{\color{primary} Câu 30.} Kết hợp các phép biến đổi trong Chương VIII, hãy tìm kết quả của bài toán: Một hộp có 13 bi trắng, 14 bi xanh và 15 bi vàng. Lấy ngẫu nhiên 1 viên bi. Xác suất để lấy được viên bi KHÔNG phải màu trắng là:\par
\hspace*{1.5em}\textbf{A.} 29/42\quad\quad \textbf{B.} 14/42\quad\quad \textbf{C.} 13/42\quad\quad \textbf{D.} 15/42\par\vspace{2pt}
\textbf{\color{primary} Câu 31.} Một bài kiểm tra định kỳ nội dung Chương VIII nêu câu hỏi: Viết ngẫu nhiên một số tự nhiên có một chữ số (từ 0 đến 9). Xác suất của biến cố 'Số viết ra là ước của 6' bằng. Hãy chọn đáp số đúng:\par
\hspace*{1.5em}\textbf{A.} 1/5\quad\quad \textbf{B.} 3/10\quad\quad \textbf{C.} 2/5\quad\quad \textbf{D.} 1/2\par\vspace{2pt}
\textbf{\color{primary} Câu 32.} Tìm lời giải đúng cho bài toán thông hiểu: Trong một lô 800 chiếc bóng đèn có 15 bóng bị lỗi. Chọn ngẫu nhiên 1 bóng đèn. Xác suất chọn phải bóng đèn bị lỗi là:\par
\hspace*{1.5em}\textbf{A.} 1/15\quad\quad \textbf{B.} 785/800\quad\quad \textbf{C.} 15/800\quad\quad \textbf{D.} 15/100\par\vspace{2pt}
\textbf{\color{primary} Câu 33.} Sử dụng mối liên hệ giữa các đại lượng trong Chương VIII: Một hộp có 16 bi trắng, 17 bi xanh và 18 bi vàng. Lấy ngẫu nhiên 1 viên bi. Xác suất để lấy được viên bi KHÔNG phải màu trắng là:\par
\hspace*{1.5em}\textbf{A.} 18/51\quad\quad \textbf{B.} 16/51\quad\quad \textbf{C.} 17/51\quad\quad \textbf{D.} 35/51\par\vspace{2pt}
\textbf{\color{primary} Câu 34.} Hãy tính toán và chọn phương án đúng đối với bài tập: Viết ngẫu nhiên một số tự nhiên có một chữ số (từ 0 đến 9). Xác suất của biến cố 'Số viết ra là ước của 6' bằng:\par
\hspace*{1.5em}\textbf{A.} 3/10\quad\quad \textbf{B.} 1/5\quad\quad \textbf{C.} 1/2\quad\quad \textbf{D.} 2/5\par\vspace{2pt}
\textbf{\color{primary} Câu 35.} Dựa trên các tính chất trọng tâm của Chương VIII, xét bài toán: Trong một lô 700 chiếc bóng đèn có 15 bóng bị lỗi. Chọn ngẫu nhiên 1 bóng đèn. Xác suất chọn phải bóng đèn bị lỗi là:\par
\hspace*{1.5em}\textbf{A.} 1/15\quad\quad \textbf{B.} 15/100\quad\quad \textbf{C.} 685/700\quad\quad \textbf{D.} 15/700\par\vspace{2pt}
\textbf{\color{primary} Câu 36.} Thực hành giải toán (Mở đầu về tính xác suất của biến cố): Một hộp có 19 bi trắng, 20 bi xanh và 21 bi vàng. Lấy ngẫu nhiên 1 viên bi. Xác suất để lấy được viên bi KHÔNG phải màu trắng là. Đáp án nào dưới đây thỏa mãn?\par
\hspace*{1.5em}\textbf{A.} 20/60\quad\quad \textbf{B.} 21/60\quad\quad \textbf{C.} 19/60\quad\quad \textbf{D.} 41/60\par\vspace{2pt}
\textbf{\color{primary} Câu 37.} Kiểm tra kỹ năng biến đổi và suy luận qua bài toán: Viết ngẫu nhiên một số tự nhiên có một chữ số (từ 0 đến 9). Xác suất của biến cố 'Số viết ra là ước của 6' bằng:\par
\hspace*{1.5em}\textbf{A.} 1/5\quad\quad \textbf{B.} 2/5\quad\quad \textbf{C.} 1/2\quad\quad \textbf{D.} 3/10\par\vspace{2pt}
\textbf{\color{primary} Câu 38.} Cho giả thiết thuộc phạm vi Chương VIII: Trong một lô 600 chiếc bóng đèn có 15 bóng bị lỗi. Chọn ngẫu nhiên 1 bóng đèn. Xác suất chọn phải bóng đèn bị lỗi là. Kết luận đúng là:\par
\hspace*{1.5em}\textbf{A.} 15/100\quad\quad \textbf{B.} 1/15\quad\quad \textbf{C.} 15/600\quad\quad \textbf{D.} 585/600\par\vspace{2pt}
\textbf{\color{primary} Câu 39.} Rút gọn hoặc tính giá trị theo yêu cầu sau của Chương VIII: Một hộp có 22 bi trắng, 23 bi xanh và 24 bi vàng. Lấy ngẫu nhiên 1 viên bi. Xác suất để lấy được viên bi KHÔNG phải màu trắng là:\par
\hspace*{1.5em}\textbf{A.} 47/69\quad\quad \textbf{B.} 22/69\quad\quad \textbf{C.} 24/69\quad\quad \textbf{D.} 23/69\par\vspace{2pt}
\textbf{\color{primary} Câu 40.} Xác định phương án đúng sau khi thực hiện phép tính / suy luận: Viết ngẫu nhiên một số tự nhiên có một chữ số (từ 0 đến 9). Xác suất của biến cố 'Số viết ra là ước của 6' bằng:\par
\hspace*{1.5em}\textbf{A.} 3/10\quad\quad \textbf{B.} 2/5\quad\quad \textbf{C.} 1/5\quad\quad \textbf{D.} 1/2\par\vspace{2pt}
\subsection*{\color{accent} Phần III. Mức độ Khó (Vận dụng \& Vận dụng cao -- Câu 41 đến Câu 50)}
\textbf{\color{primary} Câu 41.} Để ước lượng số cá trong hồ, người ta bắt lên 40 con cá, đánh dấu rồi thả lại xuống hồ. Sau đó bắt ngẫu nhiên 100 con thì thấy có 8 con đã được đánh dấu. Số cá trong hồ ước lượng khoảng:\par
\hspace*{1.5em}\textbf{A.} 500 con\quad\quad \textbf{B.} 400 con\quad\quad \textbf{C.} 140 con\quad\quad \textbf{D.} 320 con\par\vspace{2pt}
\textbf{\color{primary} Câu 42.} Bài toán vận dụng (Chương VIII – Mở đầu về tính xác suất của biến cố): Gieo đồng thời hai đồng xu cân đối và đồng chất. Xác suất để hai đồng xu xuất hiện hai mặt khác nhau (một Sấp, một Ngửa) bằng. Kết quả đúng là:\par
\hspace*{1.5em}\textbf{A.} 3/4\quad\quad \textbf{B.} 1/2\quad\quad \textbf{C.} 1/3\quad\quad \textbf{D.} 1/4\par\vspace{2pt}
\textbf{\color{primary} Câu 43.} Vận dụng kiến thức nâng cao trong phạm vi Chương VIII để giải bài toán: Một hộp chứa 15 viên bi gồm bi xanh và bi đỏ, biết xác suất lấy ngẫu nhiên được 1 viên bi xanh là 1/3. Số viên bi đỏ có trong hộp là:\par
\hspace*{1.5em}\textbf{A.} 10 viên bi đỏ\quad\quad \textbf{B.} 14 viên bi đỏ\quad\quad \textbf{C.} 5 viên bi đỏ\quad\quad \textbf{D.} 8 viên bi đỏ\par\vspace{2pt}
\textbf{\color{primary} Câu 44.} Phân tích và giải bài toán tổng hợp sau của Chương VIII: Để ước lượng số cá trong hồ, người ta bắt lên 40 con cá, đánh dấu rồi thả lại xuống hồ. Sau đó bắt ngẫu nhiên 100 con thì thấy có 8 con đã được đánh dấu. Số cá trong hồ ước lượng khoảng:\par
\hspace*{1.5em}\textbf{A.} 500 con\quad\quad \textbf{B.} 400 con\quad\quad \textbf{C.} 320 con\quad\quad \textbf{D.} 140 con\par\vspace{2pt}
\textbf{\color{primary} Câu 45.} Thử thách tư duy trong phạm vi kiến thức Chương VIII: Gieo đồng thời hai đồng xu cân đối và đồng chất. Xác suất để hai đồng xu xuất hiện hai mặt khác nhau (một Sấp, một Ngửa) bằng. Đáp số chính xác là:\par
\hspace*{1.5em}\textbf{A.} 3/4\quad\quad \textbf{B.} 1/3\quad\quad \textbf{C.} 1/2\quad\quad \textbf{D.} 1/4\par\vspace{2pt}
\textbf{\color{primary} Câu 46.} Sử dụng phương pháp giải toán của Chương VIII để tìm đáp án cho bài toán: Một hộp chứa 24 viên bi gồm bi xanh và bi đỏ, biết xác suất lấy ngẫu nhiên được 1 viên bi xanh là 1/3. Số viên bi đỏ có trong hộp là:\par
\hspace*{1.5em}\textbf{A.} 11 viên bi đỏ\quad\quad \textbf{B.} 16 viên bi đỏ\quad\quad \textbf{C.} 23 viên bi đỏ\quad\quad \textbf{D.} 8 viên bi đỏ\par\vspace{2pt}
\textbf{\color{primary} Câu 47.} Xét bài toán vận dụng thực tiễn / nâng cao thuộc Chương VIII: Để ước lượng số cá trong hồ, người ta bắt lên 40 con cá, đánh dấu rồi thả lại xuống hồ. Sau đó bắt ngẫu nhiên 100 con thì thấy có 8 con đã được đánh dấu. Số cá trong hồ ước lượng khoảng:\par
\hspace*{1.5em}\textbf{A.} 500 con\quad\quad \textbf{B.} 320 con\quad\quad \textbf{C.} 400 con\quad\quad \textbf{D.} 140 con\par\vspace{2pt}
\textbf{\color{primary} Câu 48.} Tìm kết quả cuối cùng của bài toán phân loại học sinh khá giỏi (Chương VIII): Gieo đồng thời hai đồng xu cân đối và đồng chất. Xác suất để hai đồng xu xuất hiện hai mặt khác nhau (một Sấp, một Ngửa) bằng:\par
\hspace*{1.5em}\textbf{A.} 1/2\quad\quad \textbf{B.} 1/4\quad\quad \textbf{C.} 3/4\quad\quad \textbf{D.} 1/3\par\vspace{2pt}
\textbf{\color{primary} Câu 49.} Kết hợp các tính chất đã học trong Chương VIII, hãy giải quyết bài toán: Một hộp chứa 33 viên bi gồm bi xanh và bi đỏ, biết xác suất lấy ngẫu nhiên được 1 viên bi xanh là 1/3. Số viên bi đỏ có trong hộp là:\par
\hspace*{1.5em}\textbf{A.} 14 viên bi đỏ\quad\quad \textbf{B.} 32 viên bi đỏ\quad\quad \textbf{C.} 11 viên bi đỏ\quad\quad \textbf{D.} 22 viên bi đỏ\par\vspace{2pt}
\textbf{\color{primary} Câu 50.} Câu hỏi vận dụng cao cuối Chương VIII: Để ước lượng số cá trong hồ, người ta bắt lên 40 con cá, đánh dấu rồi thả lại xuống hồ. Sau đó bắt ngẫu nhiên 100 con thì thấy có 8 con đã được đánh dấu. Số cá trong hồ ước lượng khoảng. Phương án đúng là:\par
\hspace*{1.5em}\textbf{A.} 140 con\quad\quad \textbf{B.} 320 con\quad\quad \textbf{C.} 500 con\quad\quad \textbf{D.} 400 con\par\vspace{2pt}
\subsection*{\color{primary} Bảng đáp án Chương VIII: Mở đầu về tính xác suất của biến cố}
\begin{center}
\renewcommand{\arraystretch}{1.25}
\begin{tabularx}{\textwidth}{*{10}{>{\centering\arraybackslash}X}}
\toprule
\textbf{\color{primary} Cột 1} & \textbf{\color{primary} Cột 2} & \textbf{\color{primary} Cột 3} & \textbf{\color{primary} Cột 4} & \textbf{\color{primary} Cột 5} & \textbf{\color{primary} Cột 6} & \textbf{\color{primary} Cột 7} & \textbf{\color{primary} Cột 8} & \textbf{\color{primary} Cột 9} & \textbf{\color{primary} Cột 10} \\
\midrule
\textbf{1.A} & \textbf{2.B} & \textbf{3.B} & \textbf{4.D} & \textbf{5.D} & \textbf{6.A} & \textbf{7.C} & \textbf{8.D} & \textbf{9.C} & \textbf{10.A} \\
\rowcolor{rowtint} \textbf{11.D} & \textbf{12.D} & \textbf{13.C} & \textbf{14.D} & \textbf{15.C} & \textbf{16.A} & \textbf{17.A} & \textbf{18.C} & \textbf{19.B} & \textbf{20.D} \\
\textbf{21.A} & \textbf{22.C} & \textbf{23.B} & \textbf{24.A} & \textbf{25.A} & \textbf{26.B} & \textbf{27.D} & \textbf{28.A} & \textbf{29.C} & \textbf{30.A} \\
\rowcolor{rowtint} \textbf{31.C} & \textbf{32.C} & \textbf{33.D} & \textbf{34.D} & \textbf{35.D} & \textbf{36.D} & \textbf{37.B} & \textbf{38.C} & \textbf{39.A} & \textbf{40.B} \\
\textbf{41.A} & \textbf{42.B} & \textbf{43.A} & \textbf{44.A} & \textbf{45.C} & \textbf{46.B} & \textbf{47.A} & \textbf{48.A} & \textbf{49.D} & \textbf{50.C} \\
\bottomrule
\end{tabularx}
\end{center}\vspace{10pt}
\section*{\color{primary} CHƯƠNG IX: TAM GIÁC ĐỒNG DẠNG}
\textit{Trọng tâm kiến thức SGK: Hai tam giác đồng dạng, ba trường hợp đồng dạng của hai tam giác, định lí Pythagore và ứng dụng, các trường hợp đồng dạng của hai tam giác vuông.}\par\vspace{4pt}
\subsection*{\color{accent} Phần I. Mức độ Dễ (Nhận biết -- Câu 1 đến Câu 20)}
\textbf{\color{primary} Câu 1.} Cho $\Delta$ABC $\sim$ $\Delta$DEF theo tỉ số đồng dạng k = 2/3. Khi đó tỉ số chu vi của $\Delta$ABC và $\Delta$DEF bằng:\par
\hspace*{1.5em}\textbf{A.} 4/9\quad\quad \textbf{B.} 2/3\quad\quad \textbf{C.} 3/2\quad\quad \textbf{D.} 9/4\par\vspace{2pt}
\textbf{\color{primary} Câu 2.} Thực hiện yêu cầu nhận biết sau: Tam giác ABC vuông tại A có hai cạnh góc vuông AB = 6 cm và AC = 8 cm. Theo định lí Pythagore, độ dài cạnh huyền BC bằng.\par
\hspace*{1.5em}\textbf{A.} 10 cm\quad\quad \textbf{B.} 14 cm\quad\quad \textbf{C.} 50 cm\quad\quad \textbf{D.} 12 cm\par\vspace{2pt}
\textbf{\color{primary} Câu 3.} Trong bài học thuộc Chương IX, xét nội dung: Cho $\Delta$ABC $\sim$ $\Delta$MNP theo tỉ số đồng dạng k = 4. Tỉ số diện tích S(ABC) / S(MNP) bằng. Đáp án đúng là:\par
\hspace*{1.5em}\textbf{A.} 4\quad\quad \textbf{B.} 8\quad\quad \textbf{C.} 1/16\quad\quad \textbf{D.} 16\par\vspace{2pt}
\textbf{\color{primary} Câu 4.} Dựa vào kiến thức chuẩn SGK của Chương IX (Tam giác đồng dạng), cho $\Delta$ABC $\sim$ $\Delta$DEF theo tỉ số đồng dạng k = 5/3. Khi đó tỉ số chu vi của $\Delta$ABC và $\Delta$DEF bằng:\par
\hspace*{1.5em}\textbf{A.} 3/5\quad\quad \textbf{B.} 5/3\quad\quad \textbf{C.} 25/9\quad\quad \textbf{D.} 9/25\par\vspace{2pt}
\textbf{\color{primary} Câu 5.} Chọn kết quả đúng nhất khi xét: Tam giác ABC vuông tại A có hai cạnh góc vuông AB = 6 cm và AC = 8 cm. Theo định lí Pythagore, độ dài cạnh huyền BC bằng:\par
\hspace*{1.5em}\textbf{A.} 50 cm\quad\quad \textbf{B.} 12 cm\quad\quad \textbf{C.} 10 cm\quad\quad \textbf{D.} 14 cm\par\vspace{2pt}
\textbf{\color{primary} Câu 6.} Khi làm bài tập cơ bản về tam giác đồng dạng, một học sinh gặp câu hỏi: Cho $\Delta$ABC $\sim$ $\Delta$MNP theo tỉ số đồng dạng k = 7. Tỉ số diện tích S(ABC) / S(MNP) bằng. Kết quả nào dưới đây đúng?\par
\hspace*{1.5em}\textbf{A.} 49\quad\quad \textbf{B.} 1/49\quad\quad \textbf{C.} 14\quad\quad \textbf{D.} 7\par\vspace{2pt}
\textbf{\color{primary} Câu 7.} Xác định giá trị hoặc khẳng định đúng cho bài toán sau: Cho $\Delta$ABC $\sim$ $\Delta$DEF theo tỉ số đồng dạng k = 8/3. Khi đó tỉ số chu vi của $\Delta$ABC và $\Delta$DEF bằng:\par
\hspace*{1.5em}\textbf{A.} 9/64\quad\quad \textbf{B.} 64/9\quad\quad \textbf{C.} 3/8\quad\quad \textbf{D.} 8/3\par\vspace{2pt}
\textbf{\color{primary} Câu 8.} Áp dụng trực tiếp định nghĩa và quy tắc trong Chương IX: Tam giác ABC vuông tại A có hai cạnh góc vuông AB = 6 cm và AC = 8 cm. Theo định lí Pythagore, độ dài cạnh huyền BC bằng:\par
\hspace*{1.5em}\textbf{A.} 10 cm\quad\quad \textbf{B.} 14 cm\quad\quad \textbf{C.} 12 cm\quad\quad \textbf{D.} 50 cm\par\vspace{2pt}
\textbf{\color{primary} Câu 9.} Kiểm tra nhanh kiến thức nền tảng: Cho $\Delta$ABC $\sim$ $\Delta$MNP theo tỉ số đồng dạng k = 10. Tỉ số diện tích S(ABC) / S(MNP) bằng:\par
\hspace*{1.5em}\textbf{A.} 100\quad\quad \textbf{B.} 10\quad\quad \textbf{C.} 1/100\quad\quad \textbf{D.} 20\par\vspace{2pt}
\textbf{\color{primary} Câu 10.} Cho dữ kiện nhận biết sau đây: Cho $\Delta$ABC $\sim$ $\Delta$DEF theo tỉ số đồng dạng k = 11/3. Khi đó tỉ số chu vi của $\Delta$ABC và $\Delta$DEF bằng. Phương án chính xác là:\par
\hspace*{1.5em}\textbf{A.} 121/9\quad\quad \textbf{B.} 9/121\quad\quad \textbf{C.} 3/11\quad\quad \textbf{D.} 11/3\par\vspace{2pt}
\textbf{\color{primary} Câu 11.} Hãy chỉ ra kết quả đúng trong bốn lựa chọn A, B, C, D đối với yêu cầu: Tam giác ABC vuông tại A có hai cạnh góc vuông AB = 6 cm và AC = 8 cm. Theo định lí Pythagore, độ dài cạnh huyền BC bằng:\par
\hspace*{1.5em}\textbf{A.} 14 cm\quad\quad \textbf{B.} 50 cm\quad\quad \textbf{C.} 10 cm\quad\quad \textbf{D.} 12 cm\par\vspace{2pt}
\textbf{\color{primary} Câu 12.} Câu hỏi ôn tập lý thuyết và kỹ năng cơ bản (Chương IX): Cho $\Delta$ABC $\sim$ $\Delta$MNP theo tỉ số đồng dạng k = 13. Tỉ số diện tích S(ABC) / S(MNP) bằng:\par
\hspace*{1.5em}\textbf{A.} 1/169\quad\quad \textbf{B.} 169\quad\quad \textbf{C.} 13\quad\quad \textbf{D.} 26\par\vspace{2pt}
\textbf{\color{primary} Câu 13.} Theo nội dung sách giáo khoa Kết nối tri thức với cuộc sống (Chương IX), cho $\Delta$ABC $\sim$ $\Delta$DEF theo tỉ số đồng dạng k = 14/3. Khi đó tỉ số chu vi của $\Delta$ABC và $\Delta$DEF bằng:\par
\hspace*{1.5em}\textbf{A.} 14/3\quad\quad \textbf{B.} 3/14\quad\quad \textbf{C.} 196/9\quad\quad \textbf{D.} 9/196\par\vspace{2pt}
\textbf{\color{primary} Câu 14.} Tìm đáp số đúng của bài toán cơ bản sau: Tam giác ABC vuông tại A có hai cạnh góc vuông AB = 6 cm và AC = 8 cm. Theo định lí Pythagore, độ dài cạnh huyền BC bằng:\par
\hspace*{1.5em}\textbf{A.} 50 cm\quad\quad \textbf{B.} 14 cm\quad\quad \textbf{C.} 10 cm\quad\quad \textbf{D.} 12 cm\par\vspace{2pt}
\textbf{\color{primary} Câu 15.} Xét đối tượng toán học trong Chương IX: Cho $\Delta$ABC $\sim$ $\Delta$MNP theo tỉ số đồng dạng k = 16. Tỉ số diện tích S(ABC) / S(MNP) bằng. Giá trị hoặc tính chất tương ứng là:\par
\hspace*{1.5em}\textbf{A.} 32\quad\quad \textbf{B.} 256\quad\quad \textbf{C.} 16\quad\quad \textbf{D.} 1/256\par\vspace{2pt}
\textbf{\color{primary} Câu 16.} Hoàn thành phép xác định cơ bản sau đây: Cho $\Delta$ABC $\sim$ $\Delta$DEF theo tỉ số đồng dạng k = 17/3. Khi đó tỉ số chu vi của $\Delta$ABC và $\Delta$DEF bằng:\par
\hspace*{1.5em}\textbf{A.} 17/3\quad\quad \textbf{B.} 3/17\quad\quad \textbf{C.} 9/289\quad\quad \textbf{D.} 289/9\par\vspace{2pt}
\textbf{\color{primary} Câu 17.} Trong các khẳng định hoặc giá trị dưới đây, đâu là đáp án đúng cho câu hỏi: Tam giác ABC vuông tại A có hai cạnh góc vuông AB = 6 cm và AC = 8 cm. Theo định lí Pythagore, độ dài cạnh huyền BC bằng?\par
\hspace*{1.5em}\textbf{A.} 10 cm\quad\quad \textbf{B.} 12 cm\quad\quad \textbf{C.} 50 cm\quad\quad \textbf{D.} 14 cm\par\vspace{2pt}
\textbf{\color{primary} Câu 18.} Bài tập khởi động Chương IX: Cho $\Delta$ABC $\sim$ $\Delta$MNP theo tỉ số đồng dạng k = 19. Tỉ số diện tích S(ABC) / S(MNP) bằng:\par
\hspace*{1.5em}\textbf{A.} 1/361\quad\quad \textbf{B.} 38\quad\quad \textbf{C.} 361\quad\quad \textbf{D.} 19\par\vspace{2pt}
\textbf{\color{primary} Câu 19.} Quan sát dữ kiện đề bài: Cho $\Delta$ABC $\sim$ $\Delta$DEF theo tỉ số đồng dạng k = 20/3. Khi đó tỉ số chu vi của $\Delta$ABC và $\Delta$DEF bằng. Hãy chọn phương án đúng:\par
\hspace*{1.5em}\textbf{A.} 3/20\quad\quad \textbf{B.} 400/9\quad\quad \textbf{C.} 9/400\quad\quad \textbf{D.} 20/3\par\vspace{2pt}
\textbf{\color{primary} Câu 20.} Sử dụng công thức hoặc tính chất đầu bài của Chương IX để giải quyết yêu cầu: Tam giác ABC vuông tại A có hai cạnh góc vuông AB = 6 cm và AC = 8 cm. Theo định lí Pythagore, độ dài cạnh huyền BC bằng:\par
\hspace*{1.5em}\textbf{A.} 10 cm\quad\quad \textbf{B.} 50 cm\quad\quad \textbf{C.} 14 cm\quad\quad \textbf{D.} 12 cm\par\vspace{2pt}
\subsection*{\color{accent} Phần II. Mức độ Trung bình (Thông hiểu -- Câu 21 đến Câu 40)}
\textbf{\color{primary} Câu 21.} Cho $\Delta$ABC $\sim$ $\Delta$DEF theo tỉ số k = 2/3. Biết chu vi $\Delta$ABC bằng 12 cm, chu vi của $\Delta$DEF bằng:\par
\hspace*{1.5em}\textbf{A.} 24 cm\quad\quad \textbf{B.} 18 cm\quad\quad \textbf{C.} 8 cm\quad\quad \textbf{D.} 27 cm\par\vspace{2pt}
\textbf{\color{primary} Câu 22.} Vận dụng quy tắc trong Chương IX (Tam giác đồng dạng) để giải bài toán: Cho $\Delta$ABC $\sim$ $\Delta$A'B'C' theo tỉ số đồng dạng k = 1/3. Biết diện tích $\Delta$ABC bằng 12 cm$^2$, diện tích $\Delta$A'B'C' bằng:\par
\hspace*{1.5em}\textbf{A.} 24 cm$^2$\quad\quad \textbf{B.} 108 cm$^2$\quad\quad \textbf{C.} 4.0 cm$^2$\quad\quad \textbf{D.} 36 cm$^2$\par\vspace{2pt}
\textbf{\color{primary} Câu 23.} Biến đổi và tính toán theo kiến thức Chương IX: Bộ ba số đo nào sau đây là độ dài ba cạnh của một tam giác vuông. Kết quả thu được là:\par
\hspace*{1.5em}\textbf{A.} 4 cm, 5 cm, 6 cm\quad\quad \textbf{B.} 6 cm, 7 cm, 8 cm\quad\quad \textbf{C.} 9 cm, 12 cm, 15 cm\quad\quad \textbf{D.} 5 cm, 10 cm, 12 cm\par\vspace{2pt}
\textbf{\color{primary} Câu 24.} Xét bài toán mức độ thông hiểu sau: Cho $\Delta$ABC $\sim$ $\Delta$DEF theo tỉ số k = 2/3. Biết chu vi $\Delta$ABC bằng 24 cm, chu vi của $\Delta$DEF bằng. Đáp án đúng là:\par
\hspace*{1.5em}\textbf{A.} 36 cm\quad\quad \textbf{B.} 54 cm\quad\quad \textbf{C.} 48 cm\quad\quad \textbf{D.} 14 cm\par\vspace{2pt}
\textbf{\color{primary} Câu 25.} Để hoàn thành phiếu học tập Chương IX, hãy xác định: Cho $\Delta$ABC $\sim$ $\Delta$A'B'C' theo tỉ số đồng dạng k = 1/6. Biết diện tích $\Delta$ABC bằng 18 cm$^2$, diện tích $\Delta$A'B'C' bằng:\par
\hspace*{1.5em}\textbf{A.} 36 cm$^2$\quad\quad \textbf{B.} 108 cm$^2$\quad\quad \textbf{C.} 648 cm$^2$\quad\quad \textbf{D.} 3.0 cm$^2$\par\vspace{2pt}
\textbf{\color{primary} Câu 26.} Phân tích dữ kiện bài toán: Bộ ba số đo nào sau đây là độ dài ba cạnh của một tam giác vuông. Giá trị cần tìm bằng:\par
\hspace*{1.5em}\textbf{A.} 6 cm, 7 cm, 8 cm\quad\quad \textbf{B.} 8 cm, 15 cm, 17 cm\quad\quad \textbf{C.} 5 cm, 10 cm, 12 cm\quad\quad \textbf{D.} 4 cm, 5 cm, 6 cm\par\vspace{2pt}
\textbf{\color{primary} Câu 27.} Áp dụng tính chất đã học ở Chương IX vào bài tập sau: Cho $\Delta$ABC $\sim$ $\Delta$DEF theo tỉ số k = 2/3. Biết chu vi $\Delta$ABC bằng 36 cm, chu vi của $\Delta$DEF bằng:\par
\hspace*{1.5em}\textbf{A.} 72 cm\quad\quad \textbf{B.} 81 cm\quad\quad \textbf{C.} 20 cm\quad\quad \textbf{D.} 54 cm\par\vspace{2pt}
\textbf{\color{primary} Câu 28.} Thực hiện các bước suy luận hợp lý để giải quyết yêu cầu: Cho $\Delta$ABC $\sim$ $\Delta$A'B'C' theo tỉ số đồng dạng k = 1/9. Biết diện tích $\Delta$ABC bằng 24 cm$^2$, diện tích $\Delta$A'B'C' bằng:\par
\hspace*{1.5em}\textbf{A.} 1944 cm$^2$\quad\quad \textbf{B.} 2.7 cm$^2$\quad\quad \textbf{C.} 48 cm$^2$\quad\quad \textbf{D.} 216 cm$^2$\par\vspace{2pt}
\textbf{\color{primary} Câu 29.} Cho bài toán tự luyện (Tam giác đồng dạng): Bộ ba số đo nào sau đây là độ dài ba cạnh của một tam giác vuông. Phương án nào sau đây chính xác?\par
\hspace*{1.5em}\textbf{A.} 4 cm, 5 cm, 6 cm\quad\quad \textbf{B.} 5 cm, 12 cm, 13 cm\quad\quad \textbf{C.} 5 cm, 10 cm, 12 cm\quad\quad \textbf{D.} 6 cm, 7 cm, 8 cm\par\vspace{2pt}
\textbf{\color{primary} Câu 30.} Kết hợp các phép biến đổi trong Chương IX, hãy tìm kết quả của bài toán: Cho $\Delta$ABC $\sim$ $\Delta$DEF theo tỉ số k = 2/3. Biết chu vi $\Delta$ABC bằng 48 cm, chu vi của $\Delta$DEF bằng:\par
\hspace*{1.5em}\textbf{A.} 72 cm\quad\quad \textbf{B.} 108 cm\quad\quad \textbf{C.} 96 cm\quad\quad \textbf{D.} 26 cm\par\vspace{2pt}
\textbf{\color{primary} Câu 31.} Một bài kiểm tra định kỳ nội dung Chương IX nêu câu hỏi: Cho $\Delta$ABC $\sim$ $\Delta$A'B'C' theo tỉ số đồng dạng k = 1/12. Biết diện tích $\Delta$ABC bằng 30 cm$^2$, diện tích $\Delta$A'B'C' bằng. Hãy chọn đáp số đúng:\par
\hspace*{1.5em}\textbf{A.} 60 cm$^2$\quad\quad \textbf{B.} 2.5 cm$^2$\quad\quad \textbf{C.} 360 cm$^2$\quad\quad \textbf{D.} 4320 cm$^2$\par\vspace{2pt}
\textbf{\color{primary} Câu 32.} Tìm lời giải đúng cho bài toán thông hiểu: Bộ ba số đo nào sau đây là độ dài ba cạnh của một tam giác vuông:\par
\hspace*{1.5em}\textbf{A.} 5 cm, 10 cm, 12 cm\quad\quad \textbf{B.} 6 cm, 7 cm, 8 cm\quad\quad \textbf{C.} 7 cm, 24 cm, 25 cm\quad\quad \textbf{D.} 4 cm, 5 cm, 6 cm\par\vspace{2pt}
\textbf{\color{primary} Câu 33.} Sử dụng mối liên hệ giữa các đại lượng trong Chương IX: Cho $\Delta$ABC $\sim$ $\Delta$DEF theo tỉ số k = 2/3. Biết chu vi $\Delta$ABC bằng 60 cm, chu vi của $\Delta$DEF bằng:\par
\hspace*{1.5em}\textbf{A.} 32 cm\quad\quad \textbf{B.} 90 cm\quad\quad \textbf{C.} 120 cm\quad\quad \textbf{D.} 135 cm\par\vspace{2pt}
\textbf{\color{primary} Câu 34.} Hãy tính toán và chọn phương án đúng đối với bài tập: Cho $\Delta$ABC $\sim$ $\Delta$A'B'C' theo tỉ số đồng dạng k = 1/15. Biết diện tích $\Delta$ABC bằng 36 cm$^2$, diện tích $\Delta$A'B'C' bằng:\par
\hspace*{1.5em}\textbf{A.} 72 cm$^2$\quad\quad \textbf{B.} 540 cm$^2$\quad\quad \textbf{C.} 8100 cm$^2$\quad\quad \textbf{D.} 2.4 cm$^2$\par\vspace{2pt}
\textbf{\color{primary} Câu 35.} Dựa trên các tính chất trọng tâm của Chương IX, xét bài toán: Bộ ba số đo nào sau đây là độ dài ba cạnh của một tam giác vuông:\par
\hspace*{1.5em}\textbf{A.} 5 cm, 10 cm, 12 cm\quad\quad \textbf{B.} 9 cm, 12 cm, 15 cm\quad\quad \textbf{C.} 6 cm, 7 cm, 8 cm\quad\quad \textbf{D.} 4 cm, 5 cm, 6 cm\par\vspace{2pt}
\textbf{\color{primary} Câu 36.} Thực hành giải toán (Tam giác đồng dạng): Cho $\Delta$ABC $\sim$ $\Delta$DEF theo tỉ số k = 2/3. Biết chu vi $\Delta$ABC bằng 72 cm, chu vi của $\Delta$DEF bằng. Đáp án nào dưới đây thỏa mãn?\par
\hspace*{1.5em}\textbf{A.} 144 cm\quad\quad \textbf{B.} 162 cm\quad\quad \textbf{C.} 38 cm\quad\quad \textbf{D.} 108 cm\par\vspace{2pt}
\textbf{\color{primary} Câu 37.} Kiểm tra kỹ năng biến đổi và suy luận qua bài toán: Cho $\Delta$ABC $\sim$ $\Delta$A'B'C' theo tỉ số đồng dạng k = 1/18. Biết diện tích $\Delta$ABC bằng 42 cm$^2$, diện tích $\Delta$A'B'C' bằng:\par
\hspace*{1.5em}\textbf{A.} 13608 cm$^2$\quad\quad \textbf{B.} 84 cm$^2$\quad\quad \textbf{C.} 756 cm$^2$\quad\quad \textbf{D.} 2.3 cm$^2$\par\vspace{2pt}
\textbf{\color{primary} Câu 38.} Cho giả thiết thuộc phạm vi Chương IX: Bộ ba số đo nào sau đây là độ dài ba cạnh của một tam giác vuông. Kết luận đúng là:\par
\hspace*{1.5em}\textbf{A.} 5 cm, 10 cm, 12 cm\quad\quad \textbf{B.} 4 cm, 5 cm, 6 cm\quad\quad \textbf{C.} 8 cm, 15 cm, 17 cm\quad\quad \textbf{D.} 6 cm, 7 cm, 8 cm\par\vspace{2pt}
\textbf{\color{primary} Câu 39.} Rút gọn hoặc tính giá trị theo yêu cầu sau của Chương IX: Cho $\Delta$ABC $\sim$ $\Delta$DEF theo tỉ số k = 2/3. Biết chu vi $\Delta$ABC bằng 84 cm, chu vi của $\Delta$DEF bằng:\par
\hspace*{1.5em}\textbf{A.} 126 cm\quad\quad \textbf{B.} 168 cm\quad\quad \textbf{C.} 44 cm\quad\quad \textbf{D.} 189 cm\par\vspace{2pt}
\textbf{\color{primary} Câu 40.} Xác định phương án đúng sau khi thực hiện phép tính / suy luận: Cho $\Delta$ABC $\sim$ $\Delta$A'B'C' theo tỉ số đồng dạng k = 1/21. Biết diện tích $\Delta$ABC bằng 48 cm$^2$, diện tích $\Delta$A'B'C' bằng:\par
\hspace*{1.5em}\textbf{A.} 1008 cm$^2$\quad\quad \textbf{B.} 96 cm$^2$\quad\quad \textbf{C.} 21168 cm$^2$\quad\quad \textbf{D.} 2.3 cm$^2$\par\vspace{2pt}
\subsection*{\color{accent} Phần III. Mức độ Khó (Vận dụng \& Vận dụng cao -- Câu 41 đến Câu 50)}
\textbf{\color{primary} Câu 41.} Một cột cờ thẳng đứng có bóng trên mặt đất dài 6 m. Cùng thời điểm đó, một thanh sắt cao 1,5 m cắm vuông góc với mặt đất có bóng dài 2 m. Chiều cao của cột cờ là:\par
\hspace*{1.5em}\textbf{A.} 4.5 m\quad\quad \textbf{B.} 8.0 m\quad\quad \textbf{C.} 7.5 m\quad\quad \textbf{D.} 9.0 m\par\vspace{2pt}
\textbf{\color{primary} Câu 42.} Bài toán vận dụng (Chương IX – Tam giác đồng dạng): Cho tam giác ABC có diện tích bằng 45 cm$^2$. Trên cạnh AB lấy điểm M sao cho AM = 2/3 AB, qua M kẻ đường thẳng song song với BC cắt AC tại N. Diện tích tam giác AMN bằng. Kết quả đúng là:\par
\hspace*{1.5em}\textbf{A.} 20 cm$^2$\quad\quad \textbf{B.} 30 cm$^2$\quad\quad \textbf{C.} 25 cm$^2$\quad\quad \textbf{D.} 15 cm$^2$\par\vspace{2pt}
\textbf{\color{primary} Câu 43.} Vận dụng kiến thức nâng cao trong phạm vi Chương IX để giải bài toán: Một cột cờ thẳng đứng có bóng trên mặt đất dài 10 m. Cùng thời điểm đó, một thanh sắt cao 1,5 m cắm vuông góc với mặt đất có bóng dài 2 m. Chiều cao của cột cờ là:\par
\hspace*{1.5em}\textbf{A.} 15.0 m\quad\quad \textbf{B.} 7.5 m\quad\quad \textbf{C.} 11.5 m\quad\quad \textbf{D.} 13.3 m\par\vspace{2pt}
\textbf{\color{primary} Câu 44.} Phân tích và giải bài toán tổng hợp sau của Chương IX: Cho tam giác ABC có diện tích bằng 63 cm$^2$. Trên cạnh AB lấy điểm M sao cho AM = 2/3 AB, qua M kẻ đường thẳng song song với BC cắt AC tại N. Diện tích tam giác AMN bằng:\par
\hspace*{1.5em}\textbf{A.} 21 cm$^2$\quad\quad \textbf{B.} 28 cm$^2$\quad\quad \textbf{C.} 42 cm$^2$\quad\quad \textbf{D.} 35 cm$^2$\par\vspace{2pt}
\textbf{\color{primary} Câu 45.} Thử thách tư duy trong phạm vi kiến thức Chương IX: Một cột cờ thẳng đứng có bóng trên mặt đất dài 14 m. Cùng thời điểm đó, một thanh sắt cao 1,5 m cắm vuông góc với mặt đất có bóng dài 2 m. Chiều cao của cột cờ là. Đáp số chính xác là:\par
\hspace*{1.5em}\textbf{A.} 10.5 m\quad\quad \textbf{B.} 18.7 m\quad\quad \textbf{C.} 15.5 m\quad\quad \textbf{D.} 21.0 m\par\vspace{2pt}
\textbf{\color{primary} Câu 46.} Sử dụng phương pháp giải toán của Chương IX để tìm đáp án cho bài toán: Cho tam giác ABC có diện tích bằng 45 cm$^2$. Trên cạnh AB lấy điểm M sao cho AM = 2/3 AB, qua M kẻ đường thẳng song song với BC cắt AC tại N. Diện tích tam giác AMN bằng:\par
\hspace*{1.5em}\textbf{A.} 15 cm$^2$\quad\quad \textbf{B.} 20 cm$^2$\quad\quad \textbf{C.} 25 cm$^2$\quad\quad \textbf{D.} 30 cm$^2$\par\vspace{2pt}
\textbf{\color{primary} Câu 47.} Xét bài toán vận dụng thực tiễn / nâng cao thuộc Chương IX: Một cột cờ thẳng đứng có bóng trên mặt đất dài 18 m. Cùng thời điểm đó, một thanh sắt cao 1,5 m cắm vuông góc với mặt đất có bóng dài 2 m. Chiều cao của cột cờ là:\par
\hspace*{1.5em}\textbf{A.} 19.5 m\quad\quad \textbf{B.} 13.5 m\quad\quad \textbf{C.} 24.0 m\quad\quad \textbf{D.} 27.0 m\par\vspace{2pt}
\textbf{\color{primary} Câu 48.} Tìm kết quả cuối cùng của bài toán phân loại học sinh khá giỏi (Chương IX): Cho tam giác ABC có diện tích bằng 63 cm$^2$. Trên cạnh AB lấy điểm M sao cho AM = 2/3 AB, qua M kẻ đường thẳng song song với BC cắt AC tại N. Diện tích tam giác AMN bằng:\par
\hspace*{1.5em}\textbf{A.} 35 cm$^2$\quad\quad \textbf{B.} 42 cm$^2$\quad\quad \textbf{C.} 21 cm$^2$\quad\quad \textbf{D.} 28 cm$^2$\par\vspace{2pt}
\textbf{\color{primary} Câu 49.} Kết hợp các tính chất đã học trong Chương IX, hãy giải quyết bài toán: Một cột cờ thẳng đứng có bóng trên mặt đất dài 22 m. Cùng thời điểm đó, một thanh sắt cao 1,5 m cắm vuông góc với mặt đất có bóng dài 2 m. Chiều cao của cột cờ là:\par
\hspace*{1.5em}\textbf{A.} 23.5 m\quad\quad \textbf{B.} 29.3 m\quad\quad \textbf{C.} 33.0 m\quad\quad \textbf{D.} 16.5 m\par\vspace{2pt}
\textbf{\color{primary} Câu 50.} Câu hỏi vận dụng cao cuối Chương IX: Cho tam giác ABC có diện tích bằng 45 cm$^2$. Trên cạnh AB lấy điểm M sao cho AM = 2/3 AB, qua M kẻ đường thẳng song song với BC cắt AC tại N. Diện tích tam giác AMN bằng. Phương án đúng là:\par
\hspace*{1.5em}\textbf{A.} 20 cm$^2$\quad\quad \textbf{B.} 15 cm$^2$\quad\quad \textbf{C.} 30 cm$^2$\quad\quad \textbf{D.} 25 cm$^2$\par\vspace{2pt}
\subsection*{\color{primary} Bảng đáp án Chương IX: Tam giác đồng dạng}
\begin{center}
\renewcommand{\arraystretch}{1.25}
\begin{tabularx}{\textwidth}{*{10}{>{\centering\arraybackslash}X}}
\toprule
\textbf{\color{primary} Cột 1} & \textbf{\color{primary} Cột 2} & \textbf{\color{primary} Cột 3} & \textbf{\color{primary} Cột 4} & \textbf{\color{primary} Cột 5} & \textbf{\color{primary} Cột 6} & \textbf{\color{primary} Cột 7} & \textbf{\color{primary} Cột 8} & \textbf{\color{primary} Cột 9} & \textbf{\color{primary} Cột 10} \\
\midrule
\textbf{1.B} & \textbf{2.A} & \textbf{3.D} & \textbf{4.B} & \textbf{5.C} & \textbf{6.A} & \textbf{7.D} & \textbf{8.A} & \textbf{9.A} & \textbf{10.D} \\
\rowcolor{rowtint} \textbf{11.C} & \textbf{12.B} & \textbf{13.A} & \textbf{14.C} & \textbf{15.B} & \textbf{16.A} & \textbf{17.A} & \textbf{18.C} & \textbf{19.D} & \textbf{20.A} \\
\textbf{21.B} & \textbf{22.B} & \textbf{23.C} & \textbf{24.A} & \textbf{25.C} & \textbf{26.B} & \textbf{27.D} & \textbf{28.A} & \textbf{29.B} & \textbf{30.A} \\
\rowcolor{rowtint} \textbf{31.D} & \textbf{32.C} & \textbf{33.B} & \textbf{34.C} & \textbf{35.B} & \textbf{36.D} & \textbf{37.A} & \textbf{38.C} & \textbf{39.A} & \textbf{40.C} \\
\textbf{41.A} & \textbf{42.A} & \textbf{43.B} & \textbf{44.B} & \textbf{45.A} & \textbf{46.B} & \textbf{47.B} & \textbf{48.D} & \textbf{49.D} & \textbf{50.A} \\
\bottomrule
\end{tabularx}
\end{center}\vspace{10pt}
\section*{\color{primary} CHƯƠNG X: MỘT SỐ HÌNH KHỐI TRONG THỰC TIỄN}
\textit{Trọng tâm kiến thức SGK: Hình chóp tam giác đều và hình chóp tứ giác đều (đỉnh, cạnh bên, mặt bên, trung đoạn, diện tích xung quanh và thể tích).}\par\vspace{4pt}
\subsection*{\color{accent} Phần I. Mức độ Dễ (Nhận biết -- Câu 1 đến Câu 20)}
\textbf{\color{primary} Câu 1.} Một hình chóp tứ giác đều có diện tích đáy S = 12 cm$^2$ và chiều cao h = 4 cm. Thể tích của hình chóp đó bằng:\par
\hspace*{1.5em}\textbf{A.} 144 cm$^3$\quad\quad \textbf{B.} 48 cm$^3$\quad\quad \textbf{C.} 16 cm$^3$\quad\quad \textbf{D.} 24 cm$^3$\par\vspace{2pt}
\textbf{\color{primary} Câu 2.} Hãy trả lời câu hỏi cơ bản sau: Các mặt bên của hình chóp tam giác đều và hình chóp tứ giác đều là những hình gì?\par
\begin{itemize}\setlength{\itemsep}{1pt}\setlength{\parskip}{0pt}
  \item[] \textbf{A.} Các tam giác vuông
  \item[] \textbf{B.} Các hình thang cân bằng nhau
  \item[] \textbf{C.} Các hình chữ nhật bằng nhau
  \item[] \textbf{D.} Các tam giác cân bằng nhau
\end{itemize}\vspace{2pt}
\textbf{\color{primary} Câu 3.} Trong bài học thuộc Chương X, xét nội dung: Một hình chóp tam giác đều có chu vi đáy bằng 30 cm và độ dài trung đoạn d = 7 cm. Diện tích xung quanh của hình chóp bằng. Đáp án đúng là:\par
\hspace*{1.5em}\textbf{A.} 210 cm$^2$\quad\quad \textbf{B.} 70 cm$^2$\quad\quad \textbf{C.} 105 cm$^2$\quad\quad \textbf{D.} 420 cm$^2$\par\vspace{2pt}
\textbf{\color{primary} Câu 4.} Dựa vào kiến thức chuẩn SGK của Chương X (Một số hình khối trong thực tiễn), một hình chóp tứ giác đều có diện tích đáy S = 21 cm$^2$ và chiều cao h = 7 cm. Thể tích của hình chóp đó bằng:\par
\hspace*{1.5em}\textbf{A.} 49 cm$^3$\quad\quad \textbf{B.} 73 cm$^3$\quad\quad \textbf{C.} 441 cm$^3$\quad\quad \textbf{D.} 147 cm$^3$\par\vspace{2pt}
\textbf{\color{primary} Câu 5.} Chọn phương án trả lời đúng nhất: Các mặt bên của hình chóp tam giác đều và hình chóp tứ giác đều là những hình gì?\par
\begin{itemize}\setlength{\itemsep}{1pt}\setlength{\parskip}{0pt}
  \item[] \textbf{A.} Các hình thang cân bằng nhau
  \item[] \textbf{B.} Các hình chữ nhật bằng nhau
  \item[] \textbf{C.} Các tam giác vuông
  \item[] \textbf{D.} Các tam giác cân bằng nhau
\end{itemize}\vspace{2pt}
\textbf{\color{primary} Câu 6.} Khi làm bài tập cơ bản về một số hình khối trong thực tiễn, một học sinh gặp câu hỏi: Một hình chóp tam giác đều có chu vi đáy bằng 48 cm và độ dài trung đoạn d = 10 cm. Diện tích xung quanh của hình chóp bằng. Kết quả nào dưới đây đúng?\par
\hspace*{1.5em}\textbf{A.} 240 cm$^2$\quad\quad \textbf{B.} 160 cm$^2$\quad\quad \textbf{C.} 960 cm$^2$\quad\quad \textbf{D.} 480 cm$^2$\par\vspace{2pt}
\textbf{\color{primary} Câu 7.} Xác định giá trị hoặc khẳng định đúng cho bài toán sau: Một hình chóp tứ giác đều có diện tích đáy S = 30 cm$^2$ và chiều cao h = 10 cm. Thể tích của hình chóp đó bằng:\par
\hspace*{1.5em}\textbf{A.} 300 cm$^3$\quad\quad \textbf{B.} 100 cm$^3$\quad\quad \textbf{C.} 150 cm$^3$\quad\quad \textbf{D.} 900 cm$^3$\par\vspace{2pt}
\textbf{\color{primary} Câu 8.} Áp dụng trực tiếp định nghĩa và quy tắc trong Chương X: Các mặt bên của hình chóp tam giác đều và hình chóp tứ giác đều là những hình gì:\par
\begin{itemize}\setlength{\itemsep}{1pt}\setlength{\parskip}{0pt}
  \item[] \textbf{A.} Các tam giác vuông
  \item[] \textbf{B.} Các hình thang cân bằng nhau
  \item[] \textbf{C.} Các tam giác cân bằng nhau
  \item[] \textbf{D.} Các hình chữ nhật bằng nhau
\end{itemize}\vspace{2pt}
\textbf{\color{primary} Câu 9.} Kiểm tra nhanh kiến thức nền tảng: Một hình chóp tam giác đều có chu vi đáy bằng 66 cm và độ dài trung đoạn d = 13 cm. Diện tích xung quanh của hình chóp bằng:\par
\hspace*{1.5em}\textbf{A.} 858 cm$^2$\quad\quad \textbf{B.} 429 cm$^2$\quad\quad \textbf{C.} 286 cm$^2$\quad\quad \textbf{D.} 1716 cm$^2$\par\vspace{2pt}
\textbf{\color{primary} Câu 10.} Cho dữ kiện nhận biết sau đây: Một hình chóp tứ giác đều có diện tích đáy S = 39 cm$^2$ và chiều cao h = 13 cm. Thể tích của hình chóp đó bằng. Phương án chính xác là:\par
\hspace*{1.5em}\textbf{A.} 1521 cm$^3$\quad\quad \textbf{B.} 507 cm$^3$\quad\quad \textbf{C.} 253 cm$^3$\quad\quad \textbf{D.} 169 cm$^3$\par\vspace{2pt}
\textbf{\color{primary} Câu 11.} Hãy chỉ ra kết quả đúng trong bốn lựa chọn A, B, C, D đối với yêu cầu: Các mặt bên của hình chóp tam giác đều và hình chóp tứ giác đều là những hình gì:\par
\begin{itemize}\setlength{\itemsep}{1pt}\setlength{\parskip}{0pt}
  \item[] \textbf{A.} Các hình chữ nhật bằng nhau
  \item[] \textbf{B.} Các tam giác vuông
  \item[] \textbf{C.} Các hình thang cân bằng nhau
  \item[] \textbf{D.} Các tam giác cân bằng nhau
\end{itemize}\vspace{2pt}
\textbf{\color{primary} Câu 12.} Câu hỏi ôn tập lý thuyết và kỹ năng cơ bản (Chương X): Một hình chóp tam giác đều có chu vi đáy bằng 84 cm và độ dài trung đoạn d = 16 cm. Diện tích xung quanh của hình chóp bằng:\par
\hspace*{1.5em}\textbf{A.} 2688 cm$^2$\quad\quad \textbf{B.} 672 cm$^2$\quad\quad \textbf{C.} 448 cm$^2$\quad\quad \textbf{D.} 1344 cm$^2$\par\vspace{2pt}
\textbf{\color{primary} Câu 13.} Theo nội dung sách giáo khoa Kết nối tri thức với cuộc sống (Chương X), một hình chóp tứ giác đều có diện tích đáy S = 48 cm$^2$ và chiều cao h = 16 cm. Thể tích của hình chóp đó bằng:\par
\hspace*{1.5em}\textbf{A.} 256 cm$^3$\quad\quad \textbf{B.} 2304 cm$^3$\quad\quad \textbf{C.} 384 cm$^3$\quad\quad \textbf{D.} 768 cm$^3$\par\vspace{2pt}
\textbf{\color{primary} Câu 14.} Tìm đáp số đúng của bài toán cơ bản sau: Các mặt bên của hình chóp tam giác đều và hình chóp tứ giác đều là những hình gì:\par
\begin{itemize}\setlength{\itemsep}{1pt}\setlength{\parskip}{0pt}
  \item[] \textbf{A.} Các hình chữ nhật bằng nhau
  \item[] \textbf{B.} Các tam giác vuông
  \item[] \textbf{C.} Các tam giác cân bằng nhau
  \item[] \textbf{D.} Các hình thang cân bằng nhau
\end{itemize}\vspace{2pt}
\textbf{\color{primary} Câu 15.} Xét đối tượng toán học trong Chương X: Một hình chóp tam giác đều có chu vi đáy bằng 102 cm và độ dài trung đoạn d = 19 cm. Diện tích xung quanh của hình chóp bằng. Giá trị hoặc tính chất tương ứng là:\par
\hspace*{1.5em}\textbf{A.} 969 cm$^2$\quad\quad \textbf{B.} 1938 cm$^2$\quad\quad \textbf{C.} 646 cm$^2$\quad\quad \textbf{D.} 3876 cm$^2$\par\vspace{2pt}
\textbf{\color{primary} Câu 16.} Hoàn thành phép xác định cơ bản sau đây: Một hình chóp tứ giác đều có diện tích đáy S = 57 cm$^2$ và chiều cao h = 19 cm. Thể tích của hình chóp đó bằng:\par
\hspace*{1.5em}\textbf{A.} 361 cm$^3$\quad\quad \textbf{B.} 541 cm$^3$\quad\quad \textbf{C.} 3249 cm$^3$\quad\quad \textbf{D.} 1083 cm$^3$\par\vspace{2pt}
\textbf{\color{primary} Câu 17.} Trong các khẳng định hoặc giá trị dưới đây, đâu là đáp án đúng cho câu hỏi: Các mặt bên của hình chóp tam giác đều và hình chóp tứ giác đều là những hình gì?\par
\begin{itemize}\setlength{\itemsep}{1pt}\setlength{\parskip}{0pt}
  \item[] \textbf{A.} Các tam giác vuông
  \item[] \textbf{B.} Các hình chữ nhật bằng nhau
  \item[] \textbf{C.} Các hình thang cân bằng nhau
  \item[] \textbf{D.} Các tam giác cân bằng nhau
\end{itemize}\vspace{2pt}
\textbf{\color{primary} Câu 18.} Bài tập khởi động Chương X: Một hình chóp tam giác đều có chu vi đáy bằng 120 cm và độ dài trung đoạn d = 22 cm. Diện tích xung quanh của hình chóp bằng:\par
\hspace*{1.5em}\textbf{A.} 5280 cm$^2$\quad\quad \textbf{B.} 2640 cm$^2$\quad\quad \textbf{C.} 880 cm$^2$\quad\quad \textbf{D.} 1320 cm$^2$\par\vspace{2pt}
\textbf{\color{primary} Câu 19.} Quan sát dữ kiện đề bài: Một hình chóp tứ giác đều có diện tích đáy S = 66 cm$^2$ và chiều cao h = 22 cm. Thể tích của hình chóp đó bằng. Hãy chọn phương án đúng:\par
\hspace*{1.5em}\textbf{A.} 726 cm$^3$\quad\quad \textbf{B.} 4356 cm$^3$\quad\quad \textbf{C.} 1452 cm$^3$\quad\quad \textbf{D.} 484 cm$^3$\par\vspace{2pt}
\textbf{\color{primary} Câu 20.} Sử dụng công thức hoặc tính chất đầu bài của Chương X để giải quyết yêu cầu: Các mặt bên của hình chóp tam giác đều và hình chóp tứ giác đều là những hình gì:\par
\begin{itemize}\setlength{\itemsep}{1pt}\setlength{\parskip}{0pt}
  \item[] \textbf{A.} Các hình chữ nhật bằng nhau
  \item[] \textbf{B.} Các tam giác vuông
  \item[] \textbf{C.} Các tam giác cân bằng nhau
  \item[] \textbf{D.} Các hình thang cân bằng nhau
\end{itemize}\vspace{2pt}
\subsection*{\color{accent} Phần II. Mức độ Trung bình (Thông hiểu -- Câu 21 đến Câu 40)}
\textbf{\color{primary} Câu 21.} Một hình chóp tứ giác đều có cạnh đáy bằng 4 cm và trung đoạn bằng 6 cm. Diện tích xung quanh của hình chóp tứ giác đều đó là:\par
\hspace*{1.5em}\textbf{A.} 24 cm$^2$\quad\quad \textbf{B.} 48 cm$^2$\quad\quad \textbf{C.} 64 cm$^2$\quad\quad \textbf{D.} 96 cm$^2$\par\vspace{2pt}
\textbf{\color{primary} Câu 22.} Vận dụng quy tắc trong Chương X (Một số hình khối trong thực tiễn) để giải bài toán: Một hình chóp tứ giác đều có cạnh đáy bằng 5 cm và trung đoạn bằng 6 cm. Diện tích toàn phần của hình chóp đó bằng:\par
\hspace*{1.5em}\textbf{A.} 150 cm$^2$\quad\quad \textbf{B.} 145 cm$^2$\quad\quad \textbf{C.} 85 cm$^2$\quad\quad \textbf{D.} 60 cm$^2$\par\vspace{2pt}
\textbf{\color{primary} Câu 23.} Biến đổi và tính toán theo kiến thức Chương X: Một hình chóp tứ giác đều có thể tích V = 64 cm$^3$ và cạnh đáy bằng 4 cm. Chiều cao của hình chóp đó bằng. Kết quả thu được là:\par
\hspace*{1.5em}\textbf{A.} 12 cm\quad\quad \textbf{B.} 4 cm\quad\quad \textbf{C.} 16 cm\quad\quad \textbf{D.} 48 cm\par\vspace{2pt}
\textbf{\color{primary} Câu 24.} Xét bài toán mức độ thông hiểu sau: Một hình chóp tứ giác đều có cạnh đáy bằng 7 cm và trung đoạn bằng 9 cm. Diện tích xung quanh của hình chóp tứ giác đều đó là. Đáp án đúng là:\par
\hspace*{1.5em}\textbf{A.} 126 cm$^2$\quad\quad \textbf{B.} 252 cm$^2$\quad\quad \textbf{C.} 175 cm$^2$\quad\quad \textbf{D.} 63 cm$^2$\par\vspace{2pt}
\textbf{\color{primary} Câu 25.} Để hoàn thành phiếu học tập Chương X, hãy xác định: Một hình chóp tứ giác đều có cạnh đáy bằng 8 cm và trung đoạn bằng 9 cm. Diện tích toàn phần của hình chóp đó bằng:\par
\hspace*{1.5em}\textbf{A.} 208 cm$^2$\quad\quad \textbf{B.} 576 cm$^2$\quad\quad \textbf{C.} 352 cm$^2$\quad\quad \textbf{D.} 144 cm$^2$\par\vspace{2pt}
\textbf{\color{primary} Câu 26.} Phân tích dữ kiện bài toán: Một hình chóp tứ giác đều có thể tích V = 112 cm$^3$ và cạnh đáy bằng 4 cm. Chiều cao của hình chóp đó bằng. Giá trị cần tìm bằng:\par
\hspace*{1.5em}\textbf{A.} 21 cm\quad\quad \textbf{B.} 84 cm\quad\quad \textbf{C.} 28 cm\quad\quad \textbf{D.} 7 cm\par\vspace{2pt}
\textbf{\color{primary} Câu 27.} Áp dụng tính chất đã học ở Chương X vào bài tập sau: Một hình chóp tứ giác đều có cạnh đáy bằng 10 cm và trung đoạn bằng 12 cm. Diện tích xung quanh của hình chóp tứ giác đều đó là:\par
\hspace*{1.5em}\textbf{A.} 120 cm$^2$\quad\quad \textbf{B.} 240 cm$^2$\quad\quad \textbf{C.} 480 cm$^2$\quad\quad \textbf{D.} 340 cm$^2$\par\vspace{2pt}
\textbf{\color{primary} Câu 28.} Thực hiện các bước suy luận hợp lý để giải quyết yêu cầu: Một hình chóp tứ giác đều có cạnh đáy bằng 11 cm và trung đoạn bằng 12 cm. Diện tích toàn phần của hình chóp đó bằng:\par
\hspace*{1.5em}\textbf{A.} 1452 cm$^2$\quad\quad \textbf{B.} 385 cm$^2$\quad\quad \textbf{C.} 264 cm$^2$\quad\quad \textbf{D.} 649 cm$^2$\par\vspace{2pt}
\textbf{\color{primary} Câu 29.} Cho bài toán tự luyện (Một số hình khối trong thực tiễn): Một hình chóp tứ giác đều có thể tích V = 160 cm$^3$ và cạnh đáy bằng 4 cm. Chiều cao của hình chóp đó bằng. Phương án nào sau đây chính xác?\par
\hspace*{1.5em}\textbf{A.} 10 cm\quad\quad \textbf{B.} 40 cm\quad\quad \textbf{C.} 30 cm\quad\quad \textbf{D.} 120 cm\par\vspace{2pt}
\textbf{\color{primary} Câu 30.} Kết hợp các phép biến đổi trong Chương X, hãy tìm kết quả của bài toán: Một hình chóp tứ giác đều có cạnh đáy bằng 13 cm và trung đoạn bằng 15 cm. Diện tích xung quanh của hình chóp tứ giác đều đó là:\par
\hspace*{1.5em}\textbf{A.} 780 cm$^2$\quad\quad \textbf{B.} 195 cm$^2$\quad\quad \textbf{C.} 390 cm$^2$\quad\quad \textbf{D.} 559 cm$^2$\par\vspace{2pt}
\textbf{\color{primary} Câu 31.} Một bài kiểm tra định kỳ nội dung Chương X nêu câu hỏi: Một hình chóp tứ giác đều có cạnh đáy bằng 14 cm và trung đoạn bằng 15 cm. Diện tích toàn phần của hình chóp đó bằng. Hãy chọn đáp số đúng:\par
\hspace*{1.5em}\textbf{A.} 1036 cm$^2$\quad\quad \textbf{B.} 2940 cm$^2$\quad\quad \textbf{C.} 616 cm$^2$\quad\quad \textbf{D.} 420 cm$^2$\par\vspace{2pt}
\textbf{\color{primary} Câu 32.} Tìm lời giải đúng cho bài toán thông hiểu: Một hình chóp tứ giác đều có thể tích V = 208 cm$^3$ và cạnh đáy bằng 4 cm. Chiều cao của hình chóp đó bằng:\par
\hspace*{1.5em}\textbf{A.} 52 cm\quad\quad \textbf{B.} 13 cm\quad\quad \textbf{C.} 156 cm\quad\quad \textbf{D.} 39 cm\par\vspace{2pt}
\textbf{\color{primary} Câu 33.} Sử dụng mối liên hệ giữa các đại lượng trong Chương X: Một hình chóp tứ giác đều có cạnh đáy bằng 16 cm và trung đoạn bằng 18 cm. Diện tích xung quanh của hình chóp tứ giác đều đó là:\par
\hspace*{1.5em}\textbf{A.} 288 cm$^2$\quad\quad \textbf{B.} 832 cm$^2$\quad\quad \textbf{C.} 576 cm$^2$\quad\quad \textbf{D.} 1152 cm$^2$\par\vspace{2pt}
\textbf{\color{primary} Câu 34.} Hãy tính toán và chọn phương án đúng đối với bài tập: Một hình chóp tứ giác đều có cạnh đáy bằng 17 cm và trung đoạn bằng 18 cm. Diện tích toàn phần của hình chóp đó bằng:\par
\hspace*{1.5em}\textbf{A.} 1513 cm$^2$\quad\quad \textbf{B.} 5202 cm$^2$\quad\quad \textbf{C.} 612 cm$^2$\quad\quad \textbf{D.} 901 cm$^2$\par\vspace{2pt}
\textbf{\color{primary} Câu 35.} Dựa trên các tính chất trọng tâm của Chương X, xét bài toán: Một hình chóp tứ giác đều có thể tích V = 256 cm$^3$ và cạnh đáy bằng 4 cm. Chiều cao của hình chóp đó bằng:\par
\hspace*{1.5em}\textbf{A.} 48 cm\quad\quad \textbf{B.} 16 cm\quad\quad \textbf{C.} 192 cm\quad\quad \textbf{D.} 64 cm\par\vspace{2pt}
\textbf{\color{primary} Câu 36.} Thực hành giải toán (Một số hình khối trong thực tiễn): Một hình chóp tứ giác đều có cạnh đáy bằng 19 cm và trung đoạn bằng 21 cm. Diện tích xung quanh của hình chóp tứ giác đều đó là. Đáp án nào dưới đây thỏa mãn?\par
\hspace*{1.5em}\textbf{A.} 399 cm$^2$\quad\quad \textbf{B.} 798 cm$^2$\quad\quad \textbf{C.} 1159 cm$^2$\quad\quad \textbf{D.} 1596 cm$^2$\par\vspace{2pt}
\textbf{\color{primary} Câu 37.} Kiểm tra kỹ năng biến đổi và suy luận qua bài toán: Một hình chóp tứ giác đều có cạnh đáy bằng 20 cm và trung đoạn bằng 21 cm. Diện tích toàn phần của hình chóp đó bằng:\par
\hspace*{1.5em}\textbf{A.} 840 cm$^2$\quad\quad \textbf{B.} 1240 cm$^2$\quad\quad \textbf{C.} 2080 cm$^2$\quad\quad \textbf{D.} 8400 cm$^2$\par\vspace{2pt}
\textbf{\color{primary} Câu 38.} Cho giả thiết thuộc phạm vi Chương X: Một hình chóp tứ giác đều có thể tích V = 304 cm$^3$ và cạnh đáy bằng 4 cm. Chiều cao của hình chóp đó bằng. Kết luận đúng là:\par
\hspace*{1.5em}\textbf{A.} 228 cm\quad\quad \textbf{B.} 57 cm\quad\quad \textbf{C.} 76 cm\quad\quad \textbf{D.} 19 cm\par\vspace{2pt}
\textbf{\color{primary} Câu 39.} Rút gọn hoặc tính giá trị theo yêu cầu sau của Chương X: Một hình chóp tứ giác đều có cạnh đáy bằng 22 cm và trung đoạn bằng 24 cm. Diện tích xung quanh của hình chóp tứ giác đều đó là:\par
\hspace*{1.5em}\textbf{A.} 528 cm$^2$\quad\quad \textbf{B.} 1540 cm$^2$\quad\quad \textbf{C.} 1056 cm$^2$\quad\quad \textbf{D.} 2112 cm$^2$\par\vspace{2pt}
\textbf{\color{primary} Câu 40.} Xác định phương án đúng sau khi thực hiện phép tính / suy luận: Một hình chóp tứ giác đều có cạnh đáy bằng 23 cm và trung đoạn bằng 24 cm. Diện tích toàn phần của hình chóp đó bằng:\par
\hspace*{1.5em}\textbf{A.} 12696 cm$^2$\quad\quad \textbf{B.} 1104 cm$^2$\quad\quad \textbf{C.} 1633 cm$^2$\quad\quad \textbf{D.} 2737 cm$^2$\par\vspace{2pt}
\subsection*{\color{accent} Phần III. Mức độ Khó (Vận dụng \& Vận dụng cao -- Câu 41 đến Câu 50)}
\textbf{\color{primary} Câu 41.} Cho hình chóp tứ giác đều S.ABCD có cạnh đáy AB = 6 cm và chiều cao SO = 4 cm. Độ dài trung đoạn SI của hình chóp bằng:\par
\hspace*{1.5em}\textbf{A.} 7 cm\quad\quad \textbf{B.} 6 cm\quad\quad \textbf{C.} 10 cm\quad\quad \textbf{D.} 5 cm\par\vspace{2pt}
\textbf{\color{primary} Câu 42.} Bài toán vận dụng (Chương X – Một số hình khối trong thực tiễn): Một mái nhà chóp tứ giác đều có cạnh đáy dài 8 m và trung đoạn dài 5 m. Chi phí lợp ngói bốn mặt xung quanh mái là 250.000 đồng/m$^2$. Tổng tiền ngói cần mua là. Kết quả đúng là:\par
\hspace*{1.5em}\textbf{A.} 40000.000 đồng\quad\quad \textbf{B.} 20000.000 đồng\quad\quad \textbf{C.} 16000.000 đồng\quad\quad \textbf{D.} 10000.000 đồng\par\vspace{2pt}
\textbf{\color{primary} Câu 43.} Vận dụng kiến thức nâng cao trong phạm vi Chương X để giải bài toán: Cho hình chóp tứ giác đều S.ABCD có cạnh đáy bằng 18 cm và chiều cao SO = 12 cm. Độ dài trung đoạn SI và diện tích xung quanh của hình chóp lần lượt là:\par
\begin{itemize}\setlength{\itemsep}{1pt}\setlength{\parskip}{0pt}
  \item[] \textbf{A.} SI = 15 cm; $S_{xq}$ = 1080 cm$^2$
  \item[] \textbf{B.} SI = 21 cm; $S_{xq}$ = 756 cm$^2$
  \item[] \textbf{C.} SI = 15 cm; $S_{xq}$ = 270 cm$^2$
  \item[] \textbf{D.} SI = 15 cm; $S_{xq}$ = 540 cm$^2$
\end{itemize}\vspace{2pt}
\textbf{\color{primary} Câu 44.} Phân tích và giải bài toán tổng hợp sau của Chương X: Cho hình chóp tứ giác đều S.ABCD có cạnh đáy AB = 12 cm và chiều cao SO = 8 cm. Độ dài trung đoạn SI của hình chóp bằng:\par
\hspace*{1.5em}\textbf{A.} 14 cm\quad\quad \textbf{B.} 20 cm\quad\quad \textbf{C.} 12 cm\quad\quad \textbf{D.} 10 cm\par\vspace{2pt}
\textbf{\color{primary} Câu 45.} Thử thách tư duy trong phạm vi kiến thức Chương X: Một mái nhà chóp tứ giác đều có cạnh đáy dài 8 m và trung đoạn dài 5 m. Chi phí lợp ngói bốn mặt xung quanh mái là 250.000 đồng/m$^2$. Tổng tiền ngói cần mua là. Đáp số chính xác là:\par
\hspace*{1.5em}\textbf{A.} 16000.000 đồng\quad\quad \textbf{B.} 10000.000 đồng\quad\quad \textbf{C.} 40000.000 đồng\quad\quad \textbf{D.} 20000.000 đồng\par\vspace{2pt}
\textbf{\color{primary} Câu 46.} Sử dụng phương pháp giải toán của Chương X để tìm đáp án cho bài toán: Cho hình chóp tứ giác đều S.ABCD có cạnh đáy bằng 36 cm và chiều cao SO = 24 cm. Độ dài trung đoạn SI và diện tích xung quanh của hình chóp lần lượt là:\par
\begin{itemize}\setlength{\itemsep}{1pt}\setlength{\parskip}{0pt}
  \item[] \textbf{A.} SI = 30 cm; $S_{xq}$ = 1080 cm$^2$
  \item[] \textbf{B.} SI = 30 cm; $S_{xq}$ = 4320 cm$^2$
  \item[] \textbf{C.} SI = 30 cm; $S_{xq}$ = 2160 cm$^2$
  \item[] \textbf{D.} SI = 42 cm; $S_{xq}$ = 3024 cm$^2$
\end{itemize}\vspace{2pt}
\textbf{\color{primary} Câu 47.} Xét bài toán vận dụng thực tiễn / nâng cao thuộc Chương X: Cho hình chóp tứ giác đều S.ABCD có cạnh đáy AB = 6 cm và chiều cao SO = 4 cm. Độ dài trung đoạn SI của hình chóp bằng:\par
\hspace*{1.5em}\textbf{A.} 5 cm\quad\quad \textbf{B.} 6 cm\quad\quad \textbf{C.} 10 cm\quad\quad \textbf{D.} 7 cm\par\vspace{2pt}
\textbf{\color{primary} Câu 48.} Tìm kết quả cuối cùng của bài toán phân loại học sinh khá giỏi (Chương X): Một mái nhà chóp tứ giác đều có cạnh đáy dài 8 m và trung đoạn dài 5 m. Chi phí lợp ngói bốn mặt xung quanh mái là 250.000 đồng/m$^2$. Tổng tiền ngói cần mua là:\par
\hspace*{1.5em}\textbf{A.} 20000.000 đồng\quad\quad \textbf{B.} 10000.000 đồng\quad\quad \textbf{C.} 16000.000 đồng\quad\quad \textbf{D.} 40000.000 đồng\par\vspace{2pt}
\textbf{\color{primary} Câu 49.} Kết hợp các tính chất đã học trong Chương X, hãy giải quyết bài toán: Cho hình chóp tứ giác đều S.ABCD có cạnh đáy bằng 54 cm và chiều cao SO = 36 cm. Độ dài trung đoạn SI và diện tích xung quanh của hình chóp lần lượt là:\par
\begin{itemize}\setlength{\itemsep}{1pt}\setlength{\parskip}{0pt}
  \item[] \textbf{A.} SI = 45 cm; $S_{xq}$ = 4860 cm$^2$
  \item[] \textbf{B.} SI = 45 cm; $S_{xq}$ = 9720 cm$^2$
  \item[] \textbf{C.} SI = 45 cm; $S_{xq}$ = 2430 cm$^2$
  \item[] \textbf{D.} SI = 63 cm; $S_{xq}$ = 6804 cm$^2$
\end{itemize}\vspace{2pt}
\textbf{\color{primary} Câu 50.} Câu hỏi vận dụng cao cuối Chương X: Cho hình chóp tứ giác đều S.ABCD có cạnh đáy AB = 12 cm và chiều cao SO = 8 cm. Độ dài trung đoạn SI của hình chóp bằng. Phương án đúng là:\par
\hspace*{1.5em}\textbf{A.} 20 cm\quad\quad \textbf{B.} 10 cm\quad\quad \textbf{C.} 12 cm\quad\quad \textbf{D.} 14 cm\par\vspace{2pt}
\subsection*{\color{primary} Bảng đáp án Chương X: Một số hình khối trong thực tiễn}
\begin{center}
\renewcommand{\arraystretch}{1.25}
\begin{tabularx}{\textwidth}{*{10}{>{\centering\arraybackslash}X}}
\toprule
\textbf{\color{primary} Cột 1} & \textbf{\color{primary} Cột 2} & \textbf{\color{primary} Cột 3} & \textbf{\color{primary} Cột 4} & \textbf{\color{primary} Cột 5} & \textbf{\color{primary} Cột 6} & \textbf{\color{primary} Cột 7} & \textbf{\color{primary} Cột 8} & \textbf{\color{primary} Cột 9} & \textbf{\color{primary} Cột 10} \\
\midrule
\textbf{1.C} & \textbf{2.D} & \textbf{3.C} & \textbf{4.A} & \textbf{5.D} & \textbf{6.A} & \textbf{7.B} & \textbf{8.C} & \textbf{9.B} & \textbf{10.D} \\
\rowcolor{rowtint} \textbf{11.D} & \textbf{12.B} & \textbf{13.A} & \textbf{14.C} & \textbf{15.A} & \textbf{16.A} & \textbf{17.D} & \textbf{18.D} & \textbf{19.D} & \textbf{20.C} \\
\textbf{21.B} & \textbf{22.C} & \textbf{23.A} & \textbf{24.A} & \textbf{25.A} & \textbf{26.A} & \textbf{27.B} & \textbf{28.B} & \textbf{29.C} & \textbf{30.C} \\
\rowcolor{rowtint} \textbf{31.C} & \textbf{32.D} & \textbf{33.C} & \textbf{34.D} & \textbf{35.A} & \textbf{36.B} & \textbf{37.B} & \textbf{38.B} & \textbf{39.C} & \textbf{40.C} \\
\textbf{41.D} & \textbf{42.B} & \textbf{43.D} & \textbf{44.D} & \textbf{45.D} & \textbf{46.C} & \textbf{47.A} & \textbf{48.A} & \textbf{49.A} & \textbf{50.B} \\
\bottomrule
\end{tabularx}
\end{center}\vspace{10pt}
\end{document}
